% 本文件由 LaTeX 排版 Agent 根据有效 translation.json 自动生成。
% author=Augustine of Hippo; work_id=1702; title=若望一书讲道集

\chapter{若望一书讲道集}

% structure: p0001 source=h1 target=omitted reason=page-title-already-rendered-by-parent

\Emph{第一篇讲道} \BibleRef{若一 1:1-2:11}\newline{}\newline{}\Emph{第二篇讲道} \BibleRef{若一 2:12-17}\newline{}\newline{}\Emph{第三篇讲道} \BibleRef{若一 2:18-27}\newline{}\newline{}\Emph{第四篇讲道} \BibleRef{若一 2:27-3:8}\newline{}\newline{}\Emph{第五篇讲道} \BibleRef{若一 3:9-18}\newline{}\newline{}\Emph{第六篇讲道} \BibleRef{若一 3:19-4:3}\newline{}\newline{}\Emph{第七篇讲道} \BibleRef{若一 4:4-12}\newline{}\newline{}\Emph{第八篇讲道} \BibleRef{若一 4:12-16}\newline{}\newline{}\Emph{第九篇讲道} \BibleRef{若一 4:17-21}\newline{}\newline{}\Emph{第十篇讲道} \BibleRef{若一 5:1-3}\par

\SourceRecord{Translated by H. Browne.；Nicene and Post-Nicene Fathers, First Series；Vol. 7.；Edited by Philip Schaff.；Buffalo, NY: Christian Literature Publishing Co.,；1888.；<http://www.newadvent.org/fathers/1702.htm>.}

% structure: p0001 source=h1 target=section reason=page-title

\section{论若望一书第一讲}

% structure: p0002 source=h2 target=subsection reason=relative-heading-rank; base=h2

\subsection{若一 1:1-2:11}

\BibleQuote{那从起初就有的，我们所听见过，我们亲眼看见过，我们亲手摸过的，论到生命的圣言——这生命已显现出来，我们看见了，并作证，且将原与父同在、又已显现给我们的那永生，传报给你们——我们将所见所闻的传报给你们，为使你们也与我们相通；而我们的相通是与父和祂的子耶稣基督相通的。我们给你们写这些事，为使你们的喜乐得以圆满。这就是我们从祂所听见，而传报给你们的讯息：天主是光，在祂内毫无黑暗。如果我们说我们与祂相通，却仍在黑暗中行走，我们就是说谎，不履行真理。如果我们走在光中，如同祂在光中，我们就彼此相通，而祂的子耶稣基督的血必洗净我们一切罪恶。如果我们说我们没有罪，我们就是自欺，真理不在我们内。如果我们明认我们的罪，祂是忠信和公义的，必赦免我们的罪，并洗净我们一切的不义。如果我们说我们没有犯过罪，我们就是以祂为说谎者，祂的话就不在我们内。我的孩子们，我给你们写这些事，是为叫你们不犯罪。但若有人犯了罪，我们在父那里有辩护者，就是正义的耶稣基督；祂是为我们的罪作赎罪祭，不但是为我们的，也是为全世界的。我们若遵守祂的诫命，就晓得我们认识祂。那说\InnerQuote{我认识祂}，却不遵守祂诫命的，是说谎者，真理不在他内。但谁若遵守祂的话，天主的爱就在他内得以圆满。由此我们知道我们在祂内。那说自己住在祂内的，就该照祂所行的去行。亲爱的，我给你们写的不是一条新命令，而是你们从起初就有的旧命令：这旧命令就是你们所听见的话。再者，我给你们写一条新命令，这在祂和你们内是真实的，因为黑暗正在过去，真光已在照耀。那说自己是在光中，却恼恨自己弟兄的，至今仍在黑暗中。凡爱自己弟兄的，就住在光中，在他内没有绊脚石。凡恼恨自己弟兄的，就是在黑暗中，且在黑暗中行走，不知往哪里去，因为黑暗弄瞎了他的眼睛。}\par

1. \BibleQuote{那从起初就有的，我们所听见过，我们亲眼看见过，我们亲手摸过的，论到生命的圣言。} 谁是以手触摸圣言的呢？岂不是因为\BibleQuote{圣言成了血肉，居住在我们中间}？如今这为可被触摸而成了血肉的圣言，始于童贞玛利亚，但圣言并非那时才开始，因为宗徒说：\BibleQuote{那从起初就有的。} 看他的这封书信不是为他的福音作证吗？你们近来听见那里说：\BibleQuote{在起初已有圣言，圣言与天主同在。} \BibleRef{若 1:1} 或许，\BibleQuote{关于生命的圣言}可以被当作某种关于基督的表达，而非当时被手触摸的基督身体。看看接下来什么：\BibleQuote{而生命已显现出来。} 所以基督是\BibleQuote{生命的圣言}。如何显现的呢？因为它是\BibleQuote{从起初就有的}，只是尚未向人显现，却已向天使显现，他们看见它并以它为粮。但圣经说什么？\BibleQuote{人吃了天使的食粮。} 然而\BibleQuote{生命在肉身中显现了}，因为它显示出来，使那只能被心灵看见的，也被眼睛看见，为要医治心灵。因为唯有藉心灵才能看见圣言；但肉身也被肉眼看见。我们有能看见肉身的，却没有能看见圣言的；\BibleQuote{圣言成了血肉}，使我们得以看见，为的是在我们内，那使我们能看见圣言的能力得到医治。\par

2. \BibleQuote{我们看见了，并且作证。} \BibleRef{若一 1:2} 或许有些弟兄不谙希腊文，不知道希腊文里\Quote{作证}一词是什么；然而这是一个被众人广泛使用、在宗教上备受敬重的术语；因为在我们言语中称为\Quote{证人}的，在希腊文里就是\OriginalTerm{殉道者}{martyrs}。如今，有谁没有听说过殉道者？或哪里有一位基督徒，他的口中不是每日都念及殉道者的名字？但愿这名字也常驻心中，好叫我们效法殉道者的苦难，而不是用我们的杯盏迫害他们！那么，\BibleQuote{我们看见了，并且作证}，就相当于说：我们看见了，并且是殉道者。因为他们为了见证他们所看见的，并见证他们从那些看见者所听见的，当这见证本身得罪了它所反对的人时，殉道者便遭受了他们所遭受的一切。殉道者是天主的见证人。天主乐意用人为祂作见证，好叫人也得以有天主作为他们的见证。\BibleQuote{我们看见了，}他说，\BibleQuote{并且作证。}他们在何处看见？在显现中。显现是什么意思？在太阳里，即在这白昼之光里。然而，那位创造太阳的，如何能在太阳里被看见，除非正如\BibleQuote{他在太阳里搭起了帷幕；他自己像新郎走出洞房，又像壮士欢然奔路？} \BibleRef{咏 19:5-6} 祂在太阳之前，创造了太阳，祂在晓星之前，在一切星辰之前，在一切天使之前，是真正的造物主（\BibleQuote{因为万物是由祂造成的，没有祂什么也不能造就} \BibleRef{若 1:3}），为了能被看到太阳的肉眼看见，祂将自己的帷幕搭在太阳里，也就是说，在这白昼之光的显现中显示了祂的肉身：而那位新郎的洞房就是童贞女的子宫，因为在那个童贞的子宫里，两者结合了，新郎与新娘，新郎是圣言，新娘是肉身；因为经上记载：\BibleQuote{二人成为一体} \BibleRef{创 2:24}，主在福音中说：\BibleQuote{因此他们不是两个，而是一体。} \BibleRef{玛 19:6} 依撒意亚也清楚地记得他们是两个：因为当他在基督的位格中说话时说：\BibleQuote{祂在我头上戴上冠冕如新郎，又用装饰品妆饰我如新娘。} \BibleRef{依 61:10} 说话者似乎是一个，却使自己同时成为新郎和新娘；因为\BibleQuote{不是两个，而是一体}，因为\BibleQuote{圣言成了血肉，居住在我们中间。} \BibleRef{若 1:14} 教会与那血肉结合，于是形成了整个基督，即头和身体。\par

3. \BibleQuote{我们作证，并把那与父同在、且已显示给我们的永生传报给你们：}即显示在我们中间；更清楚地说，是显示给我们。\BibleQuote{因此，我们所见所闻的，也传报给你们。} \BibleRef{若一 1:3} 那些人亲眼看见了主亲身在肉身中，亲耳听见了主口中的话，并将它们传报给我们。因此，我们也听见了，但没有看见。那么，我们就不如那些看见和听见的人幸福吗？但他怎么又说：\BibleQuote{好使你们也与我们相通}？那些人看见了，我们没有看见，但我们仍然是相通的，因为我们持守共同的信仰。因为曾有一个人，即使看见了也不相信，还要用手探摸才相信，他说：\BibleQuote{我除非将我的指头探入钉孔，触摸祂的伤痕，我决不相信。} \BibleRef{若 20:25} 祂却一度将自己交付给人的手触摸，而祂永远将自己显示给天使的视线；那位门徒触摸后，喊道：\BibleQuote{我的主，我的天主！} \BibleRef{若 20:28} 因为他触摸了那人，便承认了那神。主为了安慰我们这些如今祂坐在天上、无法用手触摸、只能以信德接近祂的人，对他说：\BibleQuote{因为你看见了，才相信；那些没有看见而相信的，才是有福的。} \BibleRef{若 20:29} 这里描述的是我们，指的是我们。那么，就让主预言要实现的真福在我们身上实现吧；让我们坚定地持守那看不见的，因为那些看见的人已告诉了我们。\BibleQuote{好使你们也}，他说，\BibleQuote{与我们相通。} 与众人相通有什么了不起？不要轻视；请看他又说什么：\BibleQuote{我们的相通是与父和祂的子耶稣基督在一起。我们给你们写这些事，是为叫你们的喜乐得以圆满。} \BibleRef{若一 1:4} 他所说的圆满喜乐，就在那种相通中，在那种爱德中，在那种合一中。\par

4. \BibleQuote{这就是我们从祂所听见，又向你们宣布的消息。}\BibleRef{若一 1:5} 这是什么？那些亲眼见过、亲手摸过生命之圣言的人——祂\Quote{从起初就存在}，并在某段时间内成为可见可触的，天主的独生子。祂为何而来？祂告诉了我们什么新事？祂愿意教导什么？祂为何这样做，使圣言成了血肉，使\Quote{超乎万有的天主}受人的侮辱，忍受被祂亲手所造的手掌掴脸？祂要教导什么？祂要显示什么？祂要宣布什么？让我们听：因为没有诫命的果实，听到基督如何诞生、基督如何受苦的故事，不过是脑中的消遣，而非坚固心智。你听到什么大事？留意你听到的果实是什么。祂要教导什么？宣布什么？听。他说：\BibleQuote{天主是光，在祂内毫无黑暗。} 到现在，他确实提到了光，但他的话是暗晦的：对我们有好处的，是他所提到的这光应该光照我们的心，使我们看见他所说的话。这就是我们所宣布的：\BibleQuote{天主是光，在祂内毫无黑暗。} 谁敢说天主内有黑暗？或者什么是光？什么是黑暗？恐怕他是说我们肉眼所涉及的东西。\BibleQuote{天主是光。}有人说：\Quote{太阳也是光，月亮也是光，蜡烛也是光。} 它应该是远超过这些、更卓越、更超绝的东西。天主与受造物之间的距离有多大，造物者与被造物之间的距离有多大，智慧与被智慧所造之物之间的距离有多大，这光就必然远超过万物。或许我们将会接近它，如果我们得知这光是什么，并致力于它，好能藉着它被光照；因为在我们自己内我们是黑暗，只有被它光照才能成为光，并且不被它蒙羞，反而被自己蒙羞。谁是被自己蒙羞的人？就是那知道自己是个罪人的人。谁是不被它蒙羞的人？就是那被它光照的人。被它光照是什么意思？那现在看见自己被罪所蒙蔽，并渴望被它光照的人，就接近它：因此圣咏上说：\BibleQuote{你们亲近祂，就必被光照；你们的脸上必不蒙羞。} 但你不会被它蒙羞，如果当它向你显示你是污秽的，你自己的污秽使你厌恶，好让你能感知它的美丽。这就是祂要教导的。\par

5. 但愿我们这样说话不算过于仓促？让宗徒自己在接下来的话中说明这一点。请记住我们开讲时所说的话：这封书信推崇爱德。他说：\Quote{天主是光，在祂内毫无黑暗。}上文又说了什么？\Quote{好使你们与我们相通，而我们的相通是与天主圣父及祂的圣子耶稣基督相通。}然而，如果\Quote{天主是光，在祂内毫无黑暗，并且我们\InnerQuote{必须}与祂相通}，那么黑暗也必须从我们这里被驱除，好在我们内造光，因为黑暗不能与光相通。为此，请看下文：\Quote{如果我们说我们与祂相通，却仍在黑暗中行走，我们就说谎。}\BibleRef{若一 1:6} 你也听到了宗徒保禄说：\Quote{光明与黑暗怎能相通？}\BibleRef{格后 6:14} 你说你与天主相通，却仍在黑暗中行走；\Quote{天主是光，在祂内毫无黑暗}：那么光明与黑暗如何能相通呢？因此，一个人可能会对自己说：我该怎么办？我怎能成为光？我生活在罪恶与不义之中。他心中仿佛涌起绝望与悲伤。除了与天主相通，别无救恩。\Quote{天主是光，在祂内毫无黑暗。}但罪恶是黑暗，正如宗徒论及魔鬼及其使者时所说的，他们是\Quote{这黑暗世界的统治者}。\BibleRef{弗 6:12} 他称他们为黑暗的，除非是指罪恶的统治者，统治着恶人。那么，我的弟兄们，我们该怎么办？必须与天主相通，否则没有永生的希望；如今\Quote{天主是光，在祂内毫无黑暗}；如今不义是黑暗；因不义我们被压制，不能与天主相通：那我们还有什么希望？我不是承诺在这几天里要说些令人喜乐的话吗？如果我不实现这个承诺，这将是悲伤。\Quote{天主是光，在祂内毫无黑暗}；罪恶是黑暗：我们会怎样？让我们听听，或许祂会安慰、提升、给予希望，使我们不致在路上灰心。因为我们正在奔跑，奔向我们的家乡；如果我们对抵达绝望，就会因这绝望而失败。但是那愿意我们抵达、愿意我们在故土得救的那一位，在路上喂养我们。那么请听：\Quote{如果我们说我们与祂相通，却仍在黑暗中行走，我们就说谎，不履行真理。}如果我们仍在黑暗中行走，就不要说我们与祂相通。\Quote{如果我们像祂一样在光明中行走，我们就彼此相通。}\BibleRef{若一 1:7} 让我们在光明中行走，如同祂在光明中一样，以便能与祂相通。那么，关于我们的罪，我们该做什么？听下文：\Quote{祂的圣子耶稣基督的血要洗净我们的一切罪。}天主赐予了极大的确据！我们欢庆逾越节，在其中倾流了主的血，这血洗净我们\Quote{脱离一切罪}！让我们确信：\Quote{那反对我们的债卷，}\BibleRef{哥 2:14} 我们奴役的契约，曾被魔鬼持有，但藉基督的血已被涂抹。\Quote{祂的圣子的血}，他说，\Quote{要洗净我们的一切罪。}\Quote{一切罪}是什么意思？请注意：就在此刻，在基督的名下，这些刚刚宣认基督、被称为婴孩的人，他们所有的罪都被洁净了。他们进来时是旧的，出去时是新的。如何进来是旧的，出去是新的？他们进来时是老人，出去时是婴孩。因为旧的生命是衰老及其昏聩，但新的生命是重生的幼年。但我们该怎么办？过去的罪已被赦免，不仅是对这些人，也是对我们；在赦免并废除了所有罪之后，由于活在这个充满诱惑的世界上，也许又沾染了一些。因此，人能做多少就做多少；让他承认自己是什么，好被那永远是什么的那一位治愈：因为祂永远是，且现在是；我们曾不是，现在则是。\par

6. 请看祂说什么；\BibleQuote{如果我们说我们没有罪，就是欺骗自己，真理也不在我们内。}\BibleRef{若一 1:8} 所以，如果你承认自己是罪人，真理就在你内：因为真理本身就是光。你的生命尚未完美地闪耀，因为你有罪；但你已经开始被光照，因为你内心有对罪的承认。请看接下来的话：\BibleQuote{如果我们承认我们的罪，祂是忠信和公义的，必赦免我们的罪，并洗净我们一切的不义。}\BibleRef{若一 1:9} 不仅是过去的罪，也许还有我们从这生活中沾染的；因为人，只要带着肉身，就不可能没有一些轻微的罪。但这些我们称为轻微的，不要轻视它们。如果你在称量时轻视它们，在数算时就要害怕。许多轻微的小罪构成一个大罪；许多水滴汇成江河；许多谷粒集成大堆。那么有什么希望呢？首先是认罪：免得有人自以为义，在天主面前——那观看实情者——人（原本没有，现在存在）昂首挺胸。所以首要的是认罪；然后是爱：因为关于爱德，经上说什么？\BibleQuote{爱德遮盖许多罪过。}\BibleRef{伯前 4:8} 现在让我们看看祂是否在论及随后加给我们的罪过时称赞爱德：因为只有爱德能消灭罪过。骄傲消灭爱德，所以谦卑坚固爱德；爱德消灭罪过。谦卑与认罪同行，就是那种我们承认自己是罪人的谦卑：这才是谦卑，不是用舌头说，好像只是为了避免骄傲，以免我们如果自称为义会得罪人。不虔诚和疯狂的人这样做：\Quote{我确实知道我是义人，但我在人前该说什么？如果我自称义人，谁能忍受？谁会容忍？让我的义德在天主面前显明吧；我却要说我是罪人，只是为了不因骄傲而被人厌恶。} 告诉人你是什么，告诉天主你是什么。因为如果你不告诉天主你是什么，天主会谴责祂在你内发现的。你不想祂谴责你吗？你就谴责自己。你想祂宽恕吗？你就承认，好使你能向天主说：\Quote{求祢转面不看我的罪过。} 也要向祂说同一圣咏中的话：\Quote{因为我承认我的过犯。} \BibleQuote{如果我们承认我们的罪，祂是忠信和公义的，必赦免我们的罪，并洗净我们一切的不义。如果我们说我们没有犯过罪，我们便是以祂为说谎者，祂的话就不在我们内。}\BibleRef{若一 1:9} 如果你说\Quote{我没有犯罪}，你便以天主为说谎者，而你想使自己真实。怎么可能天主是撒谎者，而人却是真实的，而圣经却相反说：\BibleQuote{人人都是撒谎者，唯独天主是真实的？}\BibleRef{罗 3:4} 因此，天主藉自己真实，你藉天主真实；因为藉你自己，你是撒谎者。\par

7. 恐怕祂似乎因为说了\Quote{祂是忠信和公义的，必洗净我们一切的不义}，而给人犯罪的免罚权；于是人们就会对自己说：让我们犯罪，我们想做什么就安全地做吧，基督洗净我们，祂是忠信和公义的，洗净我们一切的不义：祂从你那里拿走一种邪恶的安全感，放入一种有益的恐惧。你要安全反而对你有害；你必须警醒。因为\Quote{祂是忠信和公义的，必赦免我们的罪}，前提是你总是对自己不满，并不断改变直到你得以成全。那么接下来呢？\BibleQuote{我的孩子们，我写这些给你们，是要你们不犯罪。}\BibleRef{若一 2:1} 但或许罪从我们的必死生命中胜过我们：那时该怎么办？什么？难道现在要绝望吗？听：\BibleQuote{如果有人犯了罪，我们在父那里有一位护慰者，就是那义者耶稣基督：祂自己就是为我们的罪作补赎的。}\BibleRef{若一 2:1} 祂就是护慰者；你尽力不犯罪：如果因这生命的软弱罪胜过你，要立刻留意，立刻不高兴，立刻谴责它；当你谴责了，你就可以放心地来到审判者面前。你那里有护慰者：不要怕在你的认罪中输掉你的案子。因为如果在今生，一个人常把自己的案子交给一个雄辩的口才，就没有输掉；你把自己交托给圣言，你会输掉吗？呼喊：\Quote{我们在父那里有一位护慰者。}\par

8. 请看若望本人如何持守谦逊。他确是一位义人、一位伟人，从主的怀里汲取了祂奥秘的秘密；他，这位从主的怀里汲取而论述祂的天主性的人，说：\Quote{在起初已有圣言，圣言与天主同在。}他既是如此伟大的人，却不曾说：\Quote{你们在父那里有一位辩护者}，而是说：\Quote{谁若犯了罪，我们有一位辩护者。}他不说\Quote{你们有}，也不说\Quote{你们有我}，也不说\Quote{你们有基督本人}，而是把基督放在前面，不把自己放在前面，还说\Quote{我们}有，而非\Quote{你们}有。他宁可把自己列在罪人之中，好让基督作他的辩护者，而不愿把自己放在基督的位置上作辩护者，以致被列入骄傲的行列而受审判。弟兄们，耶稣基督——那义者——就是我们在父那里的辩护者；\Quote{祂}自己\Quote{就是我们罪的赎罪祭}。谁持守这一点，就不会制造异端；谁持守这一点，就不会制造分裂。因为分裂从何而来？当人们说\Quote{我们}是义人，当人们说\Quote{我们}圣化不洁者，\Quote{我们}使不敬虔者成义；\Quote{我们}祈求，\Quote{我们}获得。但若望怎么说？\Quote{谁若犯了罪，我们在父那里有一位辩护者，就是那义者耶稣基督。}可是有人会说：那么圣人们不为我们祈求吗？那么主教和领袖们不为民众祈求吗？是的，但请看圣经，看领袖们也将自己托付于民众的祈祷。因此宗徒对会众说：\Quote{也要为我们祈求。}\BibleRef{哥 4:3}宗徒为民众祈祷，民众为宗徒祈祷。我们为你们祈祷，弟兄们；但你们也要为我们祈祷。让所有的肢体彼此祈祷；让头为众人转求。因此，他接着堵住了那些分裂天主的教会之人的口，这并不奇怪。因为他说了：\Quote{我们有一位辩护者耶稣基督——那义者，祂就是我们罪的赎罪祭}；他注意到那些要分裂自己的人，他们会说：\Quote{看，基督在这里，看，在那里。}\BibleRef{玛 24:23}他们想把那位赎买了整个并拥有整个的基督局限于一部分。于是他接着说：\Quote{不是单为我们的罪，也是为普世的罪。}这是什么意思，弟兄们？的确，\Quote{我们在林间的田野找到了}；我们在万民中找到了教会。看，基督\Quote{是我们罪的赎罪祭；不仅为我们的，也为普世的罪。}看，你有了遍布全世界的教会；不要追随那些虚假的成义者，他们实则是使人分裂的。要留在这座充满大地的山上，因为\Quote{基督是我们罪的赎罪祭；不仅为我们的，也为普世的罪}，他用祂的血赎买了整个世界。\par

9. 他说：\Quote{由此我们知道我们认识祂，如果我们遵守祂的诫命。}\BibleRef{若一 2:3}什么诫命？\Quote{那说\InnerQuote{我认识祂}，却不遵守祂诫命的，是撒谎的人，真理不在他内。}但你还要问：什么诫命？他说：\Quote{可是，谁若遵守祂的话，天主的爱在他内就得以圆满。}\BibleRef{若一 2:5}我们来看看，这诫命是否被称为爱？因为我们问\Quote{什么诫命}，他说：\Quote{可是，谁若遵守祂的话，天主的爱在他内就得以圆满。}请看福音，看看这诫命是否是这样：\Quote{我给你们一条新诫命}，主说，\Quote{叫你们彼此相爱。}\BibleRef{若 13:34}\Quote{由此我们知道我们是在祂内，如果我们在祂内得以圆满。}他称他们为在爱中得以圆满的人。什么是爱的圆满？就是连仇敌也爱，并且爱他们是为了使他们成为弟兄。因为我们的爱不该是血肉之爱。愿一个人得现世的好处，这是好的；但即便得不到，但愿灵魂得救。你愿你的朋友得生命吗？你做得对。你因你的仇敌死亡而欢喜吗？你做得不对。但或许你愿朋友得的生命对他无益，而你为仇敌之死欢喜的事对他有益。这现世生命对任何人是否有益，是不确定的；但与天主同在的生命无疑是有益的。所以，你要爱你的仇敌，愿意他们成为你的弟兄；你要爱你的仇敌，愿意他们被召叫进入你的团契。因为那位挂在十字架上说\Quote{父啊，宽恕他们，因为他们不知道他们在做什么}的人就是这样愛的。\BibleRef{路 23:34}祂并没有说：\Quote{父啊，让他们长寿，他们固然杀我，但让他们活着。}祂以祂最仁慈的祈祷和祂那超越一切的大能，从他们身上驱逐那永永远远的死亡。他们中有许多人相信了，基督流血的事被宽恕了。起初他们狂怒时流了祂的血；如今他们相信时饮了祂的血。\Quote{由此我们知道我们是在祂内，如果我们在祂内得以成全。}主提醒有关爱仇敌的圆满：\Quote{所以你们应当是成全的，如同你们的天父是成全的一样。}\BibleRef{玛 5:48}\Quote{所以，那说住在祂内的，便应该照祂所行的去行。}\BibleRef{若一 2:6}弟兄们，这是什么意思？他怎么劝诫我们？\Quote{那说住在祂内的}——即在基督内的——\Quote{便应该照祂所行的去行。}或许这个劝诫是要我们走在水面上？决不是！所以，这是要我们行走在义的道路上。哪条路？我已经提过。祂被钉在十字架上，却仍在走这条路；这条路就是爱德之路：\Quote{父啊，宽恕他们，因为他们不知道他们在做什么。}因此，如果你学会了为你的仇敌祈祷，你就在走主的道路。\par

10. \Quote{亲爱的诸位，我给你们写的不是一条新诫命，而是你们从起初就有的旧诫命。}\BibleRef{若一 2:7} 那诫命他为什么称为\Quote{旧的？你们从起初就有的，}他说，\Quote{从起初。既是旧的}，就是因为你们已经听过了；否则他将与主矛盾，因为主说：\Quote{我给你们一条新诫命，就是你们要彼此相爱。}\BibleRef{若 13:34} 但为什么称为\Quote{旧的}诫命？不是属于旧人。但为什么？\Quote{你们从起初就有的。那旧诫命就是你们所听过的道理。}既是旧的，就是因为你们已经听过了。而他也表明同一条诫命是新的，说：\Quote{再者，我给你们写一条新诫命。}\BibleRef{若一 2:8} 不是另一条，而是同一条他称为旧的，也是新的。为什么？\Quote{这在祂和你们身上都是真的。}为什么是旧的，你们已经听见了：\Emph{即}因为你们早已知道。但为什么是新的？\Quote{因为黑暗已过，真光正在照耀。}看哪，它之所以是新的，因为黑暗属于旧人，而光属于新人。保禄宗徒怎么说？\Quote{脱去旧人，穿上新人。}\BibleRef{哥 3:9} 他又怎么说？\Quote{你们从前是黑暗，但现在在主内是光。}\BibleRef{弗 5:8}\par

11. \Quote{那说自己是在光明中的}——现在他是在说明他先前所说的一切——\Quote{那说自己是在光明中，却恼恨自己弟兄的，直到现在还是在黑暗里。}\BibleRef{若一 2:9} 怎么！我的弟兄们，我们还要向你们说多久\Quote{爱你们的仇敌}\BibleRef{玛 5:44}？看看，还有更糟的，你们是否恼恨自己的弟兄？如果你们只爱自己的弟兄，你们还不够完美；但如果你们恼恨自己的弟兄，你们是什么，你们在哪里？愿每个人查看自己的心：不要因任何难听的话而对弟兄怀恨；不要因地上的争执而变成土。因为凡恼恨自己弟兄的，不要说他行在光明中。\Quote{那说自己是在光明中，却恼恨自己弟兄的，直到现在还是在黑暗里。}\BibleRef{若一 2:9} 就这样，一个人曾是外邦人，成了基督徒；请注意：看，他在身为外邦人时是在黑暗里；现在他成了基督徒；感谢天主，大家欢欣地说；宗徒的书信被诵读，那里他喜悦地说：\Quote{你们从前是黑暗，但现在在主内是光。}\BibleRef{弗 5:8} 从前他崇拜偶像，如今他崇拜天主；从前他崇拜自己所造之物，如今他崇拜造他的主。他改变了：感谢天主，所有基督徒都欢喜地向他问候。为什么？因为从今以后他是朝拜圣父、圣子和圣神的人；憎恨魔鬼和偶像的人。然而若望仍然担忧我们的皈依者：当许多人欢喜地迎接他时，若望仍以忧虑的眼光看他。弟兄们，让我们欣然接受一位母亲的焦虑。这母亲不是无缘无故在他人欢乐时为我们忧虑：我所说的母亲，是指爱德；因为当若望说这些话时，爱德居住在他心中。为什么如此？岂不是因为在我们身上有他所惧怕的东西，即使现在人们欢喜地向我们致意？他惧怕的是什么？\Quote{那说自己是在光明中的}——这是什么？那说自己现在是基督徒的——\Quote{却恼恨自己的弟兄，直到现在还是在黑暗里。}这无需解释：如果事实不是这样，我们为之欢喜；如果是这样，我们为之哀悼。\par

12. \BibleQuote{爱弟兄的，就住在（\OriginalTerm{住}{manet}）光明中，在他里面毫无绊脚石。} \BibleRef{若一 2:10} — 我以基督恳求你们：天主在喂养我们，我们即将奉基督的名恢复身体的精力；它们已经得到某种程度的恢复，并将继续恢复：但愿心灵也得喂养。我说这话，并非因为我要讲很久；看，读经即将结束：但恐怕因疲乏，我们聆听那最必要的话时不如应有的专注。—— \BibleQuote{爱弟兄的，就住在光明中，在他里面没有绊脚石，} 或 \BibleQuote{在他里面毫无绊倒的缘由。} 谁是那被绊倒或使人绊倒的人？就是在基督内和在教会内跌倒的人。在基督内跌倒的人，如同被太阳灼伤；在教会内跌倒的人，如同被月亮灼伤。但圣咏说：\BibleQuote{白日太阳不烧伤你，夜间月亮也不害你：} \Emph{即}，你若持守爱德，那么在基督内你必无跌倒的缘由，在教会内也必无；你既不会离弃基督，也不会离弃教会。因为那离弃教会的，怎能是在基督内呢？他既不在基督的肢体里，怎能是在基督内呢？所以，那离弃基督或教会的，就被绊倒，或有了跌倒的缘由。我们如何理解圣咏说\BibleQuote{白日太阳不烧伤你，夜间月亮也不害你} 时，指的是这意思，即灼伤就是绊倒或跌倒的缘由？首先注意这比喻本身。正如一个人被某物灼伤时说：我受不了，我忍不下去，便退后；同样，那些受不了教会中某些事，便从基督的名或教会中退出的人，就是被绊倒。因为请看那些属血肉的人，他们从太阳般（的基督）那里被绊倒，基督曾论及祂的肉躯说：\BibleQuote{谁不吃人子的肉，不喝祂的血，在他内就没有生命。} \BibleRef{若 6:54} 约有七十人说：\BibleQuote{这话生硬，} 便离开祂，只有十二人留下。所有这些人都被太阳灼伤，因受不了圣言的力量而离去。于是留下了十二人。唯恐人们以为他们相信基督是给基督带来好处，而非基督将好处赐予他们；当十二人留下时，主对他们说：\BibleQuote{你们也要去吗？} 为叫你们知道，我是你们必需的，你们却不是我所必需的。但那些没有被太阳灼伤的人，以伯多禄的声音回答：\BibleQuote{主，你有永生的话，我们去投奔谁呢？} 那么，那些被教会如同月亮在夜间灼伤的人是谁？就是制造分裂的人。请听宗徒所用的词：\BibleQuote{谁软弱，我不软弱呢？谁跌倒，我不焦急呢？} \BibleRef{格后 11:29} 那么，在爱弟兄的人身上毫无绊脚石或跌倒的缘由，是什么意思呢？因为爱弟兄的，为了合一而忍受一切；因为弟兄之爱是在爱德的合一中存在的。我不知道某个人得罪了你：也许他是个坏人，或者你以为是坏人，或者你假装是坏人：难道你就因此离弃这么多好人吗？那在这些人身上显示出来的弟兄之爱是怎样的呢？他们指责阿非利加人，却离弃了全世界！怎么，全世界就没有圣人吗？或者，他们未经你审问就被定罪吗？但唉！你若爱你的弟兄们，在你内就毫无跌倒的缘由。请听圣咏所说：\BibleQuote{爱慕你法律的人，必得大平安，在他们内毫无绊脚石。} 它说，爱慕天主法律的人必得大平安，所以他们内毫无绊脚石。因此，那些被绊倒或跌倒的人，毁灭了平安。它说哪些人不被绊倒也不使人绊倒？就是爱慕天主法律的人。所以他们是在爱德中。但有人会说：\BibleQuote{这是论及爱慕天主法律的人，而不是论及弟兄们。} 请听主所说：\BibleQuote{我给你们一条新诫命：你们要彼此相爱。} \BibleRef{若 13:34} 法律不就是诫命吗？再者，他们为何不跌倒？岂不是因为他们彼此容忍吗？正如保禄所说：\BibleQuote{彼此以爱德容忍，竭力保持圣神的合一，以和平的联系相连。} \BibleRef{弗 4:2} 为表明这就是基督的法律，请听同一位宗徒称赞这法律。他说：\BibleQuote{你们要彼此担当重担，} \BibleQuote{这样就满了基督的法律。} \BibleRef{迦 6:2}\par

13. \BibleQuote{凡恼恨自己弟兄的，就是在黑暗中，且行在黑暗中，不知道往哪里去。}\BibleRef{若一 2:11} 这是一件大事，我的弟兄们：请注意，我们恳求你们。\BibleQuote{凡恼恨自己弟兄的，就行在黑暗中，不知道往哪里去，因为黑暗蒙蔽了他的眼睛。} 还有什么比这些恼恨自己弟兄的人更瞎眼的呢？为使你们知道他们是瞎眼的，他们竟绊倒在一座山上。我常常重复这些话，免得它们从你们的记忆中溜走。那\Quote{非人手从山中凿出的}石头，岂不就是基督，祂出自犹太国，不是出于人的作为吗？那石头不是打碎了地上所有的王国，即一切偶像和魔鬼的统治吗？那石头不是长大，变成一座大山，充满全地吗？我们指着这座山，就像人们指出第三天的月亮那样吗？例如，当他们希望人们看到新月时，他们说：\Quote{看，月亮！看，它在那里！}如果那里有一些不够敏锐的人说：\Quote{在哪里？}于是就伸出手指让他们看见。有时他们羞于被认为瞎眼，就说他们看见了其实没看见的东西。我们是这样指出教会吗，我的弟兄们？它不是敞开的吗？它不是明显的吗？它不是已占据了万民吗？那么多年以前应许给亚巴郎的话，即\BibleQuote{万民将因他的后裔蒙受祝福}\BibleRef{创 22:18}，这岂不已经应验了吗？这应许是给一个信徒的，而世界充满了千千万万的信徒。看，这里就是那座充满全地面的大山！看，那座城，经上论它说：\BibleQuote{建在山上的城，是不能隐藏的！}\BibleRef{玛 5:14} 但那些人却在那山上绊倒，当有人对他们说：\Quote{上去}；他们却说：\Quote{没有山}，并且宁可撞破头也不愿在那里寻求住处。昨天诵读了\Person{依撒意亚}；你们中谁若不仅以眼睛醒着，而且以耳朵醒着，不仅是肉身的耳朵，更是心灵的耳朵，就注意到这一点：\BibleQuote{在末日，上主殿宇的山岳必耸立，卓然超越群山。}\BibleRef{依 2:2} 还有什么比一座山更明显的呢？但有些山甚至不为人知，因为它们位于地球的一个地方。你们谁认识\Place{奥林帕斯山}？正如住在那里的居民不认识我们的\Place{吉达巴山}一样。这些山在地球的不同地方。但那座山却不同，因为它充满了全地面，关于它经上说：\Quote{卓然超越群山。}它是一座高耸于所有山顶之上的山。\Quote{还有，}他说，\Quote{万民都要汇流到它那里。}谁能不认识这座山呢？谁会在上面绊倒而撞破头呢？谁不知道这座建在山上的城呢？但不要惊讶它不为那些恼恨弟兄的人所知，因为他们行在黑暗中，不知道往哪里去，因为黑暗蒙蔽了他们的眼睛。他们看不见那座山：我不愿你们惊讶；他们没有眼睛。他们怎么没有眼睛呢？因为黑暗蒙蔽了他们。我们如何证明这一点呢？因为他们恼恨弟兄，因他们见怪非洲人，就自绝于普世教会；因他们为基督的平安不能容忍那些他们诽谤的人，却为\Person{多纳图}的缘故容忍那些他们定断的人。\par

\SourceRecord{Translated by H. Browne.；Nicene and Post-Nicene Fathers, First Series；Vol. 7.；Edited by Philip Schaff.；Buffalo, NY: Christian Literature Publishing Co.,；1888.；<http://www.newadvent.org/fathers/170201.htm>.}

% structure: p0001 source=h1 target=section reason=page-title

\section{\WorkTitle{若望一书}第二篇讲道}

% structure: p0002 source=h2 target=subsection reason=relative-heading-rank; base=h2

\subsection{若望一书 2:12-17}

\BibleQuote{小子们，我给你们写，因为你们的罪因祂的名已获得赦免。父老们，我给你们写，因为你们认识自始就有的那一位。青年们，我给你们写，因为你们战胜了恶者。孩子们，我给你们写，因为你们认识了父。父老们，我给你们写，因为你们认识自始就有的那一位。青年们，我给你们写，因为你们刚强，天主的言语存在于你们内，你们战胜了恶者。不要爱世界，也不可爱世界上的事物。若有人爱世界，圣父的爱就不在他内。原来世界上的一切：肉身的贪欲、眼目的贪欲，以及生活的骄傲，都不是出于圣父，而是出于世界。世界和它的贪欲都要过去；但那履行天主旨意的，却永远存在（正如天主也永远存在）。}\par

1. 凡为教训和救恩而从圣经诵读的一切，我们都当热切聆听。但最要牢记于心的是那些最具抵挡异端力量的事，因为异端者的阴险诡计无休止地欺骗那些较为软弱和疏忽的人。要记住，我们的主和救主耶稣基督为我们死了，又复活了；祂是为我们的过犯而死，为我们的成义而复活。\BibleRef{罗 4:25} 正如你们刚才听到的，关于祂在路上遇到的两个门徒，如何\BibleQuote{他们的眼睛被蒙蔽，以致认不出祂来：}\BibleRef{路 24:13} 祂发现他们对在基督内的救赎感到绝望，认为祂已经受苦，像凡人一样死去，不承认祂作为天主之子永远活着；也认为祂在肉身内死了，不能复活，如同一位先知一样：正如你们中留心的人刚才听到他们自己的话。然后\BibleQuote{祂从梅瑟开始，为他们开启圣经，}并遍历所有先知，向他们显示祂所受的一切都曾被预言，免得他们如果主复活了更加惊慌，且如果这些事未曾提前晓谕，他们更难相信祂。因为信德的坚固在于：凡在基督身上发生的事，都已预报过。门徒们不认识祂，除非\BibleQuote{在擘饼中。}确实，那在擘饼中不自己吃喝审判的人，就认识基督。\BibleRef{格前 11:29} 后来那十一人也\BibleQuote{以为看见了鬼魂。}祂将自己交付他们触摸，祂也曾将自己交付钉十字架；被敌人钉十字架，被朋友触摸；但祂是万有的医师，既治那些人的不虔，也治这些人的不信。当你们听到\WorkTitle{宗徒大事录}被诵读时，你们听到杀害基督的人中，有数千人相信了。\BibleRef{宗 2:41} 如果那些杀害者后来相信了，那些暂时疑惑的人岂不更应该相信吗？然而，即使关于他们（这事你们尤其应当留意并牢记于心，因为那将使我们坚强抵御阴险错误的事，天主乐意记载于圣经中，任何人只要稍愿显为基督徒，都不敢反对圣经），当祂将自己交付他们触摸之后，祂仍不满足，更愿借圣经来坚定信者的心：因为祂预见后来要有的我们；既然我们在祂身上没有可触摸的，却有可读的。如果那些人只是因为握住并触摸了才相信，我们该怎么办？如今，基督已升天，除非最后审判活人死人，祂不会再来。我们凭什么相信呢？岂不是凭祂愿意连那些触摸祂的人也得坚固的凭借吗？因为祂为他们开启圣经，并指示他们：基督必须受苦，凡在梅瑟法律、先知书和圣咏中关于祂所记载的一切，都当应验。祂在讲论中涵盖了全部古老圣经文本。那先前圣经中所有的一切，都在讲述基督；但只当它寻得耳朵时。祂也\BibleQuote{开启他们的理智，使他们理解圣经。}因此我们也必须为此祈祷，求祂开启我们的理智。\par

2. 但主在梅瑟法律、先知和圣咏中，关于祂自己写了什么？祂展示了什么？让祂亲自说吧。福音作者简要地记述了这一点，好让我们知道，在如此浩瀚的圣经中，我们应当相信并理解什么。的确，有许多书页和许多书卷；它们所有的内容，就是主对祂的门徒简要说的这一点。这是什么？就是\BibleQuote{基督必须受苦，第三日从死者中复活。}现在你听到了关于新郎的话，就是\BibleQuote{基督必须受苦并复活}：新郎已展示给我们。关于新娘，让我们看看祂说什么；这样，当你认识新郎和新娘时，就不会无缘无故地来参加婚宴。因为每一次庆典都是婚姻的庆典：教会正在举行婚礼。国王的儿子将要娶妻，而这位国王的儿子本身就是一位国王；参加婚礼的客人正是新娘本身。不像在肉体的婚姻中，有客人，另有一位是新妇；在教会中，那些前来参加婚礼的人，如果他们怀着善意而来，就成为了新娘。因为整个教会都是基督的新娘，其开端和初果是基督的肉身：在肉身中，新娘与新郎结合了。很有理由，当祂预示这同一肉身时，祂擘开了饼，并且\BibleQuote{在擘饼时，门徒的眼睛开了，就认出了祂。}好了，那么，主说了什么关于祂自己写在法律、先知和圣咏中的事？就是\BibleQuote{基督必须受苦。}如果祂没有再加上\BibleQuote{并复活}，那么那些眼睛被蒙蔽的人就理应哀哭；但\BibleQuote{复活}也预言了。这是为什么？为何基督必须受苦并复活？因为那首我们在上周第四日第一站特别向你们推荐的圣咏。为何基督必须受苦并复活？因为这个原因：\BibleQuote{大地四极将想起上主，归向祂；万邦的家族将朝拜在祂面前。}为了使你知道基督必须受苦并复活，在此处祂又加上了什么，以便在展示新郎之后也展示新娘？祂说：\BibleQuote{并且要因祂的名，从耶路撒冷开始，向万邦宣讲悔改和赦罪。}弟兄们，你们听到了，要牢记。任何人都不要怀疑教会是\BibleQuote{遍布万邦}：不要怀疑它从耶路撒冷开始，已经充满了万邦。我们知道葡萄树栽种的田地；但当它长大时，我们却不认识，因为它占据了全部。它从何处开始？\BibleQuote{从耶路撒冷。}它到了何处？到了\BibleQuote{万邦。}还有少数地方；它将占据全部。在此期间，当它正在占据全部时，园主似乎砍掉了一些不结实的枝条，它们便成了异端和分裂。不要让被砍掉的枝条引诱你也砍掉；反而要劝诫那些被砍掉的，好让他们重新接上。显然，基督已经受苦、复活并升天了；教会也显明了，因为\BibleQuote{因祂的名，从耶路撒冷开始，向万邦宣讲悔改和赦罪。}它从何处开始？\BibleQuote{从耶路撒冷开始。}人听到这些；愚拙而虚妄，并且（我该怎么说？）比瞎子更糟！如此大的山，他看不见；灯放在灯台上，他却闭眼不看！\par

3. 当我们对他们说：\Quote{如果你们是公教信徒，就与那教会共融，福音正是从那里传遍普世；与那\Place{耶路撒冷}共融。}当我们这样对他们说时，他们回答我们说：\Quote{我们不与那城共融，因为我们的王在那里被杀，我们的主在那里被杀。}好像他们憎恨那杀了我们主的城。犹太人杀死了他们在世上找到的祂，这些人嘲弄坐在天上的祂！谁更恶劣？是那些因以为祂是人而鄙视祂的人，还是那些嘲弄如今承认祂是天主的祂的圣事的人？但他们确实憎恨那杀了他们主的城！虔诚而慈悲的人！他们非常忧伤基督被杀，却在他们内杀害基督！但祂爱那城，并怜悯它：祂吩咐从那里开始传讲祂，\BibleQuote{从耶路撒冷开始。}祂在那里开始了宣讲祂名字的工作：而你们竟惊恐地退缩，不愿与那城共融！难怪被剪除的你们憎恨这根源。祂对祂的门徒说了什么？\BibleQuote{留在城里，直到我把我的恩许派遣到你们身上。}看，他们所憎恨的是怎样的城！也许如果杀害基督的人住在那里，他们倒会爱上它。因为显而易见，所有杀害基督的人，即犹太人，已被从那城驱逐出去了。那曾有凶暴反对基督之人的地方，如今有了朝拜基督的人。因此这些人憎恨它，因为基督徒在那里。祂的旨意是要祂的门徒留在那里，并在那里把圣神派遣给他们。教会始于何处？岂不正是始于圣神从天降下、充满聚集在一处的百二十人的地方？那十二的数变成了十倍。他们一百二十人坐着，圣神降临，\BibleQuote{充满了整个地方，有声音好像一阵强劲的风吹过，又有舌头像火一样分开。}你们已经听到了\WorkTitle{宗徒大事录}：这是今日所读的经文：\BibleQuote{他们开始说各种语言，按圣神所赐给他们的话讲论。}所有在场的人，即来自各民族的犹太人，都认出自己的语言，惊奇那些没有学问、无知的人一下子学会了不是一两种语言，而是所有民族的各种语言。所以，在那里，当所有语言响起时，就预示着所有语言都要相信。但是这些人，他们非常爱基督，因此拒绝与杀害基督的城共融，他们如此尊敬基督，以至于断言基督只局限于两种语言：拉丁语和\OriginalTerm{布匿语}{Punic}，即非洲语。基督只拥有两种语言！因为在\OriginalTerm{多纳徒}{Donatus}一边只有这两种语言，他们不再有别的了。我的弟兄们，让我们醒悟，让我们更看看天主的圣神的恩赐，让我们相信先前关于祂所说的话，并让我们看到在\WorkTitle{圣咏}中预先所说的话得以应验：\BibleQuote{没有语言，没有话语，但他们的声音在他们中间可被听见。}免得情况是这样：语言本身来到一处，而不是基督的恩赐来到所有语言中，请听接下来所说的：\BibleQuote{他们的声音传遍全地，他们的话语传到地极。}这是为何？因为\BibleQuote{祂在太阳中设置了祂的帐幕}，\Emph{i.e.}，在光明中。祂的帐幕，祂的肉身；祂的帐幕，祂的教会：它被设置在\BibleQuote{太阳中}，不是在夜里，而是在白昼。但为什么那些人认不出它呢？回到昨日读经结束的地方，看看他们为何认不出：\BibleQuote{那恨自己兄弟的，是在黑暗中行走，不知道往哪里去，因为黑暗弄瞎了他的眼睛。}因此，让我们看看接下来有什么，不要留在黑暗中。我们怎样才能不留在黑暗中？如果我们爱弟兄。我们如何证明爱弟兄的团体？凭这一点：我们不撕裂合一，我们持守爱德。\par

4. \BibleQuote{小子们，我写信给你们，因为你们的罪因祂的名得了赦免。} \BibleRef{若一 2:12} 因此，\BibleQuote{小子们}，因为在罪的赦免中你们获得了重生。但罪是因谁的名而得赦免？是因\OriginalTerm{奥斯定}{Augustine}的名吗？不，因此也不是因\OriginalTerm{多纳徒}{Donatus}的名。你们要留意看清奥斯定是谁，多纳徒是谁：不，也不是因保禄的名，也不是因伯多禄的名。因为对于那些把教会分给自己、试图从合一中制造党派的人，宗徒中的母亲爱德，为其幼小者生产之苦，袒露自己的心肠，以言语似乎撕裂自己的胸，哀哭她所见已死去并生出的孩子，召叫那些需要许多名字的人归于一个名字，把他们从对她的爱中推开，使基督受到爱，她说：\BibleQuote{保禄曾为你们被钉十字架吗？或者你们是因保禄的名受洗吗？} \BibleRef{格前 1:13} 他说了什么？\Quote{我不愿你们属于我，好使你们能与我同在：与我同在；我们全是祂的，祂为我们死了，为我们被钉在十字架上。}因此这里又说：\BibleQuote{你们的罪因祂的名得了赦免，}不是因任何人的名。\par

5. \BibleQuote{我写信给你们，父老们。}\BibleRef{若一 2:13} 为什么先称子女？\BibleQuote{因为你们的罪因祂的名得蒙赦免，}并且你们重生进入新生命，因此是子女。为什么称父老？\BibleQuote{因为你们认识了那从元始就有的：}因为元始与父职有关。基督在肉身是新的，但在神性上是古老的。我们以为祂有多古老？多大年龄？我们以为祂比祂的母亲更年长吗？确实比祂的母亲更年长，因为\BibleQuote{万物是藉着祂造的。}\BibleRef{若 1:3} 如果万物都是藉着祂造的，那么那古老者就造了那新者将由之出生的母亲。我们以为祂仅在祂的母亲之前吗？是的，而且在祂母亲的祖先之前，祂就早已存在。祂母亲的祖先乃是亚巴郎；主说：\BibleQuote{在亚巴郎以前，我是。}\BibleRef{若 8:58} 我们说在亚巴郎以前？天地在人在世之前就已造成。在这些之前已有主，不如说祂也是。因为祂说得很好：不是说\InnerQuote{在亚巴郎以前，我曾是}，而是\BibleQuote{在亚巴郎以前，我是。}因为一个人说\Quote{曾是}的，那就不再是；说\Quote{将是}的，那还没有是：祂只知道存在。作为天主，祂知道\Quote{存在}：祂不知道\Quote{曾是}和\Quote{将是}。在那里只有一日，却是永远无穷的一日。那一日没有昨天和明天夹在中间：因为当\InnerQuote{昨天}结束时，\InnerQuote{今天}开始，即将被来临的\InnerQuote{明天}所终结。那某一日是一日没有黑暗、没有黑夜、没有空间、没有度量、没有时数的。随你称它什么：你若愿意，称它为一日；你若愿意，称它为一年；你若愿意，称它为许多年。因为论到这同一件事，经上说：\BibleQuote{你的年岁不会穷尽。}但什么时候它被称为一日？那就是当主被说道：\BibleQuote{今日我生了你。}由永恒圣父所生，由永恒而生，在永恒中而生：没有开始，没有界限，没有广延的空间；因为祂是那自有者，因为祂自己就是\Quote{那自有的。}祂把这名字告诉梅瑟：\BibleQuote{你要对他们说：那自有的打发我到你们这里来。}\BibleRef{出 3:14} 那么为什么还说\Quote{在亚巴郎以前}？为什么在诺厄以前？为什么在亚当以前？请听圣经：\BibleQuote{在晨星以前，我已生了你。}总之，在天地以前。为什么？因为\BibleQuote{万物是藉着祂造的，没有祂，什么也没有造。}\BibleRef{若 1:3} 由此你们认识\BibleQuote{父老们：}因为他们因承认\BibleQuote{那从元始就有的}而成为父老。\par

6. \BibleQuote{我写信给你们，青年们。}有子女，有父老，有青年：子女，因为已重生；父老，因为他们承认元始；为什么是青年？\BibleQuote{因为你们已战胜了那恶者。}在子女身上是新生，在父老身上是古老，在青年身上是力量。如果那恶者被青年\Quote{战胜}，那么牠就与我们争战。争战，但不能得胜。为什么？因为我们强，还是因为祂在我们内强？祂在迫害者手中曾显得软弱。祂使我们坚强，祂未曾抵抗迫害祂的人。\BibleQuote{因为祂因软弱而被钉在十字架上，却因天主的德能而活着。}\BibleRef{格后 13:4}\par

7. \BibleQuote{我写信给你们，孩童们。}\BibleRef{若一 2:13} 为什么是孩童？\BibleQuote{因为你们认识了父。我写信给你们，父老们：}他加强这说法，并重复说：\BibleQuote{因为你们认识了那从元始就有的。}要记住你们是父老：如果你们忘记了\BibleQuote{那从元始就有的，}你们就失去了父职。\BibleQuote{我写信给你们，青年们。}要再三思想你们是青年：作战，使你们能得胜；得胜，使你们能得冠冕；谦卑，使你们不致在战斗中跌倒。\BibleQuote{我写信给你们，青年们，因为你们刚强，天主的圣言存留在你们内，你们已战胜了那恶者。}\par

8. 我的弟兄们，这一切——\BibleQuote{因为我们认识了那从元始就有的，因为我们刚强，因为我们认识了父，}——这一切，虽在某种程度上称赞知识，难道不是也称赞爱德吗？如果我们认识了，就让我们相爱：因为没有爱德的知识不能得救。\BibleQuote{知识使人自高自大，爱德却造就人。}\BibleRef{格前 8:1} 如果你有意宣认而不爱，你便开始像魔鬼一样。魔鬼宣认了天主子，说：\BibleQuote{我们与祢有什么相干？}\BibleRef{玛 8:29} 并且被斥退。宣认并拥抱。因为那些魔鬼因自己的不义而惧怕；你们要爱那赦免你们不义的那一位。但是我们怎能爱天主，如果我们爱世界？祂因此准备我们，让爱德居住在我们内。有两种爱：世俗之爱和天主之爱。如果世俗之爱居住在内，就没有地方让天主之爱进入；让世俗之爱让路，天主之爱便居住；让更好的占有位置。你们曾爱世界，不要再爱世界：当你们倒空了心中世俗的爱，你们将饮入神圣的爱；从此爱德开始居住在你们内，从爱德不能生出任何邪恶。因此请听他的话，他如何像一个清理场地的人那样工作。他来到人心，如同来到他要占领的一块田地：但他发现那地是什么状态？若发现丛林，他就拔除；若发现田地已清理，他就栽种。他要在那里栽种一棵树，就是爱德。他要拔除的丛林是什么？就是世俗之爱。请听他，这拔除丛林的人！\BibleQuote{不要爱世界，也不爱世界上的事物；若有人爱世界，父的爱就不在他内。}\BibleRef{若一 2:15}\par

9. 你们已听到：\Quote{谁若爱世界，圣父的爱就不在他内。}任何人不可以在他心里说这是假的，弟兄们：是天主说的；圣神曾藉宗徒发言；没有比这更真实的：\Quote{谁若爱世界，圣父的爱就不在他内。}你愿意拥有圣父的爱，好与圣子同为继承人吗？不要爱世界。摒弃邪恶的世俗之爱，好使你充满天主的爱。你是一个器皿；但你还满着。倒出你所拥有的，好领受你所没有的。的确，我们的弟兄们现在从水和圣神重生；我们也在多年前从水和圣神重生。我们不爱世界是好的，免得圣事在我们内成为定罪的标记，而非加强救恩的方法。那加强人走向救恩的，是拥有爱德的根，拥有\Quote{虔诚的能力}，而不只是\Quote{外表}。\BibleRef{弟后 3:5}外表是好的，外表是神圣的；但若没有根，外表有何用？被砍下的树枝，岂不扔进火里？要有外表，但要在根上。但你怎样被生根，才不会被拔除？藉着持守爱德，正如宗徒保禄所说：\Quote{在爱德上生根奠基。}\BibleRef{弗 3:17}爱德怎能在那世俗之爱蔓生的荒野中生根？彻底清除杂草。你即将播下一颗强大的种子：不要让田里有东西窒息种子。这些就是他所说的拔除性的话：\Quote{不要爱世界，也不要爱世界上的事。谁若爱世界，圣父的爱就不在他内。}\BibleRef{若一 2:15}\par

10. \Quote{因为凡在世界上的，就是肉身的私欲、眼目的私欲、生活的骄傲，}\BibleRef{若一 2:16}这三样是他所说的，不是出于圣父，而是出于世界。世界正在逝去，它的私欲也如此；但那承行天主旨意的人，却永远存留，就如祂永远存留一样。我为什么不该爱天主所造的呢？你想要什么？你愿意爱时间的事物，并随时间逝去；或是不爱世界，而与天主永远生活？暂时事物的河流将人匆匆冲走；但我们的主耶稣基督就像一棵生长在河边的树。祂取了肉身，死了，复活了，升了天。祂愿意将自己，可以说，种植在时间事物的河边。你是否正顺流而下冲入那直坠的深渊？紧紧抓住那棵树。世俗之爱是否正把你卷走？紧紧抓住基督。因为为了你，祂成了暂时的，好使你成为永恒的；因为祂虽成了暂时的，却仍然是永恒的。有时间的事物加给了祂，祂的永恒并未失去什么。但你生来是暂时的，且因罪成了暂时的：你因罪成为暂时的，祂却因怜悯赦罪而成为暂时的。当两人在监狱里时，罪犯和探望者之间有多大差别！因为曾有一次，一个人来找他的朋友，进去探望他，两人似乎都在监狱里；但他们有着极大的区别。一个被他的案件压着；另一个是出于人情到了那里。同样，在我们这有死的状态中，我们被我们的罪过紧紧束缚，祂出于怜悯降下：祂进入被俘者之中，是救赎者而非压迫者。主为我们倾流了祂的血，救赎了我们，改变了我们的望德。目前我们还带着肉身的可朽性，凭着信心接受未来的不朽：在海上我们被波浪颠簸，但我们已有了锚定在陆地上的望德之锚。\par

11. 但我们不要\Quote{爱世界，也不要爱世界上的事。因为世界上的是肉身的贪欲、眼目的贪欲和生活的骄矜。}这三样就是；免得有人或许说：\Quote{世界上的事，天主造了：即天地、海洋、太阳、月亮、星辰，以及天上的一切点缀。海洋的点缀是什么？一切爬行的生物。地上的呢？动物、树木、飞禽。这些是\InnerQuote{在世界上的}，是天主造了他们。那么，为什么我不该爱天主所造的呢？}愿天主的圣神在你内，使你看见这一切都是好的；但如果你爱受造之物而离弃造物之主，就有祸了！它们对你是美丽的，但形成它们的那一位要美丽多少！亲爱的，你们要注意。因为借着比喻你们可以受教：免得撒但悄悄潜入，说它常说的话：\Quote{享受天主的受造物吧；祂为什么造了这些东西，不就是为让你享受吗？}人们喝得酩酊大醉，就灭亡了，忘记自己的造物主；他们不节制地、而是贪欲地使用受造之物，造物主却被轻视了。宗徒论到这样的人说：\Quote{他们崇拜并事奉受造之物，而不事奉造物主——祂是永受赞颂的。}\BibleRef{罗 1:25} 天主并不禁止你爱这些事物，只是不要为着福乐而把情感寄托在它们上面，而是要为此认可并赞美它们，好使你爱你的造物主。我的弟兄们，这就好像一个新郎为他的新娘做了一枚戒指，新娘得到戒指后，爱那戒指胜过爱为她做戒指的新郎：那么，即使是新郎的礼物，她的灵魂岂不也要被判定犯奸淫吗？虽然她只是爱新郎所给的东西？当然，她可以爱新郎所给的；但若她说：\Quote{这枚戒指对我来说就够了，我现在不想见他的面。}那是什么样的女人？谁不讨厌这样的愚蠢？谁不判定她有奸邪的心？你爱金子胜过爱那个人，爱戒指胜过爱新郎：如果你这样，爱戒指胜过爱你的新郎，并且不想见你的新郎；那么，他给你的信物，非但不是使你与他结合，反而使你的心背离了他！因为新郎给信物，正是为了在这信物中，他自己能被爱。同样，天主给了你这一切：你要爱造了它们的那一位。祂还愿意给你更多，就是祂自己——造了这一切的那一位。但如果你爱这些——尽管是天主造的——却忽略造物主，爱世界；那么，你的爱岂不算是奸淫吗？\par

12. 因为\Quote{世界}这个名称不仅指天主所造的这座建筑物——天地、海洋、有形无形之物——而且世上的居民也被称为世界，正如我们称一座\Quote{房屋}，既指墙壁，也指住在里面的人。有时我们称赞房屋，却责怪住在里面的人。因为我们说：\Quote{好房屋}——因为它有大理石和美丽的藻井；在另一种意义上，我们也说：\Quote{好房屋}——那里没有人受冤屈，没有掠夺行为，没有压迫行为。这时我们称赞的不是建筑，而是住在建筑物内的人；然而我们称两者都为\Quote{房屋}。同样，所有爱世界的人，因为借着爱而住在世界中，就如那些虽在肉身中行走在地上、心却在高天的人住在天上；我说，所有爱世界的人被称为世界。他们只有这三样：\Quote{肉身的贪欲、眼目的贪欲、生活的虚幻自负。}因为他们贪图吃、喝、同房，享受这些快感。当然，这并不是说在这些事上没有允许的限度，或者说当说\Quote{不要爱这些事物}时，意思是你不必吃、不必喝、不必生儿育女？所说的不是这个。只是，要有限度，为了造物主的缘故，免得这些事物因你的爱而束缚你：免得你把本应使用的东西，当作享受的爱。但除非有两件事摆在你面前——这个或那个——你不会受到考验：\Quote{你要正义，还是要利益？}\Quote{我没有赖以生活的，没有吃的，没有喝的。}但如果你只能通过不义得到这些？那么，爱那不会失去的，岂不比因不义而失去自己更好？你看见金子的利益，却看不见信仰的损失。因此，他对我们说，这就是\Quote{肉身的贪欲}，即贪求那些属于肉身的事，如食物、肉体的同居，以及诸如此类的其他事物。\par

13. \Quote{眼目的情欲}：他所说的\Quote{眼目的情欲}，指的是所有好奇心。好奇心是何等广泛！它表现在观赏、剧场、魔鬼的圣事、邪术、与黑暗的来往中，无非是好奇心。有时它甚至诱惑天主的仆人，使他们仿佛想要行奇迹，试探天主是否会应允他们行奇迹的祈祷；这就是好奇心：这就是\Quote{眼目的情欲}；它\Quote{不是出于圣父}。如果天主赐给了能力，就行奇迹，因为祂已把行奇迹的路放在你面前：因为不要以为那些没有行过奇迹的人就不能进入天主的国。当宗徒们因魔鬼服从他们而欢喜时，主对他们说了什么？\BibleQuote{不要为这事欢喜，而要因你们的名字已登记在天上而欢喜。}\BibleRef{路 10:20}祂愿意宗徒们为此欢喜，你们也当如此欢喜。如果你的名字没有登记在天上，你实在有祸了！若你没有复活死人，你就有祸吗？若你没有步行海面，你就有祸吗？若你没有驱逐魔鬼，你就有祸吗？如果你领受了行这些事的能力，要谦卑地使用，不要骄傲。因为主甚至论及某些假先知说：\BibleQuote{他们必行神迹奇事。}\BibleRef{玛 24:24}因此，不可有\Quote{世界的野心}：\OriginalTerm{世界的野心}{Ambitio sæculi}，就是骄傲。人希望因自己的尊荣而自大：他以为自己是伟大的，无论是因为财富，或是因为某种权力。\par

14. 这三种欲望都存在，你找不到任何人性的贪欲不受诱惑，除非是肉体的情欲、眼目的情欲，或是生活的骄傲。主也曾被魔鬼用这三种试探。\BibleRef{玛 4:1}当祂禁食后饥饿时，有人对祂说：\BibleQuote{你若是天主子，就叫这些石头变成饼}，这是肉体的情欲的试探。但祂如何击退试探者，并教导祂的士兵如何作战？请注意祂对那人所说的话：\BibleQuote{人生活不只靠饼，而靠天主口中的一切话。}祂也因眼目的情欲在奇迹方面受到了试探，那时魔鬼对祂说：\BibleQuote{你跳下去，因为经上记载：祂为你吩咐了天使，他们要用手托着你，免得你的脚碰在石头上。}祂抵抗了试探者，因为行奇迹只会显得祂要么屈服了，要么出于好奇而行；因为祂作为天主，在祂愿意的时候行奇迹，但却是为了医治软弱者。因为如果祂当时行了奇迹，人们可能会以为祂只想行奇迹。但为了避免人们这样想，请注意祂的回答；当类似的试探临到你时，你也当这样说：\BibleQuote{退到我后面去，撒殚！因为经上记载：你不可试探上主你的天主。}也就是说，如果我这样做，我就是试探天主。祂说了祂要你所说的话。当敌人暗示你说：\Quote{你是什么人，什么样的基督徒？你至今行过一个奇迹吗？或借你的祈祷使死人复活吗？或医治过发热的人吗？如果你真有分量，你就会行些奇迹。}你要回答说：\BibleQuote{经上记载：你不可试探上主你的天主。}因此我不试探天主，好像我行奇迹就属于天主，不行就不属于似的。那么祂的话：\BibleQuote{欢喜吧，因为你们的名字已登记在天上}又成了什么？主如何被\Quote{生活的骄傲}试探？当魔鬼把祂带到高处，对祂说：\BibleQuote{这一切我都要给你，只要你俯伏朝拜我。}他企图用世俗王国的崇高来试探万王之王；但创造天地的天主把魔鬼踏在脚下。魔鬼被主征服，这有什么了不起呢？那么，主在回答魔鬼时，岂不正是教导你要作出怎样的回答吗？\BibleQuote{经上记载：你要朝拜上主你的天主，唯独事奉祂。}坚持这些，你就不会有世界的贪欲；没有世界的贪欲，肉体的情欲、眼目的情欲、生活的骄傲都不能征服你；你将为爱德预备地方，使你能爱天主。因为如果爱世界在那里，天主的爱就不在那里。宁可坚持天主的爱，好使天主永久长存，你也得以永久长存：因为人怎样，他的爱也怎样。爱土地，你就是土地。爱天主，我能说什么呢？你将成为一个神？我不敢自己说，让我们听圣经：\BibleQuote{我曾说过：你们是神，你们都是至高者的儿子。}如果你们愿意成为神和至高者的儿子，就\BibleQuote{不要爱世界，也不要爱世界上的事物。谁若爱世界，圣父的爱就不在他内。因为凡世界上的事物，就是肉体的情欲、眼目的情欲、生活的骄傲，这些不是出于圣父，而是出于世界。}\BibleRef{若一 2:15}\Emph{即}出于人，即爱世界的人。\BibleQuote{世界和它的情欲正在逝去；但那承行天主旨意的，却永远存留，如同天主永远存留一样。}\par

\SourceRecord{Translated by H. Browne.；Nicene and Post-Nicene Fathers, First Series；Vol. 7.；Edited by Philip Schaff.；Buffalo, NY: Christian Literature Publishing Co.,；1888.；<http://www.newadvent.org/fathers/170202.htm>.}

% structure: p0001 source=h1 target=section reason=page-title

\section{若望一书讲道集第三篇}

% structure: p0002 source=h2 target=subsection reason=relative-heading-rank; base=h2

\subsection{若望一书二章十八至二十七节}

\BibleQuote{孩子们，现在是末时了；你们曾听说假基督要来，现在已经出了许多假基督，由此我们就知道现在是末时了。他们从我们中间出去，却不是属于我们的；因为如果他们属于我们，必会留在我们中间；但他们出去，是为了显明他们都不是属于我们的。然而你们从那圣者得了傅油，并且知道一切。我写信给你们，不是因为你们不知道真理，而是因为你们知道真理，并且知道任何谎言都不是出于真理。谁是说谎者呢？不是那否认耶稣是基督的吗？[那否认父和子的，就是假基督。]凡否认子的，也没有父；凡承认子的，也有父。因此，你们要将从起初所听见的，存留在你们内。如果你们从起初所听见的留在你们内，你们也必存留在子和父内。这就是祂所应许我们的，即永生。我写了这些事给你们，论到那些迷惑你们的人；为使你们知道自己有傅油，并且你们从祂所受的傅油常留在你们内，你们就不需要任何人教导你们；因为祂的傅油教导你们一切。}\par

1. \BibleQuote{孩子们，现在是末时了。}在这段学习中，他称呼孩子们为要他们赶快长大，因为\BibleQuote{现在是末时了}。身体的年龄或身材不由人自己决定。人在肉体方面的成长不能由自己决定，正如他出生也不能由自己决定；但若出生取决于意愿，成长也取决于意愿。没有人是\BibleQuote{由水和圣神生的}\BibleRef{若 3:5}除非他愿意。因此，如果他愿意，他就成长或增加；如果他愿意，他就退步。什么是成长？在进步中前进。什么是退步？在欠缺中后退。凡知道自己已重生的人，让他听自己仍是婴孩；让他热切地吸吮母亲的乳房，他就迅速成长。他的母亲就是教会；她的双乳是神圣圣经的两约。因此，让他吸吮一切在时间中为我们的永恒救恩而作为属神真理标记所完成之事的奶，好使被滋养并强壮后，能达到吃固体食物，即\BibleQuote{在起初已有圣言，圣言与天主同在，圣言就是天主}\BibleRef{若 1:1}。我们的奶是谦卑中的基督；我们的肉是同一基督与父同等。祂用奶滋养你，为能用饼喂养你；因为用心神触摸基督，就是知道祂与父同等。\par

2. 因此祂禁止玛利亚触摸祂，并对她说：\BibleQuote{不要触摸我，因为我还没有升到父那里去。}这是什么意思？祂让自己被门徒们触摸，却避开玛利亚的触摸吗？祂不是对那怀疑的门徒说过\BibleQuote{伸过你的指头，摸摸我的伤痕}吗？那时祂升到父那里去了吗？那么为何祂禁止玛利亚，却说\BibleQuote{不要触摸我，因为我还没有升到父那里去？}或者我们要说，祂不怕被男人触摸，却怕被女人触摸？祂的触摸洁净一切血肉。祂愿意先向谁显现，就被他们害怕触摸吗？不是妇女们向男人们宣告了祂的复活，好使蛇以某种反计被击败吗？因为蛇首先通过女人向男人宣布了死亡，所以生命也通过女人向男人宣布了。那么为何祂不愿意被触摸？难道不是因为祂愿意这被理解为神性的触摸吗？神性的触摸出自纯洁的心。那以纯洁的心触摸基督的人，是理解祂与父同等的人。但谁若尚未理解基督的天主性，那人只触摸血肉，不触摸天主性。只触摸到迫害者所触摸的，即钉死祂的人，那算什么呢？但大事是理解圣言天主与天主同在，在起初，万物藉祂而造；正如祂愿意自己被认识，当祂对斐理伯说：\BibleQuote{斐理伯，这么长久我与你们在一起，你还不认识我吗？谁看见了我，就是看见了父}\BibleRef{若 14:9}。\par

3. 但恐怕有人懈怠不前，让他听：\BibleQuote{孩子们，现在是末时了。}前进，奔跑，成长；\BibleQuote{现在是末时了}。这同一个末时很长，但仍是末时。因为他用\Quote{小时}指\Quote{末期}；因为我们的主耶稣基督要在末期来临。但有人会说：怎么是末期？怎么是末时？当然假基督要先来，然后审判的日子才来。若望觉察到这些想法：恐怕人们会变得安然，并因假基督将来就以为不是末时，他对他们说：\BibleQuote{你们曾听说假基督要来，现在已经出了许多假基督。}除了是\Quote{末时}，怎能有许多假基督呢？\par

4. 他称谁为假基督呢？他接着加以解释。\Quote{我们就晓得这是最后的时刻了。} 凭什么？因为\Quote{有许多假基督已经出来了。他们从我们中间出去；} 看看这些假基督！\Quote{他们从我们中间出去：} 因此我们哀悼损失。请听安慰的话。\Quote{他们本不属于我们。} 所有的异端者、所有的分裂分子都从我们中间出去，也就是说，他们离开了教会；然而他们不会出去，如果他们属于我们。因此，在他们出去之前，他们本不属于我们。如果他们在出去之前不属于我们，那么许多人在教会内，并未出去，却仍是假基督。我们敢这样说：为什么？岂不是为了使每个人在教会内时不成为假基督吗？因为他即将描述并指明假基督，我们现在就要看到他们。每个人都应当审问自己的良心，自己是否是假基督。因为假基督在我们的语言中，意思是对抗基督。并不像有些人所认为的那样，假基督之所以这样称呼，是因为他要来在基督之前（\Emph{ante Christum}），即基督在他之后来：不是这个意思，也不是这样写的，而是\Emph{Antichristus}，即对抗基督。现在谁是基督的对抗者，你已经从宗徒本人的解释中认识到，并且明白：只有假基督才能出去；而那些不是对抗基督的人，绝不能出去。因为那不对抗基督的人，紧紧持守在祂的身体内，并被算为其中的一个肢体。肢体之间永远不会彼此对抗。整个身体由所有肢体组成。宗徒关于肢体的共融说了什么？\Quote{若一个肢体受苦，所有的肢体都一同受苦；若一个肢体受光荣，所有的肢体都一同欢乐。} \BibleRef{格前 12:26} 那么，如果一个肢体受光荣，其他肢体就一同欢乐；若它受苦，所有肢体都受苦，肢体的共融就没有假基督。并且有些人，在内心是如此地处于我们主耶稣基督的身体内——因为祂的身体仍在治疗中，完全的康健只将在死人复活时才完成——他们在基督的身体内，如同不良的体液。当这些体液被呕吐出来时，身体就舒缓了：同样，当坏人出去时，教会就舒缓了。有人说，当身体呕吐并吐出它们时：这些体液从我里面出去了，但它们不属于我。怎么不属于我呢？它们不是从我肉中割下来的，而是当它们在我内时压迫了我的胸膛。\par

5. \Quote{他们从我们中间出去；但是，} 不要悲伤，\Quote{他们本不属于我们。} 你如何证明这一点？如果他们属于我们，他们肯定会继续与我们同在。因此你可以看出，许多不属于我们的人，与我们一同领受圣事，与我们一同领受圣洗，与我们一同领受信友们所知道的祝福、圣体以及神圣圣事中的一切：他们与我们一同领受同一祭台的共融，却不属于我们。诱惑证明他们不属于我们。当诱惑像风一样吹向他们时，他们就飞散了；因为他们不是麦子。但所有这些人都会飞散，正如我们必须常常告诉你们的，一旦在审判之日，主的打谷场开始簸扬。\Quote{他们从我们中间出去，但是他们不属于我们；如果他们属于我们，他们必定会继续与我们同在。} 因为，亲爱的，你们想知道这句话多么确定吗？那些偶然出去又回来的人，不是假基督，不是对抗基督的吗？凡不是假基督的，就不可能继续留在外面。而是各人凭自己的意愿，或是假基督，或是在基督内。我们要么在肢体中，要么在不良的体液中。那改善自己的人，在身体内是一个肢体；但那坚持恶行的人，是不良的体液；当他出去时，那些曾被压迫的人就得到舒缓。\Quote{他们从我们中间出去，但是不属于我们；因为如果他们属于我们，必定会继续与我们同在；但（他们出去）是为了显明他们都不是属于我们的。} 他加上了\Quote{为了显明，} 是因为即使他们在内时不属于我们，却不明显；但通过出去，他们就显明了。\Quote{你们从圣者领受了傅油，使你们自己可以显明。} 属灵的傅油就是圣神本身，而圣事的标记在于可见的傅油。关于基督的这傅油，他说，凡拥有它的，都认识善与恶；他们不需要受教导，因为傅油本身教导他们。\par

6. \BibleQuote{我写信给你们，不是因为你们不认识真理，而是因为你们认识它，并且任何谎言都不出于真理。}\BibleRef{若一 2:21} 看，我们受劝告怎样能认识\OriginalTerm{假基督}{Antichrist}。基督是什么？是真理。祂亲自说过：\BibleQuote{我是真理。}\BibleRef{若 14:6} 但\BibleQuote{谎言不出于真理。}因此，凡说谎的人还不属于基督。祂并没有说有些谎言出于真理，有些不出于真理。留意这话句。不要自慰，不要自欺，不要迷惑自己，不要欺骗自己：\BibleQuote{谎言不出于真理。}让我们看看假基督怎样说谎，因为说谎有多种。\BibleQuote{谁是说谎者呢？不是那否认耶稣是基督的吗？}\Quote{耶稣}这个词有一个意义，\Quote{基督}这个词有另一个意义：虽然我们的救主耶稣基督是同一的，但\Quote{耶稣}是祂的本名。正如\Person{梅瑟}、\Person{厄里亚}、\Person{亚巴郎}各有本名一样，我们的主也有祂的本名\Quote{耶稣}；但\Quote{基督}是祂神圣身份的称号。正如我们说先知、司祭，同样，借着\Quote{基督}这个名，我们明白那受傅者，在他身上应有全体百姓的救赎。犹太人民曾盼望这位基督的来临，但因为他谦卑而来，就没有被认出；因为石头很小，他们绊倒在它上面而破碎了。但\BibleQuote{那石头长大，成了一个大山；}\BibleRef{达 2:35} 圣经说什么？\BibleQuote{凡绊倒在这石头上的，必要破碎；这石头落在谁身上，必要把他压得粉碎。}我们必须留意词语的区别：它说，绊倒的将被破碎；但石头落在谁身上，谁将被压得粉碎。起初，因为他谦卑而来，人们绊倒在他身上；因为他将威严地来审判，落在他身上的，他将压得粉碎。但他在将来来临时，不会将那个他在初来时没有击碎的人压得粉碎。那没有在谦卑中绊倒的人，不会畏惧那威严的。简言之，弟兄们，你们已听见：那没有在谦卑中绊倒的人，不会畏惧那威严的。因为对所有恶人，基督是绊脚石；基督所说的任何话，在他们都是苦涩的。\par

7. 因为你们要听并要看。的确，凡从教会出去、与教会的合一隔绝的人，都是假基督；但愿无人怀疑：因为宗徒亲自标记了他们：\BibleQuote{他们从我们中间出去了，但不是属于我们的；因为如果他们属于我们，必会存留在我们中间。}因此，凡不存留在我们中间，反而从我们中间出去的人，显然他们就是假基督。他们怎被证明是假基督呢？借着说谎。\BibleQuote{谁是说谎者呢？不是那否认耶稣是基督的吗？}\BibleRef{若一 2:22} 让我们询问异端者：你在哪里能找到否认耶稣是基督的异端者？看，我亲爱的，一个伟大的奥秘。注意主天主可能启示了我们什么，以及我愿灌输到你们心思中的什么。看，他们从我们中间出去，变成了\OriginalTerm{多纳徒派}{Donatistae}：我们问他们耶稣是否是基督；他们立即承认耶稣是基督。如果那否认耶稣是基督的人是假基督，那么他们既不能称我们为假基督，我们也不能称他们；因此，既不是他们从我们中间出去，也不是我们从他们中间出去。如果我们彼此没有出去，我们就在合一中；如果我们合一，那么在这城中有两个祭台是什么意思？分立的家庭、分立的婚姻是什么意思？有共同的床，却有分裂的基督？他劝告我们，他要我们承认什么是真理：——或是他们从我们中间出去，或我们从他们中间出去。但是不要想象我们是从他们中间出去的。因为我们有主产业的遗嘱，我们诵读它，在那里找到：\BibleQuote{我要将万民赐给你作产业，将地极赐给你作基业。}我们紧握基督的产业；他们不握有它，因为他们不与全地共融，不与那被主的血所救赎的普世身体共融。我们有主自己从死者中复活，祂将自己呈现给怀疑的门徒们触摸；当他们还在怀疑时，祂对他们说：\BibleQuote{基督必须受苦，第三日从死者中复活；并且要因祂的名宣讲悔改和罪过的赦免}\BibleRef{路 24:46}——在哪里？哪条路？对哪些人？——\BibleQuote{从耶路撒冷开始，直到万邦。} 我们的心思在产业的合一中安息了！凡不与此产业共融的，就是出去了。\par

8. 但愿我们不要为此忧伤：\Quote{他们从我们中间出去了，但不是属于我们的；因为如果他们属于我们，就必会留在我们中间。} \BibleRef{若一 2:19} 既然他们从我们中间出去了，他们就是假基督；若是假基督，就是撒谎者；若是撒谎者，就否认耶稣是基督。我们又回到了问题的难点。逐个问他们，他们都会承认耶稣是基督。使我们为难的是，我们过于狭隘地理解了书信中的话。无论如何，你们看到这个问题了；若不理解，这个问题会使我们和他们都陷入困境。要么我们是假基督，要么他们是假基督；他们称我们为假基督，说我们从他们中间出去了；我们同样说他们。但如今这封书信已用此标记指出了假基督：\Quote{凡否认耶稣是基督的，}那人\Quote{就是假基督。} 因此，现在让我们查问谁是否认的；让我们不看口舌，而看行为。因为若问所有人，所有人都异口同声承认耶稣是基督。暂且让舌头静默，询问生活。如果我们发现这一点，如果圣经本身告诉我们否认不仅是口舌之事，也是行为之事，那么我们无疑会发现许多假基督，他们口里承认基督，行为却与基督相悖。我们在哪里看到圣经中这样记载？请听宗徒保禄论及这样的人时说：\Quote{他们口里承认认识天主，行为上却否认祂。} \BibleRef{铎 1:16} 我们也发现这些人是假基督：凡行为上否认基督的，就是假基督。我不听他口说什么，而是看他如何生活。行为在说话，我们还需要言语吗？哪有恶人不想说好话呢？但主对这些人的话是什么？\Quote{你们这些假善人，你们既然邪恶，怎么能说出善事呢？} \BibleRef{玛 12:34} 你们的声音传入我耳中，我却察看你们的内心。我在那里看到邪恶的意愿，你们却展示虚假的果实。我知道该从哪里收集什么；我不\Quote{从蒺藜里摘无花果}，我不\Quote{从荆棘里收葡萄}，因为\Quote{每棵树凭其果实就能被认出。} \BibleRef{玛 12:7} 更虚伪的假基督是那人：他口里承认耶稣是基督，行为上却否认祂。他是撒谎者，因为他口说一套，做另一套。\par

9. 因此，如今弟兄们，若要审问行为，我们不仅发现许多假基督已经出去，而且还有许多尚未显明、根本没有出去的人。因为教会内有多少背信者、欺诈者、行邪术者、求问占卜者、奸淫者、醉酒者、放高利贷者、拐骗孩童者，以及我们无法一一列举的种种恶行；这些行为都与基督的教训相悖，与天主的话语相悖。天主的话语就是基督；凡与天主的话语相悖的，就是属于假基督。因为假基督的意思是\Quote{敌对基督}。你们想知道这些人如何公开反抗基督吗？有时他们做了恶事，有人开始责备他们；由于他们不敢亵渎基督，便亵渎那责备他们的基督的仆人；但如果你们向他们表明你们所说的是基督的话，而非自己的话，他们便竭尽全力要证明你们所说的是自己的话，而非基督的话；然而，如果你们所说的话明明是基督的话，他们甚至攻击基督，开始指责基督：\Quote{怎么，}他们说，\Quote{祂为什么把我们造成这样？} 难道人们不是每天在被指责行为不端时这样说吗？他们被邪恶意念扭曲，反指摘造物主。他们的造物主从天上向他们呼喊（因为那位再造我们的，正是造我们的那一位）：我造了你们什么？我造了人，不是贪婪；我造了人，不是偷盗；我造了人，不是奸淫。你们已听过我的作品赞美我。从三圣童口中发出的赞美诗保护他们免于火焰。上主的作品赞美上主：天、地、海赞美祂；天上的一切赞美祂，天使赞美祂，星辰赞美祂，光明赞美祂，凡游泳的、飞翔的、行走的、爬行的都赞美祂；所有这些都赞美上主。你们可曾在那赞美诗中听到贪婪赞美上主？可曾听到醉酒赞美上主？可曾听到奢靡赞美祂、轻浮赞美祂？凡你们在赞美诗中未听到赞美上主的，那东西就不是上主所造的。纠正你们自己所造的，好使天主在你们内所造的得以得救。但如果你们不愿，反而喜爱并拥抱你们的罪，你们就是敌对基督。无论你们在里面还是在外面，你们都是假基督；无论你们在里面还是在外面，你们都是糠秕。但为什么你们不在外面？因为你们还没有遇到把你们吹走的风。\par

10. 这些事如今已显明，我的弟兄们。不要有人说，我不敬拜基督，但我敬拜祂的父天主。\Quote{凡否认子的，也没有父；凡承认子的，也有父也有子。} 祂对你们这些麦粒说话：让那些曾是糠秕的人聆听，并成为麦粒。让每个人省察自己的良心，如果他是爱世界的，就让他改变；让他成为爱基督的人，免得成为假基督。如果有人对他说他是假基督，他会愤怒，认为这是对他的侮辱；也许，若那与他争论的人说他是假基督，他会威胁要提起诉讼。基督却对他说：你要忍耐；如果你被妄加评论，要与我一同喜乐，因为我也被假基督们妄加评论；但如果你是被如实评论，就要与自己的良心达成共识；如果你害怕被人这样称呼，就更要害怕成为这样的人。\par

11. \BibleQuote{你们应让那从起初所听见的，常存在你们内。如果你们从起初所听见的，常存在你们内，你们也必存在子内，并存在父内。这就是他向我们应许的恩许。}\BibleRef{若一 2:24} 因为或许你们会问关于报酬的事，说：看，我从起初所听见的，我安全地保存在我内，我依此而行；为了坚持的缘故，我承担了危险、劳苦、诱惑：有什么成果？什么报酬？他以后将给我什么？因为在这世界上，我看到我在诱惑中劳苦；我看不到这里有安息：只是必死的生命压着灵魂，可朽坏的身体把它拉向低处；但我承担一切，为的是那\BibleQuote{我从起初所听见的}\BibleRef{智 9:15} 能常存在我内；并使我向我的天主说：\Quote{由于你嘴唇的言语，我保持了艰难的道路。} 那么为了什么报酬？请听，不要气馁。如果你们在劳苦中气馁，要在所应许的报酬上坚强。有谁在葡萄园工作，却让将要得到的报酬从心中溜走呢？假设他忘记了，他的手就软了。对所应许报酬的记忆使他在工作中坚持不懈：然而那应许的人是一个可能欺骗你期望的人。何况在天主的田地里，你们应该怎样更坚强呢？因为那应许者是真理，他既没有继承者，也不会死，也不会欺骗那领受应许的人！那么，这应许是什么？让我们看看他应许了什么。是世人很爱的金子吗？或是银子？或是产业（为了产业人挥霍金子，尽管他们很爱金子）？或是美好的田地、宽敞的房子、许多奴隶、无数牲畜？这些并不是报酬——他劝我们恒心忍受劳苦所期待的报酬。这些报酬有什么名称？\BibleQuote{永生。}\BibleRef{玛 25:34} 看，这就是天主所威胁的：永火。对右边的人，他说：\BibleQuote{来吧，我父所祝福的，承受从创世以来为你们预备的国度。}\BibleRef{玛 25:34} 对左边的人，他说：\BibleQuote{滚到为魔鬼和他的使者预备的永火里去！}\BibleRef{玛 25:41}\par

12. 请记住，我的弟兄们，基督应许了我们永生：\BibleQuote{这就是他向我们应许的恩许——永生。我曾给你们写了这些，论到那些诱惑你们的人。}\BibleRef{若一 2:25} 不要让任何人诱惑你们走向死亡：要渴望永生的应许。世界能应许什么？任它应许什么，它或许应许给一个明天就要死的人。你们将以何面目去见那永远存在的主？\Quote{但一个有权势的人威胁我，所以我必须作恶。} 他威胁什么？监狱、锁链、烈火、酷刑、野兽：是的，但不是永火？要畏惧那全能者所威胁的；要爱那全能者所应许的；那么全世界在我们眼中都变得微不足道，无论它应许还是恐吓。\BibleQuote{我曾给你们写了这些，论到那些诱惑你们的人；为使你们知道你们有傅油，并且我们所领受的傅油常存在你们内。} 在傅油中我们有圣事的标记（无形之事的标记），德能本身是看不见的；看不见的傅油就是圣神；看不见的傅油就是那爱德，无论它在谁内，都将成为他的根：无论太阳多么炽热，它都不能枯萎。凡有根的都因太阳的温暖而得到滋养，而不是枯萎。\par

13. \BibleQuote{你们不需要任何人教导你们，因为祂的傅油教导你们一切事。} \BibleRef{若一 2:27} 那么，我弟兄们，我们教导你们，又有什么用处呢？如果祂的傅油教导你们一切事，看来我们是徒劳无功的。我们这样呼喊，又是什么意思呢？让我们把你们交给祂的傅油，让祂的傅油教导你们吧。但这只是我向自己提出的问题；我也向那位宗徒提出同样的问题：让他俯听一个向他请教的小子吧。我对若望本人说：那些听你讲话的人，他们拥有那傅油吗？你已经说过：\Quote{祂的傅油教导你们一切事。}那你写这封信，又有什么目的呢？你给了他们什么教导？什么训诲？什么建树？请看，弟兄们，请看一个伟大的奥秘。我们言语的声音触动耳朵，但导师在内。不要以为任何人能从人学到什么。我们能用声音劝诫；但如果没有一位内在的导师教导，我们的喧嚣就是徒然。是的，弟兄们，你们有心要明白吗？你们不是都听了这篇讲道吗？然而有多少人会从这地方空手而去！就我而论，我已对所有人讲了；但那些内在的傅油不对他们讲话、内在的圣神不教导他们的人，他们回去时仍是未受教的。来自外部的导师的教导，好像是帮助和劝告。教导人心的那一位，祂的座位在天上。因此祂自己也在福音中说：\BibleQuote{不要称地上的人为师傅，你们的师傅只有一位，就是基督。} \BibleRef{玛 23:8} 因此让祂自己在你们内说话，当没有人在你们旁边时：因为虽然有人在你身边，却没有人在你心中。但不要让你心中空无一人：让基督在你心中，让祂的傅油在你心中，免得你的心像在荒野中干渴，又没有泉源可以灌溉。所以，我说，有一位内在的导师教导：基督教导，祂的默感教导。哪里没有祂的默感和傅油，外部的言语就徒然喧嚷。同样，弟兄们，我们从外部说的话，正如农夫对树：他从外部工作，浇水、殷勤培植；让他从外部做他愿意做的，他能结出果子吗？他能给光秃的树枝披上浓荫的叶子吗？他从内在做了任何这样的事吗？但这是谁的作为呢？听那农夫（即宗徒）：你们既知道我们是什么，也听内在的导师：\BibleQuote{我栽种了，阿颇罗浇灌了，但使之生长的是天主；栽种的算不得什么，浇灌的算不得什么，只算那使之生长的天主。} \BibleRef{格前 3:6} 这样我们对你们说：无论我们栽种，还是我们浇灌，借着说话我们算不得什么；但那使之生长的，即天主：这就是\BibleQuote{祂的傅油教导你们一切事。}\par

\SourceRecord{Translated by H. Browne.；Nicene and Post-Nicene Fathers, First Series；Vol. 7.；Edited by Philip Schaff.；Buffalo, NY: Christian Literature Publishing Co.,；1888.；<http://www.newadvent.org/fathers/170203.htm>.}

% structure: p0001 source=h1 target=section reason=page-title

\section{\BibleBook{若望}{一书}第四篇讲道}

% structure: p0002 source=h2 target=subsection reason=relative-heading-rank; base=h2

\subsection{\BibleBook{若望}{一书} 2:27-3:8}

\BibleQuote{\Emph{这傅油是真实的，绝不虚假。你们要照它所教训你们的，常存在祂内。孩子们，你们要住在祂内，这样，当祂显现时，我们可以坦然无惧，在祂来临时，不致在祂面前蒙羞。你们既然知道祂是正义的，就该知道凡行正义的人，都是由祂而生的。请看，圣父赐给我们何等的爱情，使我们得称为天主的子女，而且我们也真是如此；因此世界不认识我们，是因为不认识祂。亲爱的，现在我们是天主的子女，但将来如何，还未显明；我们知道，祂显现时，我们必要相似祂，因为我们要看见祂实在怎样。凡对祂怀着这希望的，必洁净自己，如同祂是洁净的一样。凡犯罪的，也行不法；罪就是不法。你们知道，祂曾显现，为除去罪；在祂内并无罪。凡住在祂内的，就不犯罪；凡犯罪的，未曾看见祂，也未曾认识祂。孩子们，不要让人迷惑你们；行义的，才是义人，如同祂是义的一样。犯罪的，是属魔鬼，因为魔鬼从起初就犯罪。天主子所以显现，是为消灭魔鬼的作为。}}\par

\Person{弟兄们}，你们记得，昨天的讲论结束于这一点：\BibleQuote{你们不需要任何人教导你们，而是傅油本身教导你们一切。}这点，我相信你们记得，我们曾这样向你们解释：我们这些从外面向你们耳朵说话的人，就像从外面给树培土的工人，但我们不能使树生长，也不能结出果实；只有那创造、救赎并召叫你们的那一位，祂藉着信德和圣神住在你们内，必须在你们内说话，否则我们的一切言语喧哗都是徒然。这从哪里显明呢？从这一点：许多人听时，并非所有的人都信服所说的话，只有那些天主在内里向他们说话的人才信服。现在，祂在内里向他们说话的，就是那些给祂留地方的人；而那些给天主留地方的人，就是\BibleQuote{不给魔鬼留地方}的人。\BibleRef{弗 5:27}因为魔鬼愿意住在人心中，并在那里说那些能迷惑人的事。但主耶稣怎么说？\BibleQuote{这世界的首领已被赶出去。}\BibleRef{若 12:31}从何处被赶出？从天上或地上？从世界的构造中？不，而是从信者的心中。侵入者被赶出后，让救赎者住在内里：因为是同一位创造了他们的，也救赎了他们。魔鬼如今从外面攻击，不能征服那在内里占有的人。他从外面攻击，投入各种诱惑；但那人不同意诱惑，因为天主在他内说话，以及你们所听过的傅油。\par

2. \BibleQuote{这是真实的，}即这个傅油；\Emph{即}主的圣神教导人，不说谎：\BibleQuote{且不是虚假的。你们要照它所教训你们的，常存在祂内。孩子们，你们要住在祂内，这样，当祂显现时，我们可以坦然无惧，在祂来临时，不致在祂面前蒙羞。}\BibleRef{若一 3:27} \Person{弟兄们}，你们看：我们相信耶稣，我们未见过祂；那些见过、摸过、听过祂亲口说话的人宣告了祂；他们为使全人类信服真理，是由祂派去的，并非擅自前往。他们被派往何处？你们读福音时听到：\BibleQuote{你们往普天下去，向一切受造物宣传福音。}因此，门徒被派往\Quote{各处，}以奇迹和征兆证明他们所说的是他们见过的。我们相信我们未见过的那位，并期待祂来临。凡藉信德期待祂的，当祂来临时必要喜乐；那些没有信德的，当那现今看不见的来临时，必要蒙羞。那羞愧不会只有一天就过去，如同那些在某种过失中被发现、被同伴嘲笑的人通常的羞愧。那羞愧将把蒙羞者带到左边，以致对他们说：\BibleQuote{进入那为魔鬼和他的使者所预备的永火里去！}\BibleRef{玛 25:31}那么，让我们住在祂的话语中，这样当祂来临时我们就不致羞愧。因为祂自己曾在福音中对那些信祂的人说：\BibleQuote{你们如果住在我的话内，就真是我的门徒。}\BibleRef{若 8:31}他们好像问：有什么结果？祂说：\BibleQuote{你们会认识真理，而真理必会使你们获得自由。}因为目前我们的救恩在于望德，不在于事实：我们尚未拥有所应许的，但盼望它的来临。而且\BibleQuote{那应许的是信实的，}\BibleRef{希 10:23}祂不欺骗你们：只要你们不气馁，而是等候那应许。因为祂，就是真理，不能欺骗。你们不要做说谎的人，说一套做一套；保持信德，祂就保持祂的应许。但如果你们不持守信德，欺骗你们的不是那应许者，而是你们自己。\par

3. \BibleQuote{如果你们知道祂是义德的，也该知道凡行义德的都由祂而生。} \BibleRef{若一 2:29} 如今我们所有的义德，是由信德而来。完美的义德，只在天使中才有；而且天使若与天主相比，也几乎算不得完美。如果天主所造的灵魂和神灵有任何完美的义德，那是在圣洁、公义、良善的天使身上，他们从未因过失而偏离，从未因骄傲而堕落，而是恒常留存在对天主圣言的默观中，除造物主外别无甘饴。在他们内有着完美的义德。但在我们内，义德已开始藉信德、藉圣神而存在。你们曾听见圣咏被诵读时这样说：\BibleQuote{你们当以告解向主开始。} \BibleQuote{开始，} 它说；我们义德的开始便是承认罪过。你们已开始不护庇你们的罪，现在你们已开始了义德。但义德要在你们内成为完美的，是当你们不再以别的为乐，当 \BibleQuote{死亡被胜利吞灭} \BibleRef{格前 15:24}，当没有情欲的瘙痒，当没有与血肉的争战，当有胜利的棕榈枝、战胜仇敌的凯旋，那时才会有完美的义德。目前我们仍在战斗。如果我们在战斗，我们就在竞技场上；我们击打，也挨打。但谁会得胜，尚待分晓。那得胜的人，在击打时也不仗恃自己的力量，而依靠那鼓励他的天主。魔鬼与我们作战时是孤单的。如果我们与天主同在，就能战胜魔鬼；因为如果你独自与魔鬼作战，你必被战胜。他是个老练的敌人，他赢过多少棕榈枝！想想他把我们扔到了何等境地！我们生而为有死的，正是因为他首先将我们的原祖从乐园中扔了下来。既然他如此老练，该做什么呢？让全能者被呼求来帮助你对抗魔鬼的阴谋。让那不能被战胜的祂住在你内，你就能稳妥地战胜那惯于战胜人的他。但要战胜谁呢？那些天主不在其内的人。弟兄们，为使你们知道，亚当在乐园中轻看了天主的诫命，挺起脖子，仿佛想要成为自己的主人，不愿服从天主的旨意；于是他从那不死和真福中堕落了。但有一个人，一个虽然生为有死却极有经验的人，他坐在粪堆上，身上生满虫子，却战胜了魔鬼。是的，亚当本人那时也战胜了，就在约伯身上；因为约伯是他的后裔。所以，在乐园中失败的亚当，在粪堆上得胜了。在乐园中，他听从了魔鬼放进女人口中的劝说；但在粪堆上，他对厄娃说：\BibleQuote{你说话如同一个愚妇。} \BibleRef{约 2:10} 他在那里倾听，在这里回答；他在快乐时听了，在受鞭打时得胜了。所以，弟兄们，注目书信中接下来的话，因为这是它要我们牢记的，好让我们确实战胜魔鬼，但不是凭我们自己。\BibleQuote{如果你们知道祂是义德的，} 它说，\BibleQuote{也该知道凡行义德的都由祂而生：} 由天主，由基督而生。并且他说\Quote{由祂而生}时，是在鼓励我们。因此，既然我们由祂而生，我们已经是成全的了。\par

4. 请听。\BibleQuote{请看父赐给我们何等的爱，使我们得称为天主的子女，而且我们真是如此。} \BibleRef{若一 3:1} 因为那些被称为子女、却不是子女的，有那名而无其实，于他们何益？有多少人被称为医生，却不知如何医治！有多少人被称为守夜者，却整夜沉睡！同样，许多人被称为基督徒，但在行为上却名不副实；因为他们没有活出所蒙的称号，即生活、品行、信德、望德、爱德上都不是。弟兄们，你们在此听到了什么？\BibleQuote{请看父赐给我们何等的爱，使我们得称为天主的子女，而且我们真是如此；因此世界不认识我们，因为不曾认识祂。世界也不认识我们。} 有一个普世的基督徒世界，也有一个普世的不敬虔世界；因为普世都有不敬虔的人，普世也有敬虔的人；那些不认识这些。我们想，他们在什么意义上不认识他们呢？他们嘲笑那些生活良善的人。你们要注意并察看，因为你们中间可能也有这样的人。你们中每一位现在虔诚生活、轻视世俗之事、不愿去观看表演、不愿在所谓庄严风俗下大醉，甚至更糟，在圣日的幌子下使自己不洁的人——那不愿做这些事的人，是如何被那些做这些事的人嘲笑的！如果他被认识，他还会被讥讽吗？但为什么他不被认识？\BibleQuote{世界不认识祂。} 谁是\Quote{世界}？那些居住在世界中的。正如我们说\Quote{房子}，指的是住客。这些事已对你们反复说了，我们不愿再重复使你们厌烦。到现在为止，当你们听到\Quote{世界}一词，从坏的意义上，你们知道必须只将其理解为世界的爱慕者，因为他们因爱慕而居住，因居住而配得此名。因此世界不认识我们，因为它不认识祂。祂自己曾行走在这里，主耶稣基督在肉身内；祂是天主，却隐藏在软弱中。为什么祂不被认识？因为祂斥责世人所有的罪。他们因喜爱罪的快乐，而不承认天主；因喜爱那热病所激起的，而对医生行了不义。\par

5. 对我们而言，我们是什么呢？我们已由祂所生；但因我们只是处于希望之中，祂说：\Quote{亲爱的，现在我们是天主的儿女了。}现在已经是？那么我们还期待什么呢，如果已经是天主的儿女？\Quote{但尚未}，祂说，\Quote{显明我们将成为什么。}然而，除了天主的儿女，我们还能成为什么？且听下文：\Quote{我们知道，当祂显现时，我们将相似祂，因为我们必看见祂的真实面目。}请理解，我亲爱的。这是一件大事：\Quote{我们知道，当祂显现时，我们将相似祂，因为我们必看见祂的真实面目。}首先请注意，什么是所谓的\Quote{是}。你们知道什么是如此称呼的。那被称为\Quote{是}的，不仅被称为，而且实是，是不变的：它永远存留，不可改变，绝无朽坏之分：它没有长进，因它是完善的；也没有欠缺，因它是永恒的。而这是什么呢？\Quote{在起初已有圣言，圣言与天主同在，圣言就是天主。}\BibleRef{若 1:1} 而这又是什么？\Quote{祂虽具有天主的形体，却没有以自己与天主同等为应被持守不舍的。}\BibleRef{斐 2:6} 以这种方式看见基督——基督具有天主的形体，天主的圣言、圣父的独生子、与圣父同等——对恶人是不可能的。但就圣言成了血肉而言，恶人也将有能力看见祂：因为在审判之日，恶人也要看见祂；因为祂将如祂曾来受审判那样来审判。具有同一形体，是人，然而也是天主：因为\Quote{凡信赖世人的人，是可咒骂的。}\BibleRef{耶 17:5} 祂作为人来受审判，祂将作为人来审判。而如果祂不被看见，那么所写的这话又是什么？\Quote{他们要瞻望他们所刺透的。}\BibleRef{若 19:37} 因为关于不敬虔的人，经上说他们将看见并羞愧。不敬虔的人怎能不看见呢？当祂将一些人置于右边，另一些人置于左边时？对右边的人，祂将说：\Quote{我父所祝福的，你们来承受国度。}\BibleRef{玛 25:41} 对左边的人，祂将说：\Quote{进入永火。}他们将只看见仆人的形体，而看不见天主的形体。为什么？因为他们是不敬虔的；而且主亲自说：\Quote{心里洁净的人是有福的，因为他们要看见天主。}\BibleRef{玛 5:8} 因此，我的弟兄们，我们要看见一种景象，\Quote{眼睛未曾看见，耳朵未曾听见，人心也未曾想到的}\BibleRef{格前 2:9}：一种景象，超越一切地上美丽——金子、银子、树林与田野；海洋与天空的美丽，太阳与月亮的美丽，星辰的美丽，天使的美丽：超越一切，因为一切美丽都由它而来。\par

6. 那么，当我们看见这景象时，我们将\Quote{成为}什么？向我们应许了什么？\Quote{我们将相似祂，因为我们必看见祂的真实面目。}舌头已尽其所能，发出了话语；其余的事由心灵去思想。因为即使若望本人，对比那\Quote{是}的，又能说什么呢？或我们这些远不如他功德的人，又能说什么呢？因此，让我们回归祂的傅油，回归那内在教导我们无法言说之事的傅油；因为你们现在不能看见，所以你们的职责和本分在于渴望。一个虔诚基督徒的整个生命就是神圣的渴望。现在你们所渴望的，你们尚未看见；然而藉着渴望，你们变得有容量，以便当那你们可以看见的来临时，你们将被充满。因为正如，如果你们要装满一个袋子，并知道将要给予的东西有多大，你们撑开袋口或皮囊，或其他什么；你们知道要放多少，并看到袋子窄小；藉着撑开，你们使它更能容纳。同样，天主藉着拖延我们的希望，伸展我们的渴望；藉着渴望，伸展心灵；藉着伸展，使它更有容量。因此，我的弟兄们，让我们渴望，因为我们将被充满。请看保禄仿佛在拓宽他的胸怀，以便能够接纳那将要来的。他说：\Quote{这不是说我已经达到或已经成为全备的；弟兄们，我不以为自己已经获得了。}\BibleRef{斐 3:13} 那么，如果你尚未获得，你在今生做什么呢？\Quote{但有一件事：忘记背后，努力面前的，向着目标竭力追求，为得天主在基督耶稣内从高天召我来得的奖赏。}他说他努力向前或伸展自己，并说他\Quote{竭力}追求。他感到自己太小，无法容纳\Quote{眼睛未曾看见，耳朵未曾听见，人心也未曾想到的}\BibleRef{格前 2:9}。这就是我们的生命，藉着渴望来操练我们。但神圣的渴望操练我们，恰好与我们修剪对世界之爱的渴望成比例。我们已经说过：\Quote{倒空那将要被充满的。}你将被充满以善：倒出恶。假使天主想要以蜂蜜充满你：如果你满了醋，蜂蜜将放在哪里？那器皿所盛的东西必须倒出：器皿本身必须洁净；必须洁净，尽管辛苦，尽管用力擦洗，才能适合那东西，无论它是什么。让我们说蜂蜜，说金子，说酒。无论我们说什么，那无法言说的，我们愿意称它为——天主。当我们说\Quote{天主}时，我们说了什么？那一个音节就是我们期待的全部吗？因此，我们所能说的，都低于祂：让我们向祂伸展自己，以便当祂来临时，祂将充满我们。因为\Quote{我们将相似祂，因为我们必看见祂的真实面目。}\par

7. \Quote{凡对祂怀着这望德的。}你看祂将我们的位置摆在了\Quote{望德}中。你看宗徒保禄与他的同宗徒何等一致：\Quote{我们因望德而得救。但看得见的望德，就不是望德；因为人看见了的，为什么还盼望？如果我们盼望那看不见的，我们就要耐心等待。}\BibleRef{罗 8:24}这份忍耐本身便锻炼了渴望。继续前行，因为祂继续存在；坚持行走，好使你到达终点；因为你所趋向的不会移开。看：\Quote{凡对祂怀着这望德的，必洁净自己，如同祂是洁净的一样。}你看，祂说\Quote{洁净自己}，并没有取消自由意志。是谁洁净我们？岂不是天主？是的，但若你不愿意，天主也不会洁净你。因此，当你将自己的意志结合于天主，你就在洁净自己。你洁净你自己，不是凭你自己，而是凭那位来居住在你内的祂。然而，因为你在此事上用意志做了一些事，所以有某事归功于你。但归功于你，只是为了叫你如同在圣咏中那样说：\Quote{求你作我的助佑，不要离弃我。}若你说\Quote{作我的助佑}，你便做了一些事；因为你若什么都不做，祂怎能被称为\Quote{助佑}你呢？\par

8. \Quote{凡犯罪的，也是行不义。}谁也不要说：罪是一回事，不义是另一回事；谁也不要说：我是有罪的人，却不是行不义的人。因为\Quote{凡犯罪的，也是行不义；罪就是不义。}那么，关于罪和不义，我们该做什么？听祂说什么：\Quote{你们知道，祂曾显现，为除去罪过；在祂内并没有罪过。}\BibleRef{若一 3:5}那没有罪过的祂，来为除去罪过。因为若祂内有罪，罪定当从祂除去，而不是祂亲自除去。\Quote{凡住在祂内的，就不犯罪。}\BibleRef{若一 3:6}他留在祂内有多久，就不犯罪有多久。\Quote{凡犯罪的，没有看见过祂，也没有认识过祂。}这是一个大问题：\Quote{凡犯罪的，没有看见过祂，也没有认识过祂。}这不奇怪。我们没有看见过祂，但将要看见；没有认识过祂，但将要认识：我们相信一位还不认识的祂。或者，也许我们因信德已认识，尚未因实际目睹而认识？但若信德还看不见，为何我们被称为受了光照？有一种光照是因信德，一种是因目睹。目前，当我们还在旅途中，\Quote{我们因信德往来，并非因目睹}\BibleRef{格后 5:7}，亦即实际目睹。因此，我们的正义也是\Quote{因信德，而非因目睹}。我们的正义将达于圆满，当我们藉实际目睹看见时。只是在这期间，我们不要放弃那属于信德的正义，因为\Quote{义人因信德而生活}\BibleRef{罗 1:17}，正如宗徒所说。\Quote{凡住在祂内的，就不犯罪。}因为\Quote{凡犯罪的，没有看见过祂，也没有认识过祂。}那犯罪的人不信；但若一个人相信，就他的信德而言，他不犯罪。\par

9. \Quote{孩子们，不要让任何人欺骗你们。那行正义的，就是正义的，如同祂是正义的一样。}\BibleRef{若一 3:7}什么？听到我们\Quote{正义如同祂是正义的一样}，我们就以为自己和天主平等吗？你必须知道这里的\Quote{如同}是什么意思：刚才他这样说：\Quote{洁净自己，如同祂是纯洁的一样。}那么我们的纯洁与天主的纯洁相似且相等吗？我们的正义与天主的正义相等吗？谁能这样说？但\Quote{如同}一词并不总表示相同。例如，如果一个人看了这座大教堂，想建造一座较小的教堂，但保持相同的比例：比如，这座教堂宽一尺，长两尺，他也建造宽一尺、长两尺的教堂：那么人们看见他建造的\Quote{如同}这座一样。但这座教堂长一百肘，那座长三十肘：它仍是\Quote{如同}这座，却不相等。可见这个\Quote{如同}不总是指对等和平等。再比如，人的脸与其镜中影像之间有何等差异：影像上有脸，身体上也有脸；影像出于模仿，身体出于实在。我们怎么说？我们说，\Quote{如同}这里有眼睛，那里也有；\Quote{如同}这里有耳朵，那里也有耳朵。事物不同，但\Quote{如同}说的是相似。那么，我们内也有天主的肖像；但不是那与父平等的圣子所有的；然而，除非我们也按我们的尺度\Quote{如同}祂，我们毫不可能被称为相似祂。\Quote{祂洁净我们，如同祂是纯洁的一样}；但祂从永远就是纯洁的，我们因信德而纯洁。我们\Quote{正义如同祂是正义的一样}；但祂是在祂不变永恒中如此，我们因相信一位我们看不见的祂而成为正义，好使我们有一天看见祂。即使当我们的正义达于圆满，当我们与天使平等时，也还不能与祂平等。那么现在距离祂何止遥远，即使那时也不相等！\par

10. \BibleQuote{那犯罪的，是出于魔鬼，因为魔鬼从起初就犯罪。}\BibleRef{若一 3:8} \Quote{是出于魔鬼：}你们知道他是什么意思：是模仿魔鬼。因为魔鬼没有创造任何人，没有生任何人，也没有创造任何人；但凡模仿魔鬼的人，就好像是由他而生，成为魔鬼的子女；是借着模仿，而不是由他亲身所生。在什么意义上，你是亚巴郎的子女，却又不是亚巴郎所生？同样，犹太人本是亚巴郎的子女，却没有模仿亚巴郎的信德，就变成了魔鬼的子女：他们由亚巴郎的肉身而生，却没有模仿亚巴郎的信德。那么，如果那些从亚巴郎所生的人因为不模仿而失去了产业，那么你，虽不是从他而生，却因模仿而成为子女，这样你就会借着模仿成为他的子女。如果你模仿魔鬼，如同他变得骄傲且对天主不敬，那么你就是魔鬼的子女：是通过模仿，而不是因为他创造了你或生了你。\par

11. \Quote{天主子显现出来，正是为此。}现在，弟兄们，请注意！所有的罪人，作为罪人，都是出于魔鬼。亚当是由天主造的；但当他同意魔鬼时，他就由魔鬼所生；并且他生了所有像他一样的人。我们生来就带着情欲本身；甚至在我们添加个人罪行之前，我们就从那个定罪中出生。因为如果我们生来毫无罪过，为什么婴孩要跑去领受圣洗以求赦免？因此，请注意，弟兄们，两个出生源头：亚当和基督；两个人：但其中一个，是人而为人；另一个，是人而为天主。借着那人而为人，我们成了罪人；借着那人而为天主，我们得以成义。那个出生把我们抛向死亡；这个出生把我们提升到生命：那个出生带来罪恶；这个出生解除罪恶。因为基督正是为此而来，作为人，为消除人的罪恶。\Quote{天主子显现出来，正是为了毁灭魔鬼的作为。}\par

12. 其余的事，我所爱的，我交托给你们的思想，以免加重你们的负担。因为我们努力要解决的问题正是这一点——我们称自己为罪人：因为若有人说自己无罪，便是说谎。而在同一若望的书信中，我们读到：\BibleQuote{如果我们说我们没有罪，便是欺骗自己。}\BibleRef{若一 1:8} 因为你们应当记得前面的话：\BibleQuote{如果我们说我们没有罪，便是欺骗自己，真理也不在我们内。}然而，另一方面，紧接着你们又被告知：\BibleQuote{凡由天主生的，就不犯罪；那犯罪的，未曾看见祂，也不认识祂。——凡犯罪的，都是出于魔鬼。}这又使我们惊惧。我们在什么意义上是由天主生的，又在什么意义上承认自己是罪人？难道我们要说，因为我们不是由天主生的吗？那么这些圣事对于婴孩又意味着什么？若望说了什么？\BibleQuote{凡由天主生的，就不犯罪。}然而同一若望又说：\BibleQuote{如果我们说我们没有罪，便是欺骗自己，真理也不在我们内！}这是一个大问题，且令人困惑；但愿我已经使你们渴望得到解答，我所爱的。明天，因主的名，我们将根据祂所赐的，讲论这事。\par

\SourceRecord{Translated by H. Browne.；Nicene and Post-Nicene Fathers, First Series；Vol. 7.；Edited by Philip Schaff.；Buffalo, NY: Christian Literature Publishing Co.,；1888.；<http://www.newadvent.org/fathers/170204.htm>.}

% structure: p0001 source=h1 target=section reason=page-title

\section{\WorkTitle{\BibleBook{若望}{一书}第五篇讲道}}

% structure: p0002 source=h2 target=subsection reason=relative-heading-rank; base=h2

\subsection{\BibleRef{若一 3:9-18}}

\BibleQuote{凡由天主生的，就不犯罪过，因为天主的种子存留在他内；他不能犯罪，因为他由天主生。天主的子女和魔鬼的子女，在这事上可以认出：凡不行正义的，和不爱自己弟兄的，就不是出于天主。原来你们从起初所听的训令就是：我们应彼此相爱；不可像加音，他原属于邪恶者，杀了自己的弟弟。为什么杀了他？因为他的行为是邪恶的，而他弟弟的行为是正义的。弟兄们，如果世界恼恨你们，不必惊奇。我们知道，我们已出死入生了，因为我们爱弟兄们；那不爱的，就存留在死亡内。凡恼恨自己弟兄的，便是杀人的；你们也知道，凡杀人的，没有永远的生命存留在他内。我们所以认识了爱，因为那一位为我们舍掉了自己的生命；我们也应当为弟兄们舍掉生命。谁若有今世的财物，看见自己的弟兄有急难，却对他封闭自己怜悯的心肠，天主的爱怎能存在他内？孩子们，我们爱，不可只用言语，也不可用口舌，而要用行动和事实。}\par

1. 请专心听，我恳求你们，因为我们要处理的并非小事；我不怀疑，因为你们昨天已专注于此事，今天怀着更大的专心聚集。因为这并非小问题：他在这封书信中说：\BibleQuote{凡由天主生的，不犯罪，}\BibleRef{若一 3:9}而他在同一封书信上面又说：\BibleQuote{如果我们说我们没有罪，就是欺骗自己，真理就不在我们内。}\BibleRef{若一 1:8}一个人被同一封书信中的两句话所困扰，他该怎么办？如果他承认自己是罪人，他害怕被说：那么你就不是由天主生的，因为经上记载：\BibleQuote{凡由天主生的，不犯罪。}但如果他说自己是正义的，没有罪，他又会从同一封书信的另一面受到打击：\BibleQuote{如果我们说我们没有罪，就是欺骗自己，真理就不在我们内。}那么，处在两者之间，他该说什么、承认什么、宣告什么，他找不到。要宣告自己无罪，是充满危险的；并且不仅危险，也充满错误：他说：\Quote{我们欺骗自己，}他说，\Quote{如果我们说我们没有罪，真理就不在我们内。}但愿你无罪，并且这样说！那么你就会说真话，而且在说出真理时，丝毫不会害怕有任何错误。但是，如果你这样说就不对，因为你说的是谎言，他说：\Quote{真理，}他说，\Quote{不在我们内，如果我们说我们没有罪。}他不是说：\Quote{不曾有过；}以免好像是在谈论过去的生活。因为此人过去有过罪；但从他由天主生的那一刻起，他开始没有罪。如果是这样，就没有问题困扰我们。因为我们会说：我们曾是罪人，但现在已成义；我们曾有过罪，但现在没有了。他不是这样说；而是怎么说？\BibleQuote{如果我们说我们没有罪，就是欺骗自己，真理就不在我们内。}然后过了一会儿他又这么说：\BibleQuote{凡由天主生的，不犯罪。}难道\Person{若望}本人不是由天主生的吗？如果若望不是由天主生的，若望——你们听说他躺在主的怀里——有谁敢保证自己身上发生了那\OriginalTerm{重生}{regeneration}，而那男人没有得到，他却有资格躺在主的怀里？主比其他人更爱的那个男人\BibleRef{若 13:23}，难道他独独不是由圣神所生的吗？\par

2. 现在请留意这些话。我正在敦促你们，我们陷入何等困境，以致你们要集中精神，也就是说，通过你们为我们和你们自己祈祷，天主可能扩展我们，给我们出路，免得有人在祂的话语中找到自己灭亡的借口，而这话语原是为了医治和救恩而宣讲和写下的。他说：\BibleQuote{每一个犯罪的人，也做不义的事。}恐怕你们会区分：\BibleQuote{罪就是不义。}恐怕你们会说：我是个罪人，但不是做不义的人，\BibleQuote{罪就是不义。你们也知道，祂显现出来，是为了除掉罪；在祂内没有罪。}祂没有罪，对我们有什么益处呢？\BibleQuote{凡不犯罪的，住在祂内；凡犯罪的，没有见过祂，也不认识祂。孩子们，不要让任何人欺骗你们。行正义的，就是正义的，正如祂是正义的。}我们已经说过，\Quote{正如}一词通常用于某种相似，而不是相等。\BibleQuote{犯罪的是属于魔鬼，因为魔鬼从起初就犯罪。}我们也已经说过，魔鬼没有创造任何人，也没有生育任何人，但他的模仿者可以说是从他生的。\BibleQuote{天主子显现出来，是为了毁灭魔鬼的作为。}因此，祂没有罪，为要解消（或松解）罪。然后接着说：\BibleQuote{凡由天主生的，就不犯罪过，因为天主的种子存留在他内；他不能犯罪，因为他由天主生；}\BibleRef{若一 3:9}他把绳索拉紧了！——似乎，他是就某种罪而说\Quote{不犯罪}，而不是就所有的罪：以至于他所说的\Quote{凡由天主生的，不犯罪}，你可以理解为某一种特定的罪，那由天主生的人不能犯；而这种罪是，如果人犯了它，它就会确认其余的罪。这是什么罪？行违反诫命的事。诫命是什么？\BibleQuote{我给你们一条新诫命，就是你们要彼此相爱。}\BibleRef{若 13:34}请注意！基督的这条诫命被称为\Quote{爱}。藉着这爱，罪被解除。如果这爱不被遵守，不持守它就立即成为一条严重的罪，且是一切罪的根源。\par

3. 你们要留意，弟兄们；我们提出了一点，对于有良好理解力的人来说，问题已经解决了。但难道我们只与跑得更快的人同行吗？那些走得较慢的人不可被撇下。让我们用我们所能的言语多方探讨，好使所有人都能理解。因为我认为，弟兄们，凡不是无端来到教会，不是在教会中寻求属世之物，不是来这里办理世俗事务的人，都是关心自己灵魂的人；他来这里，是为了抓住应许给他的、他能达到的某种永恒之物；他必须思想如何在路上行走，免得被撇下，免得后退，免得走迷，免得因跛行而达不到目标。因此，凡认真的人，让他或慢或快，但不要离开正道。我说这话，是因为说\BibleQuote{凡由天主所生的，就不犯罪}，很可能是指某一种特定的罪；否则便与另一处经文矛盾：\BibleQuote{如果我们说我们没有罪，便是欺骗自己，真理也不在我们内}。这样，问题就可能得到解决。有一种罪，凡由天主所生的人不能犯；这种罪若不犯，其他罪便得赦免；若犯了，其他罪便被确认。这是什么罪呢？就是违反基督的诫命，违反新约。什么是新诫命？\BibleQuote{我给你们一条新诫命：你们要彼此相爱}。\BibleRef{若 13:34}凡违反爱德和兄弟之爱的人，不要妄自夸口说他是由天主所生的；但谁在兄弟之爱中，就有某些罪是他不能犯的，尤其是恨他的弟兄。那么，关于他其他的罪，就如经上所说：\BibleQuote{如果我们说我们没有罪，便是欺骗自己，真理也不在我们内}，他又如何呢？让他从圣经的另一处听取使他心安的教导：\BibleQuote{爱德遮盖许多罪过}。\BibleRef{伯前 4:8}\par

4. 因此，我们推崇爱德；这封书信推崇爱德。主复活后，向伯多禄提出了什么问题？岂不是\BibleQuote{你爱我吗？}\BibleRef{若 21:15}问一次不算；第二次也只问了同一个问题，第三次也没有别的。虽然到了第三次，伯多禄，如同不明白这问题的用意，便忧愁起来，因为主似乎不信他；但无论第一次、第二次还是第三次，主都问了这个问题。三次的畏惧否认，三次的爱德承认。看，伯多禄爱主。他要为主做什么？不要以为他在圣咏中不知如何是好：\BibleQuote{我该用什么报答主赐给我的一切恩惠？}在圣咏中说这话的人，注意到天主为他行了何等大事，便寻求应当用什么报答天主，却找不着。因为无论你要报答什么，都是从祂那里领受的。那么，他找到了什么作为回报呢？我们所说的，弟兄们，正是从他领受的，他找到的只有这个作为回报：\BibleQuote{我要举起救恩之杯，呼求主的名。}因为谁给了他救恩之杯呢？岂不就是他要回报的那一位？现在，举起救恩之杯，呼求主的名，就是充满爱德；而且充满到如此程度，以至你不仅不恨你的弟兄，还准备为你的弟兄而死。这就是圆满的爱德：你准备为你的弟兄而死。主亲自在祂自己身上显示了这一点，祂为众人而死，为那些钉死祂的人祈祷，说：\BibleQuote{父啊，宽恕他们吧，因为他们不知道自己在做什么。}\BibleRef{路 23:34}但若只有祂做了这事，祂就不是师傅了，除非祂有门徒。跟随祂的门徒也做了这事。众人用石头砸死斯德望，他跪下说：\BibleQuote{主，不要将这罪归到他们身上。}\BibleRef{宗 7:59}他爱那些杀害他的人，因为他也是为他们而死。请听宗徒保禄的话：\BibleQuote{我自己，也愿为你们的灵魂耗尽。}\BibleRef{格后 12:15}因为斯德望在死于他们手时，求宽恕的那些人中就有他。这就是圆满的爱德。若有人有如此大的爱德，甚至准备为弟兄而死，那个人就有圆满的爱德。但爱德刚出生时，就已经完全圆满了吗？它出生是为了变得圆满；出生后，它被滋养；滋养后，它被加强；加强后，它被成全；当它达到圆满时，它说什么呢？\BibleQuote{对我来说，活着就是基督，死了就是利益。我渴望离世，与基督同在，那好得无比；然而，肉身存留对你们是必要的。}\BibleRef{斐 1:21}为了那些他准备为之而死的人，他愿意活下去。\par

5. 为使你们知道，这就是那完美爱德，由天主所生的人不违背它，不犯此罪；这是主对\Person{伯多禄}所说的：\Quote{伯多禄，你爱我吗？} 他回答说：\Quote{我爱。} 主没说：你若爱我，就向我行善。因为当主带着可朽血肉时，他曾饥饿、口渴；那时他饥饿口渴时，被人接待为客；那些有财物的人，曾以他们的财产服侍他，正如我们在福音中所读到的。\Person{匝凯}曾接待他作客；他因接待了那位医生而从自己的疾病中痊愈。什么疾病？贪财之病。因为他很富有，又是税吏长。请看这个从贪财之病中痊愈的人：\Quote{我把我财产的一半给穷人；如果我曾向任何人勒索过什么，我以四倍偿还。} \BibleRef{路 19:8} 他保留另一半，不是为了享用，而是为了偿还债务。那时他接待医生作客，因为主有肉身的软弱，人们可以向它行这样的善事；这之所以如此，是因为主愿意赐给那些向他行善的人这同一恩赐；益处乃归于行善的人，而非他自己。因为，侍奉他的天使尚且不需要这些人的善行，难道主反需要吗？连他的仆人\Person{厄里亚}也不需要；天主曾在一个场合藉乌鸦给他送饼和肉 \BibleRef{列上 17:4}，然而为祝福一位虔诚的寡妇，天主的仆人被差遣，那位天主暗中供养的人，竟由寡妇供养。但是，尽管藉着这些天主的仆人，那些体贴他们需要的人为自己获得了益处，因着主在福音中最清楚宣称的那个赏报：\Quote{谁因义人的名接纳一个义人，必得义人的赏报；谁因先知的名接纳一个先知，必得先知的赏报；谁若只因门徒的名，给这些小子中的任何一个一杯凉水喝，我实在告诉你们，他决不失其赏报。} \BibleRef{玛 10:41} 然而，虽然行这些事的人是为自己的益处，但当他即将升天时，却不能再向他行这样的善事。爱主的\Person{伯多禄}能向他献上什么呢？听：\Quote{你牧放我的羊。} \Emph{即}：为弟兄们做那我对你们所做的事。我用我的血赎回了所有人；不要犹豫为真理的宣认而死，好使其余的人效法你。\par

6. 但如上所述，弟兄们，这就是完美爱德。由天主所生的人拥有它。请注意，我所亲爱的，看看我在说什么。看，一个人领受了那重生的圣事，即圣洗圣事；他拥有这圣事，且是伟大的圣事，神圣、圣洁、无法言喻。想一想这是何等圣事！藉着赦免一切罪恶，使他成为一个新人！尽管如此，他还是要好好省察内心，看看那里是否彻底完成了那在身体上所做的事；他要看看自己是否拥有爱德，然后才说：我是由天主所生的。然而，若他没有爱德，他确实有士兵的印记，却像逃兵一样游荡。让他拥有爱德；否则，不要让他说他是由天主所生的。但他说：我拥有圣事。请听宗徒的话：\Quote{我若能说各种奥秘，且有全备的信德，甚至能移山，却没有爱德，我算不了什么。} \BibleRef{格前 13:2}\par

7. 此点，若你们记得，我们在开始诵读此书信时已使你们明白：其中没有比爱德更受称赞的。即使它似乎谈论许多其他事物，但仍回归于此，无论它说什么，终将一切引向爱德。让我们看看它在此是否如此。注意：\BibleQuote{凡由天主所生的，不犯罪。}我们问：什么罪？因为若你理解为一切罪，这与另一处经文相悖：\BibleQuote{若我们说我们没有罪，就是欺骗自己，真理也不在我们内。}那么让他说明是什么罪；让他教导我们；免得我可能轻率地说此处的罪是违反爱德，因为他在上面说过：\BibleQuote{凡恨自己弟兄的，就是在黑暗中，在黑暗中行走，不知道往哪里去，因为黑暗弄瞎了他的眼睛。}\BibleRef{若一 2:11}但或许他在后面说了些什么，而且明确提到了爱德？看，这些话的循环有此目的，有此结果。\BibleQuote{凡由天主所生的，不犯罪，因为天主的种子留在他内。}\BibleRef{若一 3:9}天主的\Quote{种子}，\Emph{即}圣言：因此宗徒说：\BibleQuote{我在基督耶稣内藉福音生了你们。他不能犯罪，因为他是由天主所生的。}\BibleRef{格前 4:15}让他告诉我们这事，让我们看看我们在何事上不能犯罪。\BibleQuote{在此天主的孩子与魔鬼的孩子显明出来：凡不行正义的，就不是出于天主，不爱自己弟兄的，也是如此。}\BibleRef{若一 3:10}啊，如今确实显明了他所说的：\BibleQuote{不爱自己弟兄的。}因此，唯有爱德区分天主的孩子与魔鬼的孩子。让他们众人都以基督十字架的记号划十字；让他们都回答\Quote{阿们}；让众人高唱\Quote{阿肋路亚}；让众人领受洗礼；让众人来到教会；让众人建造教会的墙壁：但只有藉着爱德，才能分辨天主的孩子与魔鬼的孩子。凡有爱德的，是由天主所生的；凡没有的，不是由天主所生的。一个伟大的标记，一个伟大的区分！任你拥有什么，若唯独缺少此，对你毫无益处；若缺少其他，若拥有此，你便成全了法律。\BibleQuote{因为爱人的，就成全了法律，}宗徒说；再者，\BibleQuote{爱德是法律的满全。}\BibleRef{罗 13:8}我认为这就是福音中那商人寻找的珍珠，他\BibleQuote{找到一颗宝贵的珍珠，就去卖掉一切所有的，买了它。}\BibleRef{玛 13:46}这就是无价珍珠——爱德，没有它，无论你有什么，对你毫无益处；若你独有此，便足够了。如今，你们藉信德看见，然后必以实际的看见而看见。因为若我们在看不见时爱，那么看见时将如何拥抱它！但我们必须在何事上操练自己？在兄弟之爱上。你或许对我说：我没有看见天主；你能对我说：我没有看见人吗？爱你的弟兄。因为若你爱那看得见的弟兄，同时你也要看见天主；因为你将看见爱德本身，天主内住在其中。\par

8. \BibleQuote{凡不行正义的，就不是出于天主，不爱自己弟兄的，也是如此。}\BibleRef{若一 3:10}\BibleQuote{这就是传报的信息：}请注意他如何证实：\BibleQuote{我们从起初就听到的信息，就是我们该彼此相爱。}他已向我们明确说明他所指的正是此事；凡违犯此诫命的，就陷于那被咒诅的罪中，那些不是由天主所生的人所堕入的。\BibleQuote{不像加音，他本就是属于那恶者，杀害了自己的弟弟。他为什么杀害了他？因为自己的行为是恶的，而弟弟的行为是正义的。}\BibleRef{若一 3:12}因此，哪里存有嫉妒，兄弟之爱便不能存在。请注意，我所爱的。凡嫉妒的，就不爱。那人的罪是魔鬼的罪；因为魔鬼因嫉妒而使人堕落。他堕落，便嫉妒那站立的。他并非希望自己站住而使人堕落，只是希望自己不独自堕落。从他所附属的这些，你们要牢记：嫉妒不能存在于爱德中。你们清楚知道，当爱德被赞扬时，\BibleQuote{爱德不嫉妒。}\BibleRef{格前 13:4}在\Person{加音}内没有爱德；若在\Person{亚伯尔}内没有爱德，天主就不会接纳他的祭献。因为当他们两人都献了祭，一个献了地里的出产，另一个献了羊群中的首生；弟兄们，你们想，是天主轻视地里的出产，而喜爱羊群的首生吗？天主不看人的手，而看内心：祂看见谁以爱德奉献，便垂顾他的祭献；看见谁以嫉妒奉献，便从祂的祭献转开眼目。因此，论及\Person{亚伯尔}的善行，他仅指爱德；论及\Person{加音}的恶行，他仅指他对自己弟兄的憎恨。他恨自己的弟兄并嫉妒他的善行还不够，因为他不愿效法，竟愿杀害。由此显示他是魔鬼的孩子，也显示另一个是天主的义人。故此人被辨别，我的弟兄们。不要看人的口舌，而要看行为与内心。若有人不为自己的弟兄行善，他显露出自己内里所有。人藉试探得以考验。\par

9. \BibleQuote{诸位弟兄，不要惊奇，如果世界恨我们。}必须常常告诉你们什么是\Quote{世界}吗？不是天，不是地，也不是天主造的这些有形之物，而是爱世界者。因我常讲这些，有人觉得我烦；但我绝非无缘无故地讲，以至于有人可以问我是否讲过，而他们答不上来。那么，即便强行灌输，也让听者心中有所持守。什么是\Quote{世界}？\Quote{世界}用于贬义时，就是爱世界者；\Quote{世界}用于赞美时，是指天和地以及其中的天主造化；因此经上说：\BibleQuote{世界是藉着祂造成的。}\BibleRef{若 1:10} 同样，世界也是大地的丰盈，正如若望自己所说：\BibleQuote{祂不仅是我们的罪的赎罪者，也是全世界的（罪的）赎罪者。}\BibleRef{若一 2:2} 他所说的\Quote{世界}是指分散在整个大地上的所有信徒。而\Quote{世界}用于贬义时，就是爱世界者。爱世界者不能爱自己的弟兄。\par

10. \BibleQuote{如果世界恨我们，我们知道}——我们知道什么？——\BibleQuote{我们已经出死入生了}——我们如何知道？\BibleQuote{因为我们爱弟兄。}\BibleRef{若一 3:14} 不要问任何人；各人当回到自己的心：若在那里找到弟兄之爱，就让他放心，因为他已经\Quote{出死入生了}。他已经站在右边；不要在意现在他的荣耀是隐藏的；当主降临的时候，他要在荣耀中显现。因为他里面有生命，但现在尚在冬季；根是活的，但枝条可以说是枯干的；里面是有生命的实质，里面有树的叶子，里面有果实；但它们等待夏天。所以，\BibleQuote{我们知道我们已经出死入生了，因为我们爱弟兄。不爱的人，就留在死亡中。}\BibleRef{若一 3:14} 恐怕你们以为仇恨或不爱是小事，弟兄们，听听下文：\BibleQuote{凡恨他弟兄的，就是杀人者。}\BibleRef{若一 3:15} 那么，若有人轻看恨弟兄，难道他也轻看杀人吗？他并未动手杀人，但在天主面前已被视为杀人者；那人还活着，而这人已被判为凶手！\BibleQuote{凡恨弟兄的，就是杀人者；你们知道，凡杀人者，没有永生住在他里面。}\BibleRef{若一 3:15}\par

11. \BibleQuote{由此我们知道爱：}\BibleRef{若一 3:16} 他指的是爱的圆满，就是我们嘱咐你们留心的那种圆满：\BibleQuote{由此我们知道爱，就是祂为我们舍命；我们也应当为弟兄舍命。}\BibleRef{若一 3:16} 看，这正是那话的来由：\BibleQuote{伯多禄，你爱我吗？你喂养我的羊。}\BibleRef{若 21:15} 因为，为使你们知道祂愿意他如此喂养祂的羊，以至于他应为羊舍命，祂立刻对他说：\BibleQuote{你年少时，自己束上带，随意往来；但年老时，你要伸出手来，别人要给你束上，带你到你不愿意去的地方。}\BibleRef{若 21:18} 福音作者说：\BibleQuote{这话是表明他将以怎样的死荣耀天主；}\BibleRef{若 21:19} 如此，祂曾对他说\Quote{喂养我的羊}的那位，祂也要教导他为祂的羊舍命。\par

12. 爱德从何处开始呢，弟兄们？\OriginalTerm{爱德}{charity} 请稍加留意：它的圆满你们已经听到了；它的终极和尺度，就是主在福音中给我们指出的：\BibleQuote{人若为自己的朋友舍命，再没有比这更大的爱。}\BibleRef{若 15:13} 所以，祂已将爱德的圆满放在福音中，这里也将其圆满呈现在我们面前。但你们会自问，对自己说：我们什么时候才能拥有\Emph{这}爱德呢？不要过早地对自己失望。也许，爱德已出生但尚未圆满；要滋养它，免得被窒息。但你或许会对我说：我凭什么知道它呢？因为它的圆满我们已经听过；现在让我们听听它的开端。祂接着说：\BibleQuote{凡有今世财物的，看见弟兄有缺乏，却关闭他慈悲的心肠，天主的爱怎能住在他内呢？}\BibleRef{若一 3:17} 看，爱德就从这里开始！即使你尚不能为弟兄舍命，现在也该能把你所有的给弟兄。现在就让爱德触动你的心肠，使你不是出于虚荣，而是出于慈悲至深之处；你要想到他现在的急需。因为如果你不能把多余的给弟兄，你能为弟兄舍命吗？你的金钱藏在怀中，盗贼可能夺去；即使盗贼不夺，你死了也要留下它，即便它活着时不离开你；你要用它做什么呢？你的弟兄饥饿，他处于困乏中；或许他焦虑不安，被债主逼迫；他是你的弟兄，你们同是被赎回的，为你们付上的代价相同，你们都因基督的血得救赎；看看你是否怜悯，如果你有今世的财物。也许你说：\Quote{这与我何干？要我拿出钱来，免得他受苦吗？}如果你的心这样回答你，圣父的爱就不在你内。如果圣父的爱不在你内，你就不是由天主生的。你凭什么夸口说自己是基督徒？你有这名号，却没有行为。但若行为不随从名号，即便任何人称你为外邦人，也要用行为证明你是基督徒。因为若你以行为不显明自己是基督徒，众人仍可称你为基督徒；但名号有什么益处呢？那实际并不存在。\BibleQuote{凡有今世财物的，看见弟兄有缺乏，却关闭他慈悲的心肠，天主的爱怎能住在他内呢？}\BibleRef{若一 3:17} 接着他说：\BibleQuote{我的孩子们，我们不要只在言语上爱，也不只在舌头上，而要行为上、真理上爱。}\BibleRef{若一 3:18}\par

13. 我想，如今这事已向你们显明了，我的弟兄们：这伟大而又最关紧要的奥秘与奥迹。凡\WorkTitle{圣经}都昭示了爱德的力量；但我不知道是否在任何其他经卷中，比在这封书信里更充分地昭示了它。我们在主内恳求并劝请你们，要将已听见的记在记忆中，并以热忱聆听那尚待讲述的，直到这封书信完结。但要为那好种子敞开你们的心：拔出荆棘，免得我们撒在你们里的被窒息，反倒使庄稼生长，使田园主欢喜，并为你们预备谷仓如同为麦子，而非为糠秕预备火。\par

\SourceRecord{Translated by H. Browne.；Nicene and Post-Nicene Fathers, First Series；Vol. 7.；Edited by Philip Schaff.；Buffalo, NY: Christian Literature Publishing Co.,；1888.；<http://www.newadvent.org/fathers/170205.htm>.}

% structure: p0001 source=h1 target=section reason=page-title

\section{论若望一书第六篇讲道}

% structure: p0002 source=h2 target=subsection reason=relative-heading-rank; base=h2

\subsection{\BibleRef{若一 3:19-4:3}}

\BibleQuote{我们也知道我们属于真理，并在祂面前安定我们的心。因为如果我们心中责备我们，天主比我们的心更大，且知道一切。亲爱的，如果我们心中不责备我们，我们就对天主有信赖。并且无论我们求什么，我们都从祂领受，因为我们遵守祂的诫命，并做祂眼中喜悦的事。这就是祂的诫命：相信祂的圣子耶稣基督的名，并彼此相爱，正如祂所命令我们的。那遵守祂诫命的，就住在祂内，祂也住在他内。由此我们知道祂住在我们内，是借着祂赐给我们的圣神。亲爱的，你们不要信每一神，但要考验那些神是否出于天主，因为有许多假先知已进入这个世界。从此你们能认出天主的神：凡宣认耶稣基督是成了血肉而来的，就是出于天主；凡不宣认耶稣基督是成了血肉而来的，就不是出于天主；这就是那敌基督者，你们曾听说他要来，现在已经在世界上了。}\par

1. 弟兄们，如果你们记得，我们昨天在这句话上结束了讲道，即\BibleRef{若一 3:18}，这句话无疑应当并确实应当存留在你们心中，因为这是你们最后听到的。\BibleQuote{小孩子们，我们爱不要只用言语，也不可只用口舌，而要用行动和事实。}接着他说：\BibleQuote{我们也知道我们属于真理，并在祂面前安定我们的心。}\BibleQuote{因为如果我们心中责备我们，天主比我们的心更大，且知道一切。}他曾说：\BibleQuote{我们爱不要只用言语，也不可只用口舌，而要用行动和事实。}我们问，在什么工作或什么真理上能认出那爱天主或爱弟兄的人呢？上面他曾说过爱德达到的限度：主在福音中说：\BibleQuote{人没有比这更大的爱，即一个人为朋友舍掉自己的生命。}\BibleRef{若 15:13}宗徒也说了同样的话：\BibleQuote{正如祂为我们舍掉了生命，我们也应当为弟兄舍掉生命。}\BibleRef{若一 3:16}这就是爱德的圆满，再不能有更大的了。但因为这不是在所有人中都圆满，且那在自己身上不圆满的人不应该绝望，如果那能够圆满的已经诞生：当然，如果诞生了，就必须被哺育，并借着适合的养料被带到它应有的圆满：所以，我们询问了爱德的起始，它从哪里开始，并立刻发现：\BibleQuote{谁若有世上的财物，看见弟兄有需要，却关闭了自己的恻隐之心，圣父的爱怎能住在他内？}\BibleRef{若一 3:17}所以，我的弟兄们，爱德从这里开始：将自己的多余给那有需要的人，给那处于任何困境的人；用自己的暂时富足解救弟兄的暂时苦难。这里是爱德的第一阶段。这样开始了，如果你们用天主圣言和来世生命的希望来滋养它，你们终将到达那圆满，即准备好为你们的弟兄舍掉生命。\par

2. 但是，因为许多此类事情是由那些寻求其他对象、不爱弟兄的人所做的；让我们回到良心的证言。我们如何证明许多此类事情是由不爱弟兄的人所做的呢？在异端和分裂中，有多少人自称殉道者！他们自认为是为弟兄舍命。如果他们是为弟兄舍命，就不会与整个兄弟团体分离。再者，有多少人为了虚荣而大量施舍、大量给予，在其中只寻求人的称赞和世俗的荣誉，这充满了虚浮，毫无稳固！既然如此，有这样的人，那么兄弟之爱的证据何在？既然他愿这事被证明，并以劝诫的方式说：\Quote{我的孩子们，我们不要只在言语和舌头上爱，而要藉行为和真理来爱；}我们问，是藉什么行为、什么真理？有什么行为比施舍穷人更明显的吗？许多人这样做是出于虚荣，而非出于爱。有什么行为比为弟兄舍命更伟大的吗？这也一样，许多人愿意被人认为做了这事，但他们这样做是出于虚荣以求名誉，而非出于爱心的深处。因此，那在独有天主看见的隐密处，在天主面前，确证自己的心，并省察自己的心是否确实出于对弟兄的爱而行此事的人，才是爱他的弟兄；他的见证便是那洞察内心的眼睛，是人不能窥视的。因此，宗徒保禄因预备为弟兄们舍命，说：\Quote{我甘愿为你们的灵魂付出自己，}\BibleRef{格后 12:15}然而，因为只有天主在他心中看见这事，而非他所说话的那些必死之人，他就对他们说：\Quote{但对我来说，受你们审断，或受人的法庭审判，是极小的事。}\BibleRef{格前 4:3}同一宗徒也在某处表明，这些事常因空洞的虚荣而行，而非建立在爱的坚实基础上：因为当他论及爱德的赞美时说：\Quote{我若把我所有的财产分施给穷人，并舍身投火被焚，但我若没有爱德，为我毫无益处。}\BibleRef{格前 13:3}一个人可能没有爱德而做这事吗？可能的。因为那些分裂合一的人，是没有爱德的人。在那里寻找，你会看见许多人大量施舍穷人；会看见另一些人预备迎接死亡，甚至在没有迫害的地方投身其中：这些人无疑是没有爱德而做此事的。那么，让我们回到良心，宗徒论良心说：\Quote{因为我们的夸耀正是这个：我们良心的证言。}\BibleRef{格后 1:12}让我们回到良心，同一宗徒又说：\Quote{但每人要察验自己的行为，这样，他只在自身内有夸耀，而不在别人内。}\BibleRef{迦 6:4}因此，我们每个人都当\Quote{察验自己的行为}，看它是否从爱德的脉络涌流而出，是否以爱德为根，使善行如枝条般发芽生长。\Quote{但每人要察验自己的行为，这样，他只在自身内有夸耀，而不在别人内，}不是当别人的舌头为他作证时，而是当他自己的良心作证时。\par

3. 那么，这就是他在这里强调的。\Quote{在这一点上，我们知道我们是属于真理的，当我们在行为和真理上}我们爱，\Quote{不仅是在言语和舌头上，并在祂面前安心。}\BibleRef{若一 3:19}\Quote{在祂面前}是什么意思？就是祂看见的地方。因此主自己在福音中说：\Quote{你们要小心，不可在人前行你们的义，为叫他们看见；否则，你们在天上你们之父那里就没有赏报了。}而\Quote{不要让左手知道右手所行的}是什么意思？岂非右手代表纯洁的良心，左手代表世俗的贪欲？许多人因贪恋世俗而行许多奇事：是左手在作工，而非右手。右手应当作工，且不让左手知道，好使当我们因爱而行任何善事时，世俗的贪欲甚至不能混入其中。我们怎样才能知道这一点？你是在天主面前：省察你的心，看你做了什么，其中的意图是什么：是你的救恩，还是人虚浮的赞美？往内看，因为人不能判断他看不见的人。如果我们\Quote{安心，}就让它\Quote{在祂面前。}因为\Quote{如果我们的心责备我们，}\Emph{即}在心中控告我们，说我们没有以应有的心志做这事，\Quote{天主比我们的心大，且知道一切。}你向人隐藏你的心：你若能，向天主隐藏看看。你怎能向祂隐藏呢？有罪人怀着敬畏和忏悔向祂说：\Quote{我往何处去，才能躲避你的神？从你的面前，我往何处逃？}他寻找逃避天主审判的路，却找不着。因为天主不在何处？\Quote{我若升到天上，}他说，\Quote{你在那里；我若下到地狱，你也在那里。}你要往哪里去？你往哪里逃？你愿听劝告吗？你若想逃离祂，就逃向祂。藉着忏悔逃向祂，而非藉着隐藏逃离祂：你不能隐藏，但你能忏悔。向祂说：\Quote{你是我的避难所；}让爱在你内被滋养，这爱唯独引向生命。让你的良心为你作证，你的爱是出于天主的。如果是出于天主的，不要愿在人前显示它；因为人的赞美不能把你提升到天堂，人的指责也不能把你从那里推下去。让那加冕于你的祂看见：让祂作你的见证，你正是由祂这位审判者加冕的。\Quote{天主比我们的心大，且知道一切。}\par

4. \Quote{亲爱的，若我们的心不责备我们，我们在天主前便放心大胆：} \BibleRef{若一 3:21} —— 所谓\Quote{我们的心不责备}，是什么意思呢？就是我们的心给我们真实的回答：我们爱，并且在我们内有真诚的爱；不是虚假的，而是诚实的；寻求弟兄的救恩，不从弟兄身上期望任何利益，只期望他的救恩—— \Quote{我们在天主前放心大胆：并且我们无论求什么，必由祂获得，因为我们遵守祂的诫命。} \BibleRef{若一 3:21} —— 因此，不是在人的眼前，而是在天主亲自看见的地方，就是在心中—— \Quote{我们放心大胆，}接着，\Quote{在天主前：并且我们无论求什么，必由祂获得：}然而，是因为我们遵守祂的诫命。什么是\Quote{祂的诫命}？我们总得重复吗？\Quote{我给你们一条新命令，你们当彼此相爱。} \BibleRef{若 13:34} 他所说的正是爱德本身，他劝勉的正是这个。凡有兄弟爱德，并在天主前——在天主看见的地方——拥有它，且他的心在正义的审查下只给他一个回答，即爱德真诚的根源就在那里，以结出美果；这样人在天主前放心大胆，并且无论求什么，必由祂获得，因为他遵守祂的诫命。\par

5. 这里有一个问题摆在我们面前：因为所讨论的不是这个或那个人，也不是你或我——因为如果我向天主祈求某事而未获得，任何人都可能轻易说我：\Quote{他没有爱德：}对于现世任何人，都能轻易如此说；且任人按自己的意思论断人：——不是我们，而是那些人更被讨论，那些被公认为圣人、写作时是圣人、如今与天主同在的人。谁有爱德呢，如果保禄也没有，他曾说：\Quote{我们的口向你们格林多人是张开的，我们的心是宽大的；你们在我们内并不狭窄，}他曾说：\Quote{我甘愿为你们的灵魂耗尽自己，}他内里有如此大的恩宠，显明他有爱德。然而我们发现他祈求了却没有获得。我们还能说什么呢，弟兄们？这是一个问题：专心注视天主：这也是一个重大的问题。正如关于罪，经上说：\Quote{凡由天主生的，就不犯罪，}我们发现了这罪就是违反爱德，并且这正是那处经文严格所指：同样，我们现在询问他想要说什么。因为若只看字句，似乎清楚；若将例子考虑进来，则晦暗不明。这些话没有更清楚的了。\Quote{并且我们无论求什么，必由祂获得，因为我们遵守祂的诫命，行祂眼中喜悦的事。} \Quote{我们无论求什么，}他说，\Quote{必由祂获得。}他使我们极为为难。在另一处他也曾使我们为难，如果他意指一切罪的话：但那时我们找到了解释的空间，即他指的是某一种罪，而非一切罪；不过是一种罪，\Quote{凡由天主生的，就不犯}这罪；并且我们发现这罪不是别的，正是违反爱德。我们还有福音中的一个明显例子，当主说：\Quote{我若没有来，他们就没有罪。} \BibleRef{若 15:22} 怎么样？难道因为主这样说，犹太人当祂来时就是无罪的？那么，如果祂没有来，他们就没有罪了？难道医生的临在反使人患病，而不是退烧？连疯子也不会这样说吧？祂来正是为了医治和治愈病人。因此，当祂说：\Quote{我若没有来，他们就没有罪，}祂要我们理解的，难道不是某种特定的罪吗？因为有一种罪是犹太人本不会有的。什么罪？就是他们不信祂，当祂来了他们却拒绝祂。正如祂在那里说\Quote{罪，}并不意味我们要理解为一切罪，而是特定的罪：同样，这里也不是一切罪，以免与那处经文\Quote{若我们说我们没有罪，我们就是欺骗自己，真理也不在我们内} \BibleRef{若一 1:8} 相矛盾；而是特定的罪，即违反爱德。但在此处他更紧地约束了我们：\Quote{我们若求，}他说，\Quote{倘若我们的心不控告我们，并在天主面前回答我们说，真爱在我们内，} \Quote{我们无论求什么，必由祂获得。}\par

6. 好了：我已经告诉过你们，我亲爱的弟兄们，不要转向我们。因为我们是什么？你们又是什么？不过是众所周知的教会罢了？并且，若祂乐意，我们就在那教会内；我们中那些借着爱住在其中的，让我们坚持在那里，如果我们愿意显示我们所拥有的爱。但宗徒保禄，我们该对他有什么坏的想法呢？他不爱弟兄们！他不内里有在临于天主前良心的见证！保禄没有在他内那一切美善果实所出的爱德之根！哪个疯子会说这个？那么好了：我们在哪里发现这位宗徒祈求了却没有获得呢？他亲自说：\Quote{恐怕因了启示的丰富而过于高举我，就有一根刺加在我身上，就是撒殚的使者来击打我。为此，我曾三次求主叫我脱离这事。祂对我说：我的恩宠为你够了，因为德能在软弱中才成全。} \BibleRef{格后 12:7} 看，他祈求那\Quote{撒殚的使者}从他身上拿走，却没有被垂听。为什么？因为对他无益。他被垂听了，是为了救恩，当他没有被按他的意愿垂听时。要知道，我亲爱的，一个伟大的奥秘：我们特意提醒你们注意，免得你们在诱惑中失掉。圣人们的一切祈求都为了救恩而被垂听：他们总是在关乎他们永恒救恩的事上被垂听；这是他们所渴望的：因为在这事上，他们的祈祷常被垂听。\par

7. 但让我们区分天主垂听祈祷的不同方式。因为我们发现有些人不为自己的愿望而蒙垂听，却为救恩而蒙垂听：又有些人我们发现有为自己愿望而蒙垂听，却不为了救恩。请留意这个区别，紧记这个例子：一个人不为他的愿望而蒙垂听，却为救恩而蒙垂听。请听宗徒保禄；因为天主亲自向他显示，何为垂听祈祷归于救恩：\Quote{我的恩宠为你足够了，}祂说，\Quote{因为能力在软弱中才达到圆满。}你曾祈求，曾呼喊，曾三次呼喊；你所发出的呼声我一次全都听见了，我没有向你掩耳；我知道我该做什么：你希望它被拿走，那使你燃烧的治愈之物；我知道你所背负的软弱。那么，这里有一个人，他为了救恩而蒙垂听，但就他的意愿而言，他并未蒙垂听。我们在哪里能找到那些为意愿而蒙垂听、却不为了救恩的人呢？我们想，我们能找到某个邪恶、不虔敬的人，为他的意愿蒙天主垂听，却不为了救恩吗？如果我向你举出某人的例子，你或许会对我说：\Quote{是你称他为邪恶，因他曾是义人；他若不是义人，他的祈祷就不会被天主垂听。}我要举的例子是这样一个，他的不义和不虔敬无人可以怀疑。魔鬼本身：他求得了约伯，而且得到了。\BibleRef{约 1:11}这里你们不也听说过关于魔鬼，\Quote{犯罪的是属于魔鬼}？不是魔鬼创造的，而是罪人效仿。不也说关于他\Quote{他不站在真理中}？\BibleRef{若 8:44}他岂不就是\Quote{那古蛇}，藉着女人把毒酒灌给最初的人？\BibleRef{创 3:1}甚至在约伯的事上，为他保留了他的妻子，使她可以试探而不是安慰她的丈夫？魔鬼求得了圣人，要试探他，并且得到了：宗徒求那肉中刺从他身上拿走，却没有得到。但是宗徒比魔鬼更蒙垂听。因为宗徒是为救恩而蒙垂听，虽然不遂其愿；魔鬼是为其意愿而蒙垂听，但是为了灭亡。因为约伯被交给他试探，是为了藉着他的坚持，使魔鬼受折磨。但是，各位弟兄，这不仅在旧约经书中找到，也在福音中。当主从那人身上驱逐鬼魔时，鬼魔恳求主准许他们进入猪群。难道主没有能力命令他们甚至不能靠近那些动物吗？因为若祂不准许，他们不会反叛天地之王。但是出于某种奥秘，某种更深的意义，祂让鬼魔进入猪群，以显示魔鬼在那些过着猪一样生活的人身上掌权。于是鬼魔的请求得到垂听；难道宗徒没有被垂听？或者（更确切地说）我们是否应该这样说：宗徒被垂听，鬼魔没有被垂听？他们的意愿得以实现，他的益处得以成全。\par

8. 与此相符，我们应当明白，天主虽然不按我们的意愿赐予，却为我们的救恩而赐予。因为假设你所求的事对你有害，而医生知道它对你有害，那么怎样？不能说医生没有垂听你的请求，当你也许求冷水时，若对你有益，他立刻给你，若无益，他不给。难道他对你的请求充耳不闻，或者更确切地说，他乃是为你的益处而垂听，即便他拒绝了你的意愿？那么，各位弟兄，愿爱德在你们内；愿它存于你们内，然后你们就可以安心：即使你们所求的事没有赐予你们，你们的祈祷也蒙应允，只是你们不知道。有许多人被交在他们自己手中，以致于受害；关于他们，宗徒说：\BibleQuote{天主就任随他们陷于心中的情欲}\BibleRef{罗 1:24}。有人求了一大笔钱；他得到了，却受害了。当他没有时，他几乎无所畏惧；一旦拥有，他就成了更强者的猎物。难道那人的请求不是使他受害的应允吗？他偏要拥有那会使强盗来追逐他的东西，而当他贫穷时，无人追逐他。要学习恳求天主，使你把自己交托给医生，让他做他知道最好的事。承认疾病，让他施以治疗。只持守爱德。因为他必然要切割，必然要烧灼；你若喊叫，在切割、烧灼和磨难下哭泣却不得宽免，然而他知道腐烂有多深。你甚至希望他现在就松手，而他只考虑创口的深度；他知道要走多远。他不按你的意愿垂听你，而是为你的痊愈垂听你。所以，各位弟兄，要确信宗徒所说的是真的：\BibleQuote{因为我们不知道我们该如何祈求才对，但圣神却亲自以不可言喻的叹息为我们转求；因为祂为圣徒们转求}\BibleRef{罗 8:26}。为何说\Quote{圣神亲自为圣徒们转求}？意思岂不是圣神在你们内所成就的爱德？因为同一位宗徒又说：\BibleQuote{天主的爱德借着所赐给我们的圣神，已倾注在我们心中了}\BibleRef{罗 5:5}。是爱德在叹息，是爱德在祈祷；那赐予爱德的，不能向爱德掩耳。你们安心吧：让爱德祈求，天主的耳朵就在那里。不是你所愿望的被成就，而是那有益的被成就。因此，他说：\Quote{我们无论求什么，必从祂那里获得}，我已经说过，如果你理解为\Quote{为了救恩}，那就没有问题；若不是为了救恩，就有问题，而且是严重的问题，这问题使你成为宗徒保禄的控告者。\Quote{我们无论求什么，必从祂那里获得，因为我们遵守祂的诫命，作祂喜悦的事：}在内里，在祂看见的地方。\par

9. 这些诫命是什么？\Quote{这就是，}他说，\Quote{祂的命令：要我们信祂的子耶稣基督的名，并彼此相爱。}\BibleRef{若一 3:23} 你看，这就是诫命：你看，凡违反这条诫命的，就是犯那\Quote{凡由天主所生的}免犯的罪。\Quote{正如祂给我们命令：}要我们彼此相爱。\Quote{而遵守祂命令的}\BibleRef{若一 3:24}——你看，我们受命的无非是彼此相爱——\Quote{遵守祂命令的，就住在祂内，祂也住在他内。}我们藉着祂所赐给我们的圣神，知道祂住在我们内。这岂不显然就是圣神在人内所成就的，即叫人在内里有爱和爱德？岂不显然，如宗徒保禄所说：\Quote{天主的爱藉着所赐给我们的圣神，倾注在我们心中}\BibleRef{罗 5:5}？因为[我们的宗徒]正在谈论爱德，并说我们应当在天主眼前省察自己的心。\Quote{但如果我们的心不责备我们：}\Emph{即}若它承认，在我们内所做的一切善工，都是出于对弟兄的爱。此外，他在谈论这命令时这样说：\Quote{这就是祂的命令：要我们信祂的子耶稣基督的名，并彼此相爱，正如祂给我们命令。}\Quote{遵守祂命令的，就住在祂内，祂也住在他内。我们藉着祂所赐给我们的圣神，知道祂住在我们内。}你若真发现自己有爱德，你就有天主的圣神，好能明了：因为这是一件非常必要的事。\par

10. 在起初，\Quote{圣神降临在信者身上：他们就说起别国的话来}，是他们未曾学过的，\Quote{按圣神所赐的口才说起话来。}\BibleRef{宗 2:4} 这些是适合那时期的标记。因为必须有一种预像，显示圣神在万语中，表明天主的福音要藉万语传遍普世。这事为作预像而行，随即就过去了。如今在覆手时，为使人领受圣神，我们岂期待他们说起别国的话来？或当我们为这些婴儿覆手时，你们每人岂都察看他们是否说起别国的话来？若见他们不说，岂有人错谬到如此地步，说：这些人没有领受圣神；因为若领受了，他们必说起别国的话来，如当时所发生的那样？若如今圣神临在的见证不再藉这些奇迹赐下，那么藉什么赐下？人凭什么知道自己领受了圣神？让他省察自己的心。若他爱弟兄，天主的圣神就住在他内。让他察看，在天主眼前检验自己，看他内里是否有对和平与合一的爱，对那遍布普世的教会的爱。不要只停留在爱眼前的弟兄，因为许多弟兄我们并未看见，但在圣神的合一中，我们与他们相联。他们与我们不在一处，有什么奇怪呢？我们同属一身，同有一个头，在天上。弟兄们，我们的两只眼睛并不彼此看见；可谓彼此不相识。但在身体结构的爱德中，它们岂不相识？为向你们表明，在结合它们的爱德中，它们彼此相识：当两眼都睁开时，右眼若不注视某一对象，左眼也必不同样注视。你试试能否单用右眼看而不用左眼。它们一同注视同一对象，一同转向同一对象：它们的目标相同，位置不同。那么，凡与你同爱天主的，若目标与你相同，你就不必在乎在身体上你们被地方分开；你们心中之目同等地定在真理之光上。因此，若你要知道自己已领受圣神，就省察你的心：免得你徒有圣事，却没有圣事的德能。省察你的心。若对你的弟兄有爱，便可安心。没有天主的圣神，不可能有爱；因为保禄呼喊道：\Quote{天主的爱藉着所赐给我们的圣神，倾注在我们心中。}\BibleRef{罗 5:5}\par

11. \BibleQuote{亲爱的，不要信从任何神。}\BibleRef{若一 4:1} 因为祂曾说过：\Quote{我们所以知道祂存留于我们内，是藉着祂所赐给我们的神。} 但如何认识这同一圣神？注意这点：\BibleQuote{亲爱的，不要信从任何神，但要考验那些神，看它们是否出于天主。} 谁验证这些神？祂给我们出了一道难题，我的弟兄们！幸好祂要亲自告诉我们如何辨别。祂即将告诉我们：不要怕；但先看；注意：看，这正是狂妄的异端者用以嘲笑我们的事。注意，看祂说什么：\BibleQuote{亲爱的，不要信从任何神，但要考验那些神，看它们是否出于天主。} 在福音中，主曾以水之名称呼圣神，祂\BibleQuote{大声说：\InnerQuote{谁若渴，就到我这里来喝！信从我的人，从他的心中要流出活水的江河。}}\BibleRef{若 7:37} 但圣史解释了祂这话是指什么：因为他继续说：\Quote{祂说这话，是指那信祂的人将要领受的神。} 为什么主没有给许多人施洗？但祂说什么？\Quote{因为圣神还没有赐下，因为耶稣还没有受到光荣。} 那时，那些人领受了洗礼，却还没有领受圣神，就是主在五旬节从天上所派遣的，所以他们先等候主受到光荣，然后圣神才被赐下。甚至在祂受光荣之前，在祂派遣圣神之前，祂已邀请人预备自己领受祂所说的水：\Quote{谁渴，就让他来喝；} 并说：\Quote{信从我的人，从他的心中要流出活水的江河。} 什么意思？\Quote{活水的江河}？那水是什么？不要问我；去问福音。\Quote{但这}，它说，\Quote{祂是指那信祂的人将要领受的神说的。} 所以，圣事的水是一回事；象征圣神的另一种水是另一回事。圣事的水是可见的；圣神的水是不可见的。\Emph{那}水洗涤身体，并象征灵魂中所发生的。藉着\Emph{这}圣神，灵魂本身被洁净和滋养。这就是圣神，异端者和一切自绝于教会的人不能拥有。而那些没有公开自绝，却因邪恶而被断绝，虽然在教会内却像糠秕一样旋转而不是麦粒的人，也没有这圣神。主以水之名表示这圣神：我们从这封书信中听到：\Quote{不要信从任何神；} 撒罗满的话也作证：\Quote{远离陌生之水。} 什么意思？\Quote{水}？圣神。水是否总是意味着神？不总是：但在一些地方它意味着圣神，在一些地方它意味着圣洗，在一些地方意味着人民\BibleRef{默 17:15}，在一些地方意味着劝告：正如你在某处读到：\Quote{劝告是生命之泉，为拥有它的人。}\BibleRef{箴 16:22} 所以，在圣经的各处，\Quote{水}这个词意味着不同的事物。然而现在，通过水这个词，你听到圣神被提及，不是根据我们的解释，而是根据福音的见证，其中说：\Quote{但祂说这话，是指那信祂的人将要领受的神。} 如果这样，以水之名指的是圣神，而这封书信告诉我们：\Quote{不要信从任何神，但要考验那些神，看它们是否出于天主；} 我们要明白，这就是所说的话：\Quote{远离陌生之水，不要喝陌生之泉。} 什么意思？\Quote{不要喝陌生之泉}？不要信从陌生之神。\par

12. 于是剩下一个考验，用以证明它是圣神。他确实设立了一个记号，而且似乎很难：但我们来看看。我们得回归那爱德；正是爱德教导我们，因为它是傅油。然而，他在这里说什么？\Quote{考验那些神，看它们是否出于天主，因为有许多假先知已经来到这世界上。} 现在，这里指一切异端者和一切分裂者。那么，我如何证明那圣神？他接着说：\Quote{在这件事上，圣神可以被认识。} 唤醒你内心的耳朵。我们不知所措；我们曾说：谁知道？谁辨别？看，他即将告诉我们记号。\BibleQuote{天主的神是这样认识的：凡明认耶稣基督是在肉身内降世的神，便是出于天主；凡不明认耶稣基督是在肉身内降世的神，便不是出于天主；这就是假基督，你们曾听说祂要来，而现在已经在世界上了。}\BibleRef{若一 4:2} 我们的耳朵，可以说，正警觉着辨别诸神；但我们被告知的事，反而让我们一点也不更明白。因为他说什么？\Quote{凡明认耶稣基督是在肉身内降世的神，就是出于天主的。} 那么，异端者中间的神也是出于天主的，因为他们\Quote{明认耶稣基督是在肉身内降世}？是的，也许他们在这里起来反对我们，说：你们没有从天主来的圣神；但我们是明认\Quote{耶稣基督是在肉身内降世}的。但宗徒在这里说过，那些不明认\Quote{耶稣基督是在肉身内降世}的人没有圣神。问亚略派：他们明认\Quote{耶稣基督是在肉身内降世}；问欧诺米派；他们明认\Quote{耶稣基督是在肉身内降世}；问马其顿派；他们明认\Quote{耶稣基督是在肉身内降世}；问加塔弗利基派；他们明认\Quote{耶稣基督是在肉身内降世}；问诺瓦替安派；他们明认\Quote{耶稣基督是在肉身内降世}。那么，所有这些异端都有圣神吗？那么他们不是假先知吗？那里没有欺骗，没有诱惑吗？他们确实是假基督；因为\Quote{他们从我们中间出去，却不属于我们。}\par

13. 那么，我们该怎么办？凭着什么辨别他们？要十分留意；让我们同心前行，叩门。爱德亲自守望；因为不是别的，正是她叩门，也是她开门：你很快就会因我们的主耶稣基督的名明白。你们已经听见上面说过：\Quote{凡否认耶稣基督是在肉身内降来的，便是假基督。}在那里我们也曾问：谁否认？因为我们既不否认，那些人也不否认。我们发现有些人在行为上否认；我们从宗徒那里引证了证据，他说：\Quote{他们声称认识天主，但在行为上却否认祂。}\BibleRef{铎 1:16}这样，现在让我们也照样从行为上——而不是从口舌上——查究。什么神不是出于天主？就是那\Quote{否认耶稣基督是在肉身内降来的}神。什么神是出于天主？就是那\Quote{明认耶稣基督是在肉身内降来的}神。谁明认耶稣基督是在肉身内降来的？如今，弟兄们，注意目标！让我们看行为，不要停留在口舌的喧哗上。我们要问基督为什么在肉身内降来，这样我们就能找出那些否认祂在肉身内降来的人。如果你停留在口舌上，怎么，你会听到许多异端明认基督是在肉身内降来的；但真理定那些人的罪。基督为什么在肉身内降来？祂不是天主吗？经上不是论祂说：\Quote{在起初已有圣言，圣言与天主同在，圣言就是天主？}\BibleRef{若 1:1}不是祂喂养天使，不是祂喂养天使吗？祂不是这样来到这里，而未曾离开那里吗？祂不是这样升天，而未曾舍弃我们吗？那么祂为什么在肉身内降来？因为我们需要将复活的希望显示给我们。祂是天主，祂在肉身内降来；因为天主不能死，肉身却能死；祂于是在肉身内降来，好为我们而死。但祂怎样为我们而死？\Quote{人若为自己的朋友舍掉性命，再没有比这更大的爱德了。}\BibleRef{若 15:13}因此，是爱德领祂进入肉身。所以，凡没有爱德的，便是否认基督是在肉身内降来的。此时，你要质问一切异端者：基督是在肉身内降来的吗？\Quote{祂确实降来了；我相信，我明认。}不，你否认。\Quote{我怎样否认？你听见我这样说！}不，我定你否认的罪。你用声音说，用心否认；用言语说，用行为否认。\Quote{怎么，}你说，\Quote{我怎样在行为上否认？}因为基督在肉身内降来的目的，是要为我们而死。祂为我们而死，因为其中祂教导了极大的爱德。\Quote{人若为自己的朋友舍掉性命，再没有比这更大的爱德了。}你没有爱德，因为你为自己的荣誉分裂合一。因此，你们要以此辨别那出于天主的神。敲打那些瓦器，试验它们，看是否有了裂纹而发出沉闷的声音；看它们是否发出饱满清澈的响声，看爱德是否在那里。你把自己从整个大地的合一中抽离，你以分裂肢解教会，你撕裂基督的身体。祂在肉身内降来，为要聚集归一，你却喧嚷着要散播分离。所以，这就是天主的神，它说耶稣是在肉身内降来的，它说——不是用口舌，而是用行为；它说——不是制造喧哗，而是去爱。而那否认耶稣基督是在肉身内降来的神，不是出于天主；这里否认，同样不是用口舌，而是用生命；不是用言语，而是用行为。因此，我们如何认识弟兄们，已经显明了。许多人在内里，可以说是在内里；但没有人是在外面，除非他确实是在外面。\par

14. 不但如此，为使你们知道他把事情指涉到行为，他说：\Quote{并且一切神，\OriginalTerm{取消基督}{qui solvit Christum}，即取消基督是在肉身内降来的，不是出于天主。}这里指的是行为上的取消。他给你显示了什么？\Quote{那否认的}：在他说\Quote{取消}（或\Quote{解除}）时。祂来是要聚集归一，你却来取消。你要把基督的肢体撕碎。怎能说你不否认基督是在肉身内降来的呢？你撕裂了天主所聚集的教会？因此，你反对基督；你是假基督。无论你在内里，还是在外头，你都是假基督：只不过，你在内里时，你是隐藏的；你在外头时，你是显露的。你取消耶稣，并否认祂在肉身内降来；你不出于天主。因此，祂在福音中说：\Quote{谁若废除这些诫命中最小的一条，并这样教训人，在天国里将称为最小的。}\BibleRef{玛 5:19}这是什么废除？什么教训？是行为上的废除和好像言语上的教训。\Quote{你教导人不可偷盗，你自己却偷盗？}\BibleRef{罗 2:21}因此，偷盗的人在行为上废除或取消诫命，并好像这样教训人：\Quote{他在天国里将称为最小的，}即在此世的教会中。论到这种人，经上说：\Quote{他们所说的，你们要遵行；但他们所做的，你们不要做。}\BibleRef{玛 23:3} \Quote{但谁若实行，并这样教训人，他在天国里将称为大的。}由此，祂在这里说了\OriginalTerm{实行}{fecerit}，\Quote{实行}；而与此对比，祂在那里说了\OriginalTerm{取消}{solverit}，意思是\OriginalTerm{不实行}{non fecerit}，\Quote{不实行，却这样教训人}——那么，废除就是不去实行——祂教导我们什么，岂不是要我们查问人的行为，而不轻信他们的话语？事物的晦暗迫使我们说得冗长，主要的是，主愿启示的事，即使理解较慢的弟兄也能领会，因为众人都是基督的血所赎买的。我担心这书信本身不能在这些天内如我所许的完成；但按主的意愿，与其以过多的食物压垮你们的心，不如保留余下的部分。\par

\SourceRecord{Translated by H. Browne.；Nicene and Post-Nicene Fathers, First Series；Vol. 7.；Edited by Philip Schaff.；Buffalo, NY: Christian Literature Publishing Co.,；1888.；<http://www.newadvent.org/fathers/170206.htm>.}

% structure: p0001 source=h1 target=section reason=page-title

\section{论若望一书第七篇讲道}

% structure: p0002 source=h2 target=subsection reason=relative-heading-rank; base=h2

\subsection{若望一书 4:4-12}

\Quote{现在你们是出于天主的，小孩子们，你们已战胜了他，因为在你们内的那一位，比在世界上那一位更大。他们是出于世界的，因此他们谈论世界，世界也听从他们。我们是出于天主的，那认识天主的，听从我们；那不出于天主的，不听从我们。由此我们认识真理之神和谬误之神。可爱的诸位，让我们彼此相爱，因为爱是出于天主；凡爱人的，是由天主而生，并且认识天主。那不爱人的，不认识天主，因为天主是爱。天主对于我们的爱在这事上显示出来，就是天主派遣了祂的独生子到世界上来，使我们藉着祂而生活。爱就在于此，不是我们爱了天主，而是祂爱了我们，并且派遣了祂的儿子为我们赎罪。可爱的诸位，如果天主这样爱了我们，我们也应当彼此相爱。从来没有人见过天主。}\par

1. 对于所有寻求故乡的忠信者而言，这世界就像以色列人民在旷野中一般。他们确实还在漂泊，寻求他们的故乡；但有天主作为他们的向导，他们不会迷途。他们的道路就是天主的命令。因为他们四十年间所经过的行程，其实只有很少的站口，是众所周知的。他们被耽搁，是因为在接受训练，而不是被抛弃。因此，天主所应许我们的，是说不出的甘饴和美善，正如圣经所说，你们也常听我们讲过的：\Quote{眼所未见，耳所未闻，人心所未想到的。}（\BibleRef{依 64:4}，\BibleRef{格前 2:9}）我们藉着现世的劳苦被锻炼，藉着今生的试探被训练。因此，如果你不愿在这旷野中渴死，就请饮爱德吧。它是天主乐意放置在泉水，使我们在路上不致昏厥；当我们到达故乡时，我们将更丰盛地饮用它。福音书刚刚读过；现在就讲题结束的话，你们还听了别的什么，不就是关于爱德吗？因为我们在祈祷中与天主约定：如果我们愿意祂宽恕我们的罪，我们也应宽恕得罪我们的人\BibleRef{玛 6:12}。那宽恕的，不是别的，正是爱德。若从心中除去爱德，仇恨就占据它，它就不懂得宽恕。让爱德存在，它就毫无畏惧地宽恕，不受困扰。我们着手为你们讲解的这整封信，你们看，它所赞许的除了这一件事——爱德——还有别的吗？我们无需担心因多讲而令人生厌。如果爱德变得可憎，还有什么可爱呢？正是藉着爱德，其他事物才得以正当地被爱；那么，爱德本身必须如何被爱！因此，那永远不应从心中离去的，也不应从舌上离去。\par

2. \Quote{现在，}他说，\Quote{你们是出于天主的，小孩子们，你们已战胜了他：}\BibleRef{若一 4:4}战胜谁呢？岂不是那假基督吗？因为上面他说过：\Quote{凡否认耶稣基督，否认祂曾降生成人的，就不是出于天主。}如果你们记得，我们曾讲解说：所有违反爱德的人，就是否认耶稣基督曾降生成人。因为耶稣之所以必须来临，无非是为了爱德；正如主自己在福音中所称赞的爱德说：\Quote{人不能有比这更大的爱了，就是为朋友舍命。}\BibleRef{若 15:13}天主子若不穿上肉体以便死，又怎能为我们舍命呢？因此，凡违反爱德的，任他口里说什么，他的生活就否认基督曾降生成人；这就是假基督，无论他身在何处，无论他怎样来。但宗徒对那乡国的公民——我们为之叹息的——说什么呢？\Quote{你们已战胜了他。}他们是凭着什么战胜的呢？\Quote{因为那在你们内的，比那在世界上的更大。}为使他们不把胜利归于自己的力量，不因骄傲的傲慢而被战胜——因为魔鬼使谁骄傲，就战胜谁——宗徒愿意他们保持谦卑，就说：\Quote{你们已战胜了他。}现在，任何人听了这句话：\Quote{你们已战胜了，}就会抬起头，挺起颈项，希望被赞美。不要自高自大；请看那在你们内里战胜的是谁。你们为什么战胜？\Quote{因为那在你们内的，比那在世界上更大。}要谦卑，背负你们的主；做那让祂骑坐的牲畜。祂统治，祂引导，对你们是有益的。因为如果你们没有祂骑在你们身上，你们可能会昂首踢蹄；但没有统治者的自由是祸哉，因为这种自由使你们被野兽吞食！\par

3. \BibleQuote{这些是属于世界的。} \BibleRef{若一 4:5} 是谁？是那些假基督。你已经听过他们是谁。如果你不是这样的人，你认识他们；但无论谁是这样的人，却不认识。\BibleQuote{这些是属于世界的，因此他们谈论世界，世界也听从他们。} 哪些人是\Quote{谈论世界}的？注意那些反对爱德的人。看，你听主说过：\BibleQuote{如果你们宽免别人的过犯，你们的天父也必宽免你们的过犯；但如果你们不宽免别人的过犯，你们的天父也必不宽免你们的过犯。} \BibleRef{玛 6:14} 这是真理的断言；或若不是真理在说话，就反驳它。如果你是一个基督徒，并相信基督，祂说过：\Quote{我是真理。} 这个断言是真实的，是坚定的。现在请听那些\Quote{谈论世界}的人。\Quote{难道你不报复你自己吗？难道你让他说他这样对你做了吗？不：让他觉得他是在跟一个人打交道。} 每天都有这样的话被说。说这样话的人，\Quote{他们谈论世界，世界也听从他们。} 没有人说这样的话，只有那些爱世界的人；也没有人听从这样的话，只有那些爱世界的人。你已经听说，爱世界而忽略爱德，就是否认耶稣基督成了肉身。或者说，主自己却在肉身中这样做了吗？祂被掌掴时，愿意报复吗？祂被挂在十字架上时，难道没有说：\BibleQuote{父啊，宽恕他们，因为他们不知道他们所做的。} \BibleRef{路 23:34} 但如果祂——那拥有权能者——没有威胁；你为什么威吓？为什么你满腔愤怒？你是在别人的权下。祂死，是因为祂愿意死，然而祂没有威胁；你不知道你何时会死，而你还威胁吗？\par

4. \BibleQuote{我们是属于天主的。} \BibleRef{若一 4:6} 让我们看看为什么；看看是否是为了爱德以外的任何事物。\BibleQuote{我们是属于天主的：认识天主的，听从我们；不属于天主的，不听从我们。由此我们可以认出真理之神和谬误之神。} 即借此：那听从我们的，有真理之神；那不听从我们的，有谬误之神。让我们看看他劝告什么，并让我们选择在真理之神中听从他的劝告，而不是假基督、爱世界的人、世界。如果我们是由天主所生的，\BibleQuote{亲爱的，} \BibleRef{若一 4:7} 他继续说——看上文从何处：\BibleQuote{我们是属于天主的：认识天主的，听从我们；不属于天主的，不听从我们。由此我们可以认出真理之神和谬误之神。} 是的，现在他使我们热切地留意：被告知那认识天主的就听从，但不认识的就不听从；并且这就是分辨真理之神和谬误之神的准则：那么，让我们看看他将要劝告什么；在什么上我们必须听他——\BibleQuote{亲爱的，让我们彼此相爱。} \BibleRef{若一 4:7} 为什么？因为一个人劝告？\BibleQuote{因为爱德出于天主。} 他大大赞扬了爱德，因他说，\Quote{出于天主：} 但他说更多的；让我们热切地听。目前他说了：\BibleQuote{爱德出于天主；凡有爱德的，都是由天主所生，并且认识天主。那不爱的，不认识天主。} \BibleRef{若一 4:7} 为什么？\BibleQuote{因为天主是爱} [爱就是天主]。还能说什么更多呢，弟兄们？如果在这封书信的篇幅中，或是在圣经的其他所有篇幅中，没有任何赞扬爱德的言语，而只有这一件事是圣神的声音告诉我们的：\Quote{因为爱就是天主；} 我们就不应再要求更多了。\par

5. 现在看，反对爱德就是反对天主。不要说：\Quote{当我不爱我的弟兄时，我得罪的是人（注意！），得罪人是容易承受的事；只是不要让我得罪天主。} 当你得罪爱德时，你怎能不得罪天主呢？\Quote{爱就是天主。} 是\Quote{我们}说这话吗？如果我们说了，\Quote{爱就是天主，} 恐怕你们中有人会被冒犯而说：他说了什么？他是什么意思，\Quote{爱就是天主}？天主\Quote{赐给了}爱，作为恩赐，天主赋予了爱。\Quote{爱德出于天主：爱就是天主。} 看，弟兄们，你们有天主的圣经：这封书信是正典；在万国中被诵读，被全地之权威所持守，它建立了整个大地。你们在此被天主圣神告知：\Quote{爱就是天主。} 现在，如果你敢，就反对天主，并拒绝爱你的弟兄！\par

6. 那么，刚才说\Quote{爱出自天主}；现在又说\Quote{爱就是天主}？这是什么意思呢？因为天主是圣父、圣子、圣神：圣子是出自天主的圣子，圣神是出自天主的圣神；这三者是一个天主，不是三个天主。如果圣子是天主，圣神是天主，而且那拥有圣神的人是在爱中，因此\Quote{爱就是天主}；但\Quote{就是天主}，是因为\Quote{出自天主}。因为你们在书信中两者都见到；既有\Quote{爱出自天主}，又有\Quote{爱就是天主}。关于圣父，圣经从不说祂是\Quote{出自天主}；但当你听到\Quote{出自天主}这个说法时，要么指圣子，要么指圣神。因为宗徒说：\BibleQuote{天主的爱，藉着所赐与我们的圣神，已倾注在我们心中。}\BibleRef{罗 5:5}我们应明白那存在于爱中的，是圣神。因为正是这位圣神，恶人不能领受，正是祂是那泉源，关于它圣经说：\BibleQuote{你水泉的活水当属你独有，不可与外人共享。}\BibleRef{箴 5:16}因为凡不爱天主的，都是外人，都是假基督。虽然他们来到教会，却不能被列入天主的儿女中；那生命的泉源不属于他们。拥有圣洗甚至可能属于恶人；拥有先知之恩甚至可能属于恶人。我们发现国王\Person{撒乌耳}拥有先知之恩：他正在迫害圣\Person{达味}，却被先知之神充满，开始说预言。\EditorialNote{撒慕尔纪上 19}领受主身体和血的圣事甚至可能属于恶人：因为关于这种人经上说：\BibleQuote{那不相称地吃和喝的人，就是吃喝自己的定案。}\BibleRef{格前 11:29}拥有基督的名号甚至可能属于恶人；\Emph{i.e.}，恶人也可以被称为基督徒，正如那些经上所说的人：\BibleQuote{他们玷污了自己的天主之名。}\BibleRef{则 36:20}我说，拥有一切这些圣事甚至可能属于恶人；但拥有爱德而同时是个恶人，是不可能的。那么，这就是那特殊的恩赐，这就是那唯独属于一个人的\Quote{泉源}。圣神劝诫你饮此泉，圣神劝诫你饮祂自己。\par

7. \BibleQuote{天主对我们的爱在这事上已显示出来：}\BibleRef{若一 4:9}请看，为了使我们能爱天主，我们有劝诫。除非祂先爱了我们，我们能爱祂吗？如果我们爱得迟缓，我们不要迟缓于回报爱。祂先爱了我们；即使这样，我们仍不爱。祂爱了不义的人，但祂除去了不义；祂爱了不义的人，但不是为了不义而聚集他们；祂爱了患病的人，但祂探望他们是为了使他们痊愈。所以，\Quote{爱就是天主。}\BibleQuote{天主对我们的爱在这事上已显示出来，就是天主把自己的独生子打发到世界上来，好使我们藉着祂得到生命。}\BibleRef{若一 4:9}主亲自说：\BibleQuote{人若为自己的朋友舍掉性命，再没有比这更大的爱情了。}\BibleRef{若 15:13}基督对我们的爱由此证明：祂为我们死了；圣父对我们的爱如何证明？在于祂\Quote{派遣了自己的独生子}为我们而死。宗徒\Person{保禄}也说：\BibleQuote{祂没有怜惜自己的儿子，反而为我们众人把祂交出来，岂不也把一切与祂一同赐给我们吗？}\BibleRef{罗 8:32}看，圣父交出了基督；\Person{犹达斯}也交出了祂；这不像是同样的事吗？\Person{犹达斯}是\Quote{ \Emph{traditor} }，即交出者（或叛徒）：圣父也是如此吗？决不可以！你会说。我不这样说，但宗徒说：祂没有怜惜自己的儿子，反而\Quote{ \Emph{tradidit Eum} }为我们众人把祂交出来。圣父交出了祂，圣子也交出了自己。同一位宗徒说：\BibleQuote{祂爱我，为我舍弃了自己。}\BibleRef{迦 2:20}如果圣父交出了圣子，圣子交出了自己，\Person{犹达斯}做了什么？有出自圣父的\Quote{ \Emph{traditio} }（交出）；有出自圣子的\Quote{ \Emph{traditio} }；有出自\Person{犹达斯}的\Quote{ \Emph{traditio} }。所做的事相同，但如何区别圣父交出圣子、圣子交出自己、和门徒\Person{犹达斯}交出他的老师？在于此：圣父和圣子是出于爱而做，\Person{犹达斯}却是出于背叛的出卖。你们看，不是人做什么事要予以考虑，而是他以怎样的心意和意愿去做。我们在同一行为中找到圣父和\Person{犹达斯}；我们祝福圣父，憎恶\Person{犹达斯}。为什么祝福圣父而憎恶\Person{犹达斯}？我们祝福爱德，憎恶不义。基督的被交出给人类带来了多么大的益处！\Person{犹达斯}是否想到这一点而交出祂？天主想到的是我们的得救，藉此我们被赎回；\Person{犹达斯}想到的是他出卖主所得的价钱。圣子自己想到的是祂为我们所付的代价，\Person{犹达斯}想到的是他收到出卖祂的价钱。因此，不同的意向使所做的事不同。虽然事情是一样的，但若以不同的意向来衡量，我们会发现一件值得爱，另一件应受谴责；一件值得赞扬，另一件应受憎恶。爱德的力量就是这样。看，唯独爱德能辨别，唯独爱德能区分人的行为。\par

8. 这事在行为相似的情况下我们已经说过。在行为不同的情形中，我们看到一个人因爱德而变得严厉，又因不义而变得谄媚温柔。父亲责打孩子，而拐骗者却抚爱孩子。如果你比较这两件事，责打与抚爱，谁不会选择抚爱而避免责打呢？如果你看行为者，责打的是爱德，抚爱的是不义。看我们坚持的：人的行为只能由爱德之根来辨别。因为许多事外表看起来很好，却并非出于爱德之根。荆棘也会开花；有些行为看起来粗鲁、野蛮，却是出于爱德的训导而为纪律所做。简而言之，给你一条简短诫命：\Quote{爱，并做你所愿：}你若缄默，因爱而缄默；你若呼喊，因爱而呼喊；你若纠正，因爱而纠正；你若宽恕，因爱而宽恕：让爱德之根在内，从这根上只能长出善来。\par

9. \Quote{在这一点上，爱——天主对我们的爱显明于此，因为天主差遣了祂的独生子来到这世界，使我们能藉着祂而生活。——在这一点上是爱，不是我们爱了天主，而是祂爱了我们：} \BibleRef{若一 4:9} 我们并没有先爱祂：祂爱我们正是为此，好叫我们爱祂；并且差遣祂的儿子为我们的罪作赎罪祭：\OriginalTerm{祭献者}{litatorem}，即一位献祭者。祂为我们的罪献上了祭品。祂在哪里找到这祭品？祂在哪里找到祂要献上的纯洁牺牲？祂没有找到别的，祂把自己献上了。\Quote{亲爱的，如果天主这样爱了我们，我们也应当彼此相爱。\BibleRef{若一 4:11}伯多禄，}祂说，\Quote{你爱我吗？}他回答说，\Quote{我爱。}\Quote{你喂养我的羊。}\par

10. \Quote{从来没有人见过天主：} \BibleRef{若一 4:12} 祂是不可见的；不该用眼睛，而该用心去寻求祂。正如我们若想看见太阳，就应洁净肉身的眼睛；想看见天主，就应洁净那能看见天主的眼睛。这眼睛在哪里？听福音说：\Quote{心里洁净的人是有福的，因为他们要看见天主。} \BibleRef{玛 5:8} 但不要让任何人按自己肉眼的欲望来想象天主。因为这样他会为自己构想一个巨大的形体，或某种不可测度的庞大之物，如同他用肉眼所见的光，向各方延伸；他尽其所能地赋予它一片又一片的空间；或者，他想象一位形貌可敬的老人。你不可想象这些事。若你想看见天主，有一样东西你可以想象：\Quote{天主是爱。} 爱有什么面貌？有什么形状？有什么身高？有什么脚？有什么手？没有人能说。然而它有脚，因为这些脚把人带到教堂；它有手，因为这些手伸向穷人；它有眼睛，因此我们才能顾念穷人：经上记载：\Quote{顾念穷乏和可怜人的人是有福的。} 它有耳朵，关于这耳朵，主说：\Quote{有耳可听的，就听吧！} \BibleRef{路 8:8} 这些肢体并不是按部位分开的，而是藉着理解，凡有爱德的人立刻看到整体。你居住，你就会被居住；你居住，你就会被居住。因为我的弟兄们，你们怎么说？谁爱那看不见的事物呢？如今当爱德受到赞美时，你们为何举手、欢呼、赞美？我给你们看了什么？我展示了颜色的光彩吗？我提出了金银吗？我从隐藏的宝库中挖出了宝石吗？我向你们的眼睛显示了这类东西吗？我说话时我的面容改变了吗？我处于肉身中；我以同样的形貌来到你们面前；你们以同样的形貌来到这里；爱德受到赞美，你们欢呼。显然你们什么也没有看见。然而正如你们赞美时它使你们喜悦，同样愿它使你们喜悦，好叫你们能在心中持守它。因为请留心听我说，弟兄们；我激励你们所有人，按天主所赐的能力，朝向一个伟大的宝藏。如果给你们看一只美丽的小花瓶，浮雕的、镶金的、精工细作的，它迷住了你们的眼睛，吸引你们内心的渴望，你们因工匠的手艺、银子的重量、金属的光彩而喜悦；难道你们每个人不会说：\Quote{哦，我要是有那只花瓶就好了！} 但你们这样说也是徒然，因为你们是否拥有它并不由你们决定。或者，如果有人想要它，他可能会想从别人的房子里偷来。爱德被赞美给你们；如果它使你们喜悦，你们就拥有它，占有它：无需抢夺任何人，无需想购买它；它是白白赐予的，无需代价。拿去，拥抱它；没有比它更甜蜜的了。如果它只是被谈论时就如此，那么拥有它时又该是怎样呢？\par

11. 如果你们中有人希望持守爱德，弟兄们，首先切勿以为它是一种卑微懒散的东西；也不要以为爱德是靠某种温和——不，不是温和，而是驯服和懈怠——来保存的。绝非如此。不要以为当你没有鞭打你的仆人时你便爱他，或者当你没有管教你的儿子时你便爱他，或者当你没有斥责你的邻人时你便爱他：这不是爱德，而仅仅是软弱。愿爱德热忱地去纠正、去改善：但若有良好的品行，就让它们令你喜悦；若有不良的，就让它们被改善、被纠正。不在人身上爱他的错误，而要爱人：因为人是天主所造的，错误是人自己所造的。要爱天主所造的，不要爱人自己所造的。当你爱那一个时，你就除去了这一个；当你重视那一个时，你就改善了这一个。但即便你有时严厉，也要出于爱德，为了纠正。为此，爱德以那降临在主身上的鸽子作为预象。那鸽子的形象，是圣神降临的形象，圣神要将爱德倾注在我们心中：为何如此？鸽子没有胆汁，但用喙和翅膀为它的幼雏争战；它是一种没有苦涩的凶猛。父亲也是如此：当他责罚儿子时，是为了管教而责罚他。如我所说，拐子为了贩卖，用痛苦的甜言蜜语哄骗孩子；父亲为了纠正，却没有胆汁地责罚。你们对所有人也当如此。请看，弟兄们，一个伟大的教训，一条伟大的规则：你们每人都有孩子，或希望有孩子；或者他若完全决定没有肉身的孩子，至少属灵地渴望有孩子——哪一个父亲不纠正他的儿子？哪一个儿子不受父亲的管教？然而他看起来对儿子是严厉的。这是爱的严厉，爱德的严厉：一种像鸽子一样没有胆汁的严厉，而不是像乌鸦那样的严厉。因此，我的弟兄们，我想到要告诉你们，那些破坏爱德的人就是制造分裂的人：他们既憎恨爱德本身，也憎恨鸽子。但鸽子斥责他们：它从天而降，天开了，它停在主的头上。这是为什么呢？为叫若望听见：\Quote{这就是那施洗者。} \BibleRef{若 1:33} 走开，你们这些强盗；走开，你们这些侵犯基督产业的人！在你们自己的产业上，你们硬要作主，竟敢刻上那大主人的名号。祂认出自己的名号；祂为自己的产业辩护。祂不涂销名号，而是进入并占据产业。因此，一个人来到天主教会时，他的圣洗并不被涂销，以免统帅的名号被涂销：但在天主教会中做了什么？名号被承认；主人凭着自己的名号进入，而那强盗却是凭着他人的名号进入的。\par

\SourceRecord{Translated by H. Browne.；Nicene and Post-Nicene Fathers, First Series；Vol. 7.；Edited by Philip Schaff.；Buffalo, NY: Christian Literature Publishing Co.,；1888.；<http://www.newadvent.org/fathers/170207.htm>.}

% structure: p0001 source=h1 target=section reason=page-title

\section{论若望一书第八篇讲道}

% structure: p0002 source=h2 target=subsection reason=relative-heading-rank; base=h2

\subsection{\BibleRef{若一 4:12-16}}

\BibleQuote{如果我们彼此相爱，天主就住在我们内，祂的爱在我们内得以成全。这样我们就知道我们住在祂内，祂也住在我们内，因为祂将自己的圣神赐给了我们。我们曾看见，并且作证：圣父派遣了圣子为世界的救主。凡明认耶稣是天主圣子的，天主就住在他内，他也住在天主内。我们也认识并相信了天主对我们所有的爱。天主是爱；那住在爱内的，就住在天主内，天主也住在他内。}\par

1. 爱是一个甘饴的词，但行为更甘饴。我们不能总是谈论它，因为我们有许多事要做，各种事务将我们引向不同方向，以致我们的舌头没有闲暇时刻谈论爱；其实，我们的舌头也没有比这更好的事可做。然而，尽管我们不能总是谈论爱，却可以时常保有它。正如我们此刻所唱的\Emph{阿肋路亚}，我们能总是在做这件事吗？不是每一小时——我甚至不是说整整一小时——我们唱阿肋路亚，而是仅在一小时中短暂的片刻，然后便转向其他事情。\Emph{阿肋路亚}，正如你们已经知道的，意为\Quote{赞美主}。那用舌头赞美主的人，不能总是这样做；那以生活与行为赞美主的人，却能时刻如此。慈悲善工、爱德之情、虔诚圣洁、贞洁无瑕、节制谦逊，这些德行总是要实践的：无论我们在公开场合或在家中，在众人前或在内室，说话或沉默，忙于事务或闲散，这些德行总是要持守的，因为我所列举的一切美德都在内心。但谁能够全部列出它们呢？这些美德有如皇帝的军队，驻扎在你的心神内。因为如同皇帝借军队行其所愿，主耶稣基督一旦开始住在我们的内心（即借着信德住在心神内），便使用这些美德作为祂的臣仆。并且，借着这些不能用肉眼看见、然而当被提及便受到赞美——它们若非被爱便不会受赞美，若非被看见便不会被爱；若非被爱且被看见，它们便是以另一种眼睛，即内心的内在注视被看见——借着这些看不见的美德，肢体被可见地带动：脚行走，但走向哪里？走向那善的意志所移动之处，善的意志犹如服侍善君的士兵；手工作，但做什么？做那由圣神在内所感发的爱德所命令的事。因此，肢体被带动时是可见的；那在内命令它们的，却是不可见的；而那在内命令的是谁，几乎只有命令者以及那在内被命令者才知道。\par

2. 弟兄们，你们刚才听到福音被宣读时，如果你们不仅用肉体的耳朵，也用心灵的耳朵去听的话。它说了什么？\BibleQuote{你们应当心，不要在人前行你们的义德，为叫他们看见。}\BibleRef{玛 6:1} 祂的意思是，无论我们做什么善事，都要从人的眼前隐藏起来，害怕被人看见吗？如果你们害怕观众，你们就不会有模仿者；因此你们应当被看见。但你们不可为了被看见而行善。你们的喜乐不应以此为终点，你们的欢欣不应以此为归宿，以致你们认为当你们被看见、被称赞时，就获得了你们善工的全部果实。这算不了什么。当你们被称赞时，要轻视自己，让那藉着你们行事的那一位在你们内受称赞。因此，不要为了自己的称赞而行善，而要为那赐给你们行善能力的那一位的称赞。从你们自己你们有恶行，从天主你们有善行。另一方面，看看那些乖戾的人，他们是多么颠倒。他们行的善，总要归功于自己；若行了恶，总要归咎于天主。你们要扭转这种扭曲和颠倒的作法，它可以说把事物头足倒置，使在上者变成在下者，使在下者变成在上者。你们想让天主成为在下者，而你们成为在上者吗？你们是跌倒，而不是升高；因为祂总是在上。那么，你们行善，而天主行恶？不，你们若说得更真实，应当说：我作恶，祂行善；我所行的善是来自祂的善行，因为从我自己，我所行的尽是恶。这种宣认坚固人心，并打下爱德的坚实基础。因为如果我们应当隐藏善工，免得被人看见，那么主在山中圣训中所说的那句话又怎么办呢？祂在那里说了这话，也在稍前说了：\BibleQuote{照样，你们的光也当在人前照耀，好使他们看见你们的善行。}\BibleRef{玛 5:16} 祂并没有停在那里，没有就此结束，而是加上：\BibleQuote{并光荣你们在天之父。} 宗徒怎么说？\BibleQuote{我本人与犹太在基督内的各教会是互不相识的；他们只是听说：那从前迫害我们的，如今却传扬他曾经一度摧毁的信仰。他们就因我而光荣了天主。}\BibleRef{迦 1:22} 请看他是怎样，尽管他变得如此广为人知，却并未将善事归于自己的赞美，而是归于天主的赞美。至于他自己，他曾是摧毁教会的人，迫害者，嫉妒者，恶毒者，是他自己承认了这些，并不是我们以此责备他。保禄喜欢我们谈论他的罪过，好使那治愈如此大疾病的祂受光荣。因为正是医生的手割断了并治愈了那大疮口。那从天上来的声音打倒了迫害者，兴起了宣讲者；杀死了\Person{扫禄}，使\Person{保禄}活过来。因为\Person{扫禄}曾是一位圣者的迫害者；这人由此而得名，当他迫害基督徒时：\EditorialNote{撒慕尔纪上第十九章} 此后扫禄成了保禄。\Emph{\Person{保禄}}这个名字是什么意思？是小的。因此当他是扫禄时，他是骄傲的、自高自大的；当他是保禄时，他是谦卑的、渺小的。因此我们说：我要见你\Quote{\Emph{paulo post}}，\Emph{即}稍后。请听他是怎样成为渺小的：\Quote{我是宗徒中最小的；诸圣中最小的}，他在别处说。因此他在宗徒中如同衣袍的边穗，但外邦人的教会触摸了它，如同那患血漏的妇女，便痊愈了。\BibleRef{玛 9:20}\par

3. 那么，弟兄们，这是我要说的，这是我现在说的，如果可能，我不会留下不说：愿你们内现在有这些工作，现在有那些，按照时间，按照时辰，按照日子。你们要常常说话吗？常常沉默吗？常常休养身体吗？常常禁食吗？常常给穷人面包吗？常常给赤身者衣服吗？常常探访病人吗？常常使不和者重归于好吗？常常埋葬死者吗？不，而是现在这样，现在那样。这些事被着手进行，然后停止；但那像皇帝一样统率着内在一切力量的，既无开端也不应停止。愿内在的爱德没有间断：愿爱德的任务按照时间展现出来。愿那\BibleQuote{兄弟之爱}，正如经上所记载的，愿\BibleQuote{兄弟之爱常存。}\BibleRef{希 13:1}\par

4. 但或许你们当中有些人一直感到奇怪，当我们向你们阐释这位有福的若望的这封书信时，他为何如此强调地称赞\Quote{兄弟}之爱。他说：\Quote{爱自己兄弟的}，又说：\Quote{我们领受了命令，要彼此相爱。}他屡次提及兄弟之爱，但对天主的爱——即我们应当爱天主的爱——却没有如此频繁地提到；然而，他也并非完全闭口不谈。但关于爱仇敌，几乎整封书信他都没有提及。尽管他大力宣扬并推崇爱德，但他并没有吩咐我们要爱仇敌，而是吩咐我们要爱弟兄。然而，刚才宣读福音时，我们听到了：\Quote{因为你们若爱那爱你们的，有什么赏报呢？税吏不也是这样做吗？}\BibleRef{玛 5:46}那么，若望宗徒为何把兄弟之爱当作关乎我们达至某种完美的大事来称赞，而主却说，我们爱弟兄还不够，还应当将爱扩展到甚至爱仇敌呢？那扩展到仇敌的人，并没有越过弟兄。爱必如火焰，先抓住最近之物，然后延伸至更远之物。弟兄比你任何陌生人都更亲近。再者，你不认识的、并非反对你的人，比反对你的仇敌更贴近你。将你的爱延伸到那些近人，但不要称之为\Quote{延伸}，因为爱那些与你亲近的人，几乎就是爱你自己。将爱延伸到陌生人，那些未曾伤害你的人。甚至越过他们，达到爱你的仇敌。至少这是主所命令的。为何宗徒在这里没有提及爱仇敌呢？\par

5. 所有的爱，无论是那被称为属血肉的爱——通常不被称为\OriginalTerm{喜爱}{dilectio}，而是\OriginalTerm{爱恋}{amor}（因为\Quote{喜爱}一词通常用于更好的对象，且被理解为指更好的对象）——然而，亲爱的弟兄们，所有的爱都包含着对所受之人的善意。因为我们不应该这样去爱，也无法这样去爱（无论是\OriginalTerm{喜爱}{diligere}还是\OriginalTerm{爱恋}{amare}；后一个词是主在说\Quote{\OriginalTerm{伯多禄，你爱我吗？}{Petra, amas me?}}时所用的），我们不应当像我们所听到的饕餮之人说\Quote{我爱画眉}那样去爱人。你问他为何爱它们？为了杀死，为了吞食。他说他爱，而这样爱它们是为了使它们不再存在；这样爱它们是为了消灭它们。我们爱任何食物，也是为此目的：为了被消耗并补充我们。难道人也应当如此被爱，以致被消灭吗？但有一种良好祝愿的友善，我们希望在某时某种场合对我们所爱的人行善。如果没有什么善可行呢？善意、良好祝愿本身便足以满足那爱者。因为我们不应当希望人们不幸，以便我们能够行慈悲的工作。你给饥饿的人面包，但更好的是没有人饥饿，而你无需施与。你给赤身者衣服，但愿人人都穿衣服，这种需要不再存在。你埋葬死者，但愿那不再有死亡的永生最终来临。你调解争吵，但愿耶路撒冷永恒的平安最终在此，那里无人不和！因为所有这些事都是面对需要的服务。除去不幸的人，慈悲的工作就会终止。慈悲的工作终止了，爱德的热忱会熄灭吗？以更真实的爱，你爱那幸福的人，他没有任何善行需要你去做；这爱将更纯洁，远更无掺杂。因为如果你对不幸的人行了善，你或许希望自己在他之上，并且希望他因你向他行了善而服从于你。他有需要，你施予了；你似乎觉得自己比接受者更伟大。愿他与你平等，这样你们两人都在同一位主之下，而主什么都不缺。\par

6. 因为骄傲的灵魂在此越过了界限，而且，在某种意义上，变得贪婪。因为经上说：\BibleQuote{贪财是万恶的根源}\BibleRef{弟前 6:10}；又说：\BibleQuote{一切罪恶的开始是骄傲}\BibleRef{德 10:15}。我们或许会问，这两句话如何一致：\BibleQuote{贪财是万恶的根源}和\BibleQuote{一切罪恶的开始是骄傲}？如果骄傲是一切罪恶的开始，那么骄傲就是万恶的根源。然而确实，\BibleQuote{贪财是万恶的根源}。我们发现在骄傲中也有贪婪（或抓取）；因为人越过了界限：什么是贪婪？就是超过所需。亚当因骄傲而堕落：\BibleQuote{一切罪恶的开始是骄傲}——圣经说：他是因抓取而堕落吗？还有什么比连天主都不能使他满足的人更贪婪的呢？事实上，我的弟兄们，我们读到人如何按天主的肖像和模样受造：天主论及他说：\BibleQuote{让他统治海中的鱼，天空的飞鸟，以及地上所有爬行的生物}\BibleRef{创 1:26}。祂说，统治人吗？祂说\BibleQuote{统治}：祂赐给了他自然的权力：\BibleQuote{统治}什么？\BibleQuote{海中的鱼，天空的飞鸟，以及地上所有爬行的生物}。为什么对这些事物的统治是自然的权力？因为人从这一点得到权力：他是按天主的肖像受造的。他按天主的肖像受造是在哪方面？在理智、在心思、在内在的人方面；他理解真理，辨别是非，知道是谁造了他，能够理解他的造物主，赞美他的造物主：谁有明智，谁就有这理智。因此，当许多人因邪恶的情欲在自己内磨损了天主的肖像，因行为的乖戾扑灭了理智的火焰，圣经向他们大声呼喊说：\BibleQuote{你们不要像骡马，没有理智}。这就是说，我立你们超越骡马；你们，我按我的肖像造了你们，赐给你们权柄管理这些。为什么？因为它们没有理性的心智；但你们藉着理性的心智能够认识真理，理解超越你们的事物：当服从那在你们之上的，那么你们所被设立管理的事物将在你们之下。但因着罪，人离弃了他应该服从的那一位，他就被置于他本应超越的事物之下。\par

7. 请注意我所说的：天主、人、野兽：即，在你们之上的是天主；在你们之下的是野兽。承认那在你们之上的天主，好使那些在你们之下的野兽也承认你们。\BibleRef{达 6:22}因此，因为达内尔承认天主在他之上，狮子也承认他高于它们。但如果你不承认那在你之上的，你就轻视你的上级，你就成为你下级的下级。因此，埃及人的骄傲是如何被挫败的？藉着青蛙和苍蝇。\EditorialNote{出谷纪第八章}天主本可以派遣狮子：但是一个伟大的人可能被狮子吓到。他们越骄傲，就越藉着可鄙微弱之物的手段粉碎他们邪恶的颈项。但达内尔，狮子承认他，因为他服从天主。那些被投入猛兽场与野兽搏斗、被野兽牙齿撕裂的殉道者，难道他们不服从天主吗？或者那三个人是天主的仆人，而玛加伯兄弟不是天主的仆人吗？火承认那三个人是天主的仆人，火烧不着他们，也不伤他们的衣服；\BibleRef{达 3:50}难道它不承认玛加伯兄弟吗？\EditorialNote{玛加伯下第七章}它承认玛加伯兄弟；是的，我的弟兄们，它也承认他们。但需要鞭打，蒙主的许可：祂在圣经中说：\BibleQuote{祂鞭打祂所接纳的每一个儿子}\BibleRef{希 12:6}。你们想，我的弟兄们，铁器能刺入主的身体除非祂允许，或者祂会被钉在树上除非这是祂的旨意？难道祂的受造物不认识祂吗？或者祂为祂的信徒立下了忍耐的榜样？你们看，天主有些明显地拯救了，有些没有明显地拯救：然而祂在精神上都拯救了，在精神上没有一个被遗弃。明显地，祂似乎遗弃了一些，似乎拯救了一些。因此祂拯救了一些，免得你们以为祂没有拯救的能力。祂证实了祂有能力，以致于在祂没有拯救的地方，你们可以明白一个更隐秘的旨意，而不是猜测祂无能为力。但是，弟兄们，当我们脱离所有必死的陷阱，当诱惑的时期过去，当这世界的河流消逝，我们将再次领受那\Quote{首袍}，即那因犯罪而失去的不朽，\BibleQuote{当这可朽坏的穿上不可朽坏的}，也就是说，这肉体穿上不可朽坏的，\BibleQuote{这可死的穿上不朽}\BibleRef{格前 15:44}，那时，已成为完美天主子女的人，不再需要被试探，也不再需要被鞭打，所有的受造物都要承认他们：万物都将服从我们，如果我们在这里服从了天主。\par

8. 因此基督徒应当如此，不夸耀胜过别人\Quote{人}。因为天主赐给你胜过野兽，即比野兽更优越。这是你生来就有的；你将永远比野兽更好。若你希望比别人更好，当你看见他与你平等时，你就会嫉妒他。你应当希望所有人都是你的平等者；若你在智慧上超过某人，你应当希望他也有智慧。当他迟钝时，他向你学习；当他未受教时，他需要你；你被视为老师，他是学生；因此你似乎是优越的，因为你是老师；他是低下的，因为他是学生。除非你希望他与你平等，否则你就是希望他永远是学生。但若你希望他永远是学生，你就会成为一个嫉妒的老师。若是一个嫉妒的老师，你又怎能是老师呢？我恳求你，不要把你的嫉妒教给他。请听宗徒论及爱德的怀抱所说的话：\BibleQuote{我切愿众人都像我一样。}\BibleRef{格前 7:7} 他在什么意义上愿众人与他平等？正是在这一点上他超越众人，即因着爱德他愿众人与他平等。我再说，人已经越过了界限；他贪求超过应得之分，想要超越人，而他本是被造为超越野兽的：这就是骄傲。\par

9. 请看骄傲做出了何等伟大的事。请你们牢记在心：骄傲的行为与爱德的行为是多么相似，多么不分轩轾。爱德喂养饥饿者，骄傲也如此做；爱德喂养，是为使天主受赞美；骄傲喂养，是为使自己受赞美。爱德给赤身者衣服，骄傲也如此做；爱德守斋，骄傲也如此做；爱德埋葬死者，骄傲也如此做。一切爱德愿意并实行的善工，骄傲也向着同样的目标努力，并且，可以说，勒紧她的马匹使之齐头并进。但爱德介于两者之间，不给不良驾驶的骄傲留有余地；不是驾驶不良，而是被不良驾驶。祸哉，那以骄傲为车夫的人，因为他必会头朝下坠落！然而在行善时，是不是骄傲驱使我们，谁知道呢？谁看见呢？它在哪里？我们看见的是行为：仁慈喂养，骄傲也喂养；仁慈接待旅客，骄傲也接待；仁慈为穷人代祷，骄傲也代祷。这是怎么回事？在行为上我们看不出差别。我敢说一点，但不是我说的，保禄已经说过：爱德致死，即一个有爱德的人宣认基督的名，忍受殉道；骄傲也宣认，也忍受殉道。一个有爱德，另一个没有爱德。但没有爱德的人，且听宗徒所说：\BibleQuote{我若把我所有的财产全施舍给人，我若舍身投火被焚，但我没有爱德，为我毫无益处。}\BibleRef{格前 13:3} 因此，神圣的圣经召唤我们从外在的面貌转向内在的实质；从这在人前夸耀的表面，召唤我们转向内在。回到你自己的良心，审问它。不要考虑外面绽放的是什么，而要看地里有什么根。那里是否扎根了情欲？可能有善行的外表，但不可能有真正的善工。那里是否扎根了爱德？不要害怕：从爱德不可能生出恶来。骄傲者抚慰，爱却是严厉的。一个给人衣服，另一个击打人。因为一个给人衣服是为了取悦人，另一个击打是为了以纪律矫正。爱德的击打比骄傲的施舍更被悦纳。所以，弟兄们，请进入内心，在一切事上，无论你们做什么，都要仰望你们的天主为见证。看，祂是否看见你们以何种心态行事。如果你们的心不控告你们是为炫耀而做，那就很好：不要害怕。但当你们行善时，不要怕被别人看见。要怕的是你们为之而行的目的是为了受赞美：让别人看见它，好使天主受赞美。因为如果你把它隐藏在人的眼前，你就剥夺了人的效法，你也夺去了天主应得的赞美。有两个对象接受你的施舍：两个饥饿者；一个饥饿的是面包，另一个饥饿的是义德。在这两个饥饿的灵魂之间——正如经上所载：\BibleQuote{饥渴慕义的人是有福的，因为他们要得饱饫。}\BibleRef{玛 5:6} ——在这两个饥饿者之间，你行善工的人被安置；若爱德因着其中一人的需要而行事，它就同时怜悯了两人，它要扶助两人。因为一个渴求可以吃的东西，另一个渴求可以效法的榜样。你喂养了一个，把自己作为榜样给了另一个；这样你就施舍给了两人：你使一个因消除饥饿而感谢你，使另一个因你树立的榜样而效法你。\par

10. 那么，要显示慈悲，如同有慈心的人；因为即使在爱仇人时，你们也是在爱弟兄。不要以为若望没有给予爱仇人的诫命，因为他没有停止谈论兄弟之爱。你们爱弟兄。\Quote{怎样，}你说，\Quote{我们爱弟兄？}我问你为何爱仇人。你为何爱他？为使他在今生平安？倘若对他并不有益呢？为使他有财富？倘若财富反使他盲目呢？为使他要妻？倘若他的婚姻生活痛苦呢？为使他有儿女？倘若儿女不成器呢？因此，你所似乎希望于仇人的这些事，在你爱他时，都是不确定的；它们是不确定的。你应希望他与你同得永生；你应希望他成为你的弟兄：当你爱他时，你爱的是弟兄。因为你爱的不是他现在的样子，而是你希望他成为的样子。我从前曾对你们说过，亲爱的弟兄们，若我没记错：有一块木料摆在眼前；一位良匠看见了这木料，尚未刨光，正如从森林中砍下时的样子，他喜欢上了它，想用它造出什么来。的确，他爱它，并非为让它始终如此。在他的技艺中，他看见了它将成为什么，而非在喜好中看见它是什么；他喜欢的是他将用它造出的东西，而非它本身。同样，天主爱了我们这些罪人。我们说天主爱了罪人：因为祂说：\BibleQuote{健康的人不需要医生，有病的人才需要。}\BibleRef{玛 9:12}祂爱我们这些罪人，是为让我们仍然做罪人吗？我们的匠人看见我们如同森林中的木料，祂心中所想的是祂将用它建造的房屋，而非未经加工的木材。同样，你看见你的仇人攻击你，发怒，用言语咬你，用侮辱激怒你，用仇恨折磨你；你在他的身上注意到这一点：他是人。你看见所有这些敌对的事，是由人做的；你在他身上看见他是天主造的。他之成为人，是天主的作为；但他恨你，是他的作为；他忌恨你，是他的作为。你在心中说什么？\Quote{主，求你仁慈待他，赦免他的罪，打击他，改变他。}你爱的不是他现在的样子，而是你希望他成为的样子。因此，当你爱仇人时，你爱的是弟兄。故此，完美的爱就是爱仇人；这完美的爱包含在兄弟之爱中。没有人可以说若望宗徒的劝导稍逊一筹，而主基督的劝导更高。若望劝导我们爱弟兄；基督劝导我们甚至爱仇人。注意基督为何命令你爱仇人。是为让他们始终作仇人吗？若祂是为这目的而命令，那么你是在恨，而不是爱。注意祂自己如何爱，\Emph{i.e.}因为祂不愿他们继续作当时的迫害者，祂说：\BibleQuote{父啊，宽赦他们吧！因为他们不知道他们做的是什么。}\BibleRef{路 23:34}祂愿意他们被赦免，就是愿意他们被改变；祂愿意他们被改变，就是屈尊将仇人变成弟兄，并且实在这样做了。祂被杀死，被埋葬，复活了，升了天；派遣圣神给祂的门徒；他们开始大胆宣讲祂的名，因被钉死的那位的名字行了奇迹；杀害主的人看见了他们；那些曾愤怒地流了祂的血的人，因着相信而喝了祂的血。\par

11. 弟兄们，我已说了这些事，并且说得有点长；然而，因为爱德需要更恳切地向你们推荐，亲爱的，就应以这种方式来推荐。因为如果你们里面没有爱德，我们等于什么也没有说。但如果你们有爱德，我们就如同在火焰上加油。在那些没有爱德的人中，或许因着这些话而燃起了爱德。在一个人的身上，已有的增长了；在另一个人的身上，从未有过的开始了存在。因此我们说这些事，好叫你们不怠于爱仇人。若有人向你发怒？他发怒，你祈祷；他仇恨，你怜悯。那恨你的是他灵魂的狂热；他将痊愈，并将感谢你。医生如何爱病人？他们是爱病人吗？若他们爱他们如同病人，他们就希望他们永远生病。医生爱病人，并非为让他们仍生病，而是为让他们从生病变成健康。他们常常要忍受疯癫病人的多少折磨！如何侮辱的话！他们甚至常常被病人打。他攻击热病，却原谅了那人。我还要说什么，弟兄们？他爱他的仇人吗？不，他恨他的仇人，即疾病；因为这才是他所恨的，却爱那个打他的人：他恨的是热病。因为他是被谁或什么打的？是被疾病、病痛、热病打的。他除去那与他敌对的东西，好让那能使他感激的存留下来。同样，你也应这样做。如果你的仇人恨你，并且恨得不义；要知道是世俗的情欲在他内作王，因此他恨你。如果你也恨他，你就以恶报恶。以恶报恶有什么用？我曾为那个恨你的病人哭泣；现在，若你也恨，我就要哀哭你了。但他攻击你的财物；他夺走了你在世上拥有的不知什么东西：因此你恨他，因为他使你在世上受窘迫。不要受窘迫，要迁到高天之上；在那里你将有你的心，那里空间宽广，使你不致在永生的希望中受窘迫。想想他夺走的是什么：他甚至不能夺走它们，除非得到那\BibleQuote{鞭打他所收纳的每一个儿子}\BibleRef{希 12:6}的准许。这个你的仇人，在某种程度上是天主手中的工具，借着他你可以得到医治。如果天主知道让你受剥夺对你有益，祂就准许他；如果天主知道让你挨打对你有益，祂就准许他打你：借着那工具，祂照顾你：愿你愿他得到痊愈。\par

12. \Quote{从来没有人见过天主。}请看，亲爱的诸位：\Quote{如果我们彼此相爱，天主就存留在我们内，祂的爱在我们内得以成全。}\BibleRef{若一 4:12}开始去爱吧，你将被成全。你已经开始爱了吗？天主已开始住在你内：要爱那已开始住在你内的天主，好叫祂因更完善地住在你内而使你成全。\Quote{我们所以知道我们存留在祂内，祂也存留在我们内，就是因为祂将祂的圣神赐给了我们。}\BibleRef{若一 4:13}说得真好：感谢天主！我们得以知道祂住在我们内。然而我们又如何知道这事，即我们知道祂住在我们内呢？因为若望自己已经说了：\Quote{因为祂将祂的圣神赐给了我们。}我们怎知道祂将祂的圣神赐给了我们呢？这事本身，即祂将祂的圣神赐给了你，你如何得以知道？问问你自己的心肝吧：如果它们充满了爱德，你就有天主的圣神。我们怎知道借此你知道天主的圣神住在你内呢？\Quote{因为天主的爱藉着所赐给我们的圣神，已倾注在我们心中。}\BibleRef{罗 5:5}\par

13. \Quote{我们看见了，并且作证：天主派遣了祂的子作世界的救主。}\BibleRef{若一 4:14}你们这些患病的人，请安心吧：这样的一位医生已经来了，你们还要绝望吗？疾病很大，伤口不治，病势危急。你注意到你的病势重大，却没有注意到医生的全能吗？你绝望，但祂是全能的；那些先被治愈、然后宣告医生的人就是祂的见证人：然而他们也是在望德中而不是在现实中痊愈。因为宗徒这样说：\Quote{因为我们得救是在于望德。}\BibleRef{罗 8:24}因此，我们已开始在信德中痊愈；但我们的痊愈将得以完全，\Quote{当这可朽坏的穿上了不朽坏的，这可死的穿上了不死的时候。}\BibleRef{格前 15:53}这是望德，不是现实。但那在望德中喜乐的人，也会握住现实；而那个没有望德的人，却不能达到现实。\par

14. \Quote{凡宣认耶稣是天主子，天主就住在他内，他也住在天主内。}我们现在可以简短地说：\Quote{凡宣认的}，不是言语上的，而是行为上的；不是用口舌，而是用生活。因为许多人在言语上宣认，但在行为上否认。\Quote{我们认识了，且相信了天主对我们所怀的爱。}\BibleRef{若一 4:16}并且，你又如何知道这事呢？\Quote{爱就是天主。}祂以上已经说过，看，祂再次说。爱不能再比被称为天主更受举荐了。或许你准备轻视天主的恩赐。但你能轻视\Emph{天主}吗？\Quote{爱就是天主：那存留在爱内的，就存留在天主内，天主也存留在他内。}二者彼此内居：那扶持的和那被扶持的。你住在天主内，好使你被扶持；天主住在你内，好使祂扶持你，免得你跌倒。免得你想象你成了天主的居所，就像你的房屋支撑你的肉身那样：如果你所在的房屋从你下面抽身，你就跌倒；但如果你抽身离去，天主并不跌倒。当你离弃祂时，祂丝毫不减；当你回到祂身边时，祂丝毫未增。你得了痊愈，你并未给祂什么；你被洁净，你被更新，你被矫正：祂是病人的药，是弯曲者的准绳，是黑暗者的光，是遗弃者的居所。因此一切都加给了你：注意不要想象你的归附给天主加了什么：不，连一个奴隶也不如。难道天主会因为你不愿意、大家都不愿意就没有仆人吗？天主不需要仆人，倒是仆人需要天主：因此圣咏说：\Quote{我曾对主说：祢是我的天主。}祂是真正的主。而它说什么呢？\Quote{因为祢不需要我的好处。}你需要你仆人那里所拥有的好处。你的仆人需要从你那里所拥有的好处，好让你喂养他；你同时需要从你仆人那里所拥有的好处，好让他帮助你。你不能自己打水，不能自己做饭，不能在你马前跑，不能照料你的牲畜。你看，你需要从你仆人那里所拥有的好处，你需要他的服侍。因此你不是真正的主，因为你还需要一个下级。祂是真正的主，祂不需要从我们那里得到任何东西；如果我们不寻求祂，我们有祸了！祂不需要从我们那里得到任何东西：然而当我们在不寻求祂时，祂却寻求了我们。有一只羊迷失了，祂找到了它，欢欢喜喜地将它扛在肩上带回。\BibleRef{路 15:4}难道羊是牧者的必需，而不是牧者是羊的必需吗？——我越喜欢谈论爱德，就越不愿这封书信结束。没有谁比它更热烈地举荐爱德了。没有比传讲给你们更甘甜的东西，没有比你们饮用更滋养的东西：但只有这样，如果你们以敬虔的生活在你们内坚定这天主的恩赐。不要对祂如此浩大的恩宠忘恩负义：祂虽只有一个独生子，却不愿祂独自为子；而是为了使祂有弟兄，就收养了那些要与祂一同拥有永生的人。\par

\SourceRecord{Translated by H. Browne.；Nicene and Post-Nicene Fathers, First Series；Vol. 7.；Edited by Philip Schaff.；Buffalo, NY: Christian Literature Publishing Co.,；1888.；<http://www.newadvent.org/fathers/170208.htm>.}

% structure: p0001 source=h1 target=section reason=page-title

\section{《若望一书》第九篇讲道}

% structure: p0002 source=h2 target=subsection reason=relative-heading-rank; base=h2

\subsection{若望一书 4:17-21}

\Quote{在这事上，爱在我们内得以成全，使我们在审判的日子坦然无惧；因为祂怎样，我们在这世上也怎样。爱里没有恐惧；但完美的爱驱除恐惧，因为恐惧含有惩罚。恐惧的人在爱内尚未得到成全。我们爱祂，因为祂先爱了我们。若有人说：我爱天主，却恨自己的弟兄，便是撒谎；因为那不爱自己所看见的弟兄的，怎能爱那看不见的天主呢？我们从祂领受了这命令：爱天主的，也当爱自己的弟兄。}\par

1. 亲爱的诸位，你们记得，若望宗徒的这部书信，尚余最后一篇有待我们靠主的恩佑来处理并向你们解释。我们对此债务念念不忘；而你们也当不忘你们所应得的权利。确实，这封书信所主要、几乎唯一赞颂的爱德，一方面使我们最忠诚地偿还债务，另一方面也使你们最甜美地行使权利。我说\Quote{最甜美地行使}，因为若没有爱德，行使权利的人是苦涩的；但有了爱德，则行使权利的人是甘甜的，而被要求者尽管承担一些辛劳，爱德却使这辛劳几乎不再成为辛劳，反倒轻省。难道我们没看见，即使在不能言语、无理性的动物中，当爱不是属灵的，而是属血肉和本性的，母亲会以极大的深情将自己献给要喝应得奶水的幼崽？无论乳婴多么猛烈地冲向乳头，这对母亲来说也比幼崽不吸、不索取那因爱而应得的更好。我们常见大牛犊用头猛顶母牛的乳房，几乎把母牛的身体顶起，但母牛并不踢开它们；若幼崽不在那里吸奶，母牛的哞叫声呼唤它来乳头。若在我们内有那属灵的爱德，宗徒所说\Quote{我在你们中间成为卑微的，如同乳母抚育自己的幼儿}\BibleRef{得前 2:7}，那么当你们行使权利时，我们更爱你们。我们不喜欢懒散的人，因为我们为软弱者担忧。然而，因为某些指定的读经插入，必须在它们的圣日宣读并讲道，我们不得不中断此书信的连续诵读。现在让我们回到被打断的顺序；其余部分，圣洁的弟兄们，请你们全神贯注地领受。我不知道爱德还能怎样更壮丽地被赞美，以至于说\Quote{爱德是天主}\BibleRef{若一 4:16}。简短的赞美，却是伟大的赞美：简短于言辞，却浩大于意义！说话如此匆匆：\Quote{爱就是天主！}这也是简短的：若计算，它是一个字；若衡量，它何其伟大！\Quote{爱是天主，那住在爱内的，}他说，\Quote{就住在天主内，天主也住在他内。}让天主成为你的房屋，你也成为天主的房屋；你住在天主内，也让天主住在你内。天主住在你内，为扶持你；你住在天主内，为不致跌倒；因为宗徒论及这同一的爱德说：\Quote{爱永不会跌倒。}被天主扶着的人怎能跌倒呢？\par

2. \Quote{在我們內，愛情得以成全，使我們在審判之日能有膽量，因為正如祂是那樣，我們在這世界上也是那樣。} \BibleRef{若一 4:17} 祂告訴我們每個人如何證明自己，愛德在他內進步了多少，或者更確切地說，他在愛德中進步了多少。因為如果愛德是天主，天主既不增益也不減損：所謂愛德在你內進步，僅意味着你在他內進步。因此，問問你在愛德中進步了多少，你的心會如何回答你，好讓你明白自己長進的程度。因為祂應許向我們顯示我們在什麼上認識祂，並說：\Quote{在此，愛情在我們內得以成全。} 問，在什麼上？\Quote{使我們在審判之日能有膽量。} 凡在審判之日有膽量的人，在他內愛德得以成全。什麼是在審判之日有膽量？不怕審判日子的來臨。有些人拒絕相信有審判之日；他們不能對他們不相信會來臨的日子有膽量。我們越過這些人：願天主喚醒他們，使他們能生活；為什麼我們談論死人？他們不相信有審判之日；他們既不害怕也不渴望他們不相信的事物。有人開始相信審判之日：如果他開始相信，他也開始害怕。但因為他仍在害怕，因為他尚在審判之日沒有膽量，愛德在那人內尚未成全。但是，這就該絕望嗎？在你看到開端的人身上，為何你對終局絕望呢？我看到了怎樣的開端？（你說。）就是那害怕。聽聖經：\Quote{敬畏上主是智慧的開端。} 那麼，他已經開始害怕審判之日：藉着害怕，讓他改正自己，讓他提防他的仇敵，即他的罪；讓他開始在內心復活，並致死他在地上的肢體，如宗徒所說：\Quote{致死你們在地上的肢體。} \BibleRef{哥 3:5} 地上的肢體，他指的是屬靈的邪惡：因為他接着解釋：\Quote{貪婪、不潔} \BibleRef{弗 6:12} 以及他隨後列舉的其他。現在，隨着這個開始害怕審判之日的人，在地上肢體中死去的程度，天上的肢體也相應地升起並被堅固。但天上的肢體都是善工。當天上的肢體升起時，他開始渴望他曾經害怕的那一天。從前他害怕基督來臨，會在他身上找到祂必須定罪的惡人；現在他渴望祂來臨，因為祂會找到祂可以加冕的義人。他既已開始渴望基督來臨，那貞潔的靈魂渴望新郎的擁抱，便棄絕了那姦夫，藉着信德、望德、愛德在內心成為貞女。如今這人在審判之日有了膽量：當他祈禱\Quote{願你的國來臨} \BibleRef{瑪 6:10} 時，他不與自己爭戰。因為害怕天主的國來臨的人，是害怕他的祈禱被俯聽。怎能說一個害怕祈禱被俯聽的人是在祈禱呢？但以愛德的膽量祈禱的人，現在願意祂來臨。關於同一渴望，聖詠中有人說：\Quote{至於你，上主，要到幾時？求你回轉，救我的靈魂。} 他因被拖延而嘆息。因為有些人以忍耐接受死亡；但有些成全的人卻以忍耐忍受生活。我是什麼意思？當一個人仍渴望此生時，在死亡日來臨時，他忍耐地接受死亡：他與自己抗爭，為要隨從天主的旨意，心中渴望天主所選擇的，而非人的意志所選擇的；由於對現世生命的渴望，對死亡產生抗拒，但他仍以忍耐和剛毅面對死亡，從容赴死。但當一個人渴望，如宗徒所說，\Quote{脫離肉身，與基督同在} \BibleRef{斐 1:23} 時，那個人不是忍耐地死，而是忍耐地活着，喜樂地死。請看宗徒忍耐地活着，即他是如何以忍耐在此世不是愛生命，而是忍受生命。\Quote{脫離肉身，與基督同在，} 他說，\Quote{是遠更好；但繼續在肉身內，為了你們是必要的。} 因此，弟兄們，你們要努力，在內心裏決意以此為你們的關切，好使你們渴望審判之日。否則，愛德就不能被證明是成全的，除非當一個人開始渴望那一日。但是那有膽量的人渴望它，他的良心在完美而真誠的愛德中不感到驚慌。\par

3. \BibleQuote{在这一点上，天主的爱在我们内得以成全，使我们在审判之日能有\OriginalTerm{信赖}{boldness}。} 为什么我们会有信赖？\BibleQuote{因为祂怎样，我们在这世上也怎样。} 你们已经听到了你们信赖的根据：\BibleQuote{因为祂怎样，} 宗徒说，\BibleQuote{我们在这世上也怎样。} 他难道不是说了一件不可能的事吗？因为人怎能如同天主呢？我已经向你们解释过，\Quote{\Emph{如同}} 一词并不总是表示相等，而是表示某种相似。因为你怎么说：\Quote{如同我有耳朵，\Emph{同样}我的肖像也有？} 它果真\Emph{同样}吗？然而你却说：\Quote{\Emph{同样，如同。}} 如果我们乃是按天主的肖像受造，为何我们不\Emph{如同}天主呢？并非达到相等，而是相对于我们的程度。那么，我们在审判之日何以得到信赖呢？\BibleQuote{因为祂怎样，我们在这世上也怎样。} 我们必须将此联系到同一个\OriginalTerm{爱德}{caritas}，并理解其含义。主在福音中说：\BibleQuote{如果你们爱那爱你们的人，你们有什么赏报呢？税吏不也是这样作吗？}\BibleRef{玛 5:44} 那么祂要我们作什么呢？\BibleQuote{我却对你们说：爱你们的仇人，为迫害你们的人祈祷。} 如果祂命令我们爱仇人，祂从哪里引出一个榜样来给我们呢？从天主自己：因为祂说：\BibleQuote{好使你们成为你们在天之父的子女。} 天主如何作呢？祂爱祂的仇人，\BibleQuote{祂使太阳上升，光照恶人，也光照善人；降雨给义人，也给不义的人。} 如果这就是天主邀请我们达到的成全，即我们爱仇人如同祂爱了仇人一样；那么，这就是我们在审判之日的信赖：\BibleQuote{祂怎样，我们在这世上也怎样：} 因为，如同祂爱祂的仇人，使太阳上升光照善人也光照恶人，降雨给义人也给不义的人，同样，我们既然不能赐给他们太阳和雨水，就在我们为他们祈祷时，赐给他们我们的眼泪。\par

4. 现在，关于这同一个信赖，让我们看看他说什么。我们如何知道爱德是成全的呢？\BibleQuote{在爱德内没有\OriginalTerm{恐惧}{fear}。}\BibleRef{若一 4:18} 那么，对于那开始恐惧审判之日的人，我们该说什么呢？如果在他内的爱德是成全的，他就不会恐惧。因为成全的爱德会产生完全的正义，他就没有什么可恐惧的；他反而会有所渴望：愿不义过去，天主的国来临。因此，\BibleQuote{在爱德内没有恐惧。} 但在什么样的爱德内呢？不是在开始的爱德内；那么在哪里呢？\BibleQuote{成全的爱德，} 他说，\BibleQuote{驱逐恐惧。} 那么，让恐惧作为开始，因为\BibleQuote{敬畏上主是智慧的开端。} 恐惧，可以说，为爱德预备了地方。但一旦爱德开始居住，那为它预备地方的恐惧就被驱逐出去了。因为爱德增长多少，恐惧就减少多少；爱德越是在内里增长，恐惧就越被驱逐。爱德越大，恐惧越小；爱德越小，恐惧越大。但若没有恐惧，爱德就无法进入。正如我们在缝纫中所见，线是通过鬃毛引入的；鬃毛先进入，但除非它出来，线就不会进入它的位置；同样，恐惧先占据心灵，但恐惧并不停留那里，因为它进入只是为了引入爱德。一旦心灵有了安全感，我们在这世上和来世将有何等喜乐！即使在这世上，谁又能伤害我们，当我们充满爱德时？请看宗徒如何因这爱德而欢欣：\BibleQuote{谁能使我们与基督的爱隔绝呢？是困苦吗？是窘迫吗？是迫害吗？是饥饿吗？是赤贫吗？是危险吗？是刀剑吗？}\BibleRef{罗 8:35} 伯多禄也说：\BibleQuote{如果你们热心行善，谁能伤害你们呢？——在爱德内没有恐惧，但成全的爱德驱逐恐惧，因为恐惧含有刑罚。}\BibleRef{伯前 3:13} 罪的意识折磨心灵；成义尚未发生。其中有某种发痒、刺痛的东西。因此，在圣咏中，关于这同样的正义的成全，他说什么呢？\BibleQuote{祢将我的哀哭变为舞蹈，脱去了我的麻衣，给我披上喜乐，好使我的光荣歌颂祢，而不致缄默。} 这是什么意思？\Quote{不致缄默？} 即没有那会刺伤我良心的东西。恐惧确实刺痛；但不要恐惧：爱德进入，它治愈了恐惧所施加的创伤。对天主的恐惧像外科医生的刀一样割伤；它除去腐烂，似乎使伤口更大。看，当腐烂在身体内时，伤口较小，但危险；然后刀来了；伤口在医生切割时比先前更痛。在治疗过程中它比不治疗时更痛；但在治疗中更痛，只是为了治疗后永远不再痛。那么，让恐惧占据你的心，以便引入爱德；让疤痕取代医生的刀。祂是这样的医生，疤痕甚至不显现；只管将自己放在祂手中。因为如果你没有恐惧，你就不能成义。这是圣经所宣告的判词：\BibleQuote{没有恐惧的人不能成义。}\BibleRef{德 1:28} 所以恐惧必须先进入，以便通过它爱德来到。恐惧是治疗手术；爱德是健康状态。\BibleQuote{那恐惧的，在爱德内尚未成全。} 为什么？\BibleQuote{因为恐惧含有刑罚；} 正如外科医生的刀切割含有刑罚一样。\par

5. 但另有一句话，若无人理解，似乎与此相悖。即在圣咏某处写道：\BibleQuote{对主的敬畏是纯洁的，永远存留。}祂向我们显示一种永恒的敬畏，却是\Emph{纯洁的}。但如果祂在那里显示一种永恒的敬畏，那么这封书信或许与祂相悖，因为它说：\BibleQuote{爱内没有恐惧，但圆满的爱驱除恐惧？}让我们质问这两句天主的话。神是同一位神，虽然有两本书，两张口，两根舌头。因为这是由若望的口说出，那是由达味的口说出：但不要以为圣神不止一位。如果一口气充满两支管子，难道同一位圣神不能充满两颗心，感动两根舌头？但如果两根被同一气息充满的管子发出和谐的声音，那么被天主的圣神或气息充满的两根舌头岂能发出不和谐的声音？因此那里有和谐，有协和，只是需要一个能听的人。看，天主的圣神已吹入并充满两颗心，感动了两根舌头：我们从一个舌头听到：\BibleQuote{爱内没有恐惧；但圆满的爱驱除恐惧；}从另一个舌头听到：\BibleQuote{对主的敬畏是纯洁的，永远存留。}这是如何？音符似乎刺耳。并非如此：竖起你的耳朵：注意旋律。并非无故在一处加上了\Quote{纯洁的}这个词，在另一处没有加：因为有一种敬畏称为纯洁的，另一种敬畏不称为纯洁的。让我们区分这两种敬畏的差别，从而理解这两根管子的和谐。我们如何理解，或如何区分？请注意，我亲爱的。有人敬畏主，唯恐被投入地狱，唯恐与魔鬼在永火中焚烧。这是引入爱德的敬畏：但它来是为了离去。因为你若因惩罚而敬畏主，你还不爱惜你所如此敬畏的那一位。你并不渴求美善，却害怕邪恶。但因为你害怕邪恶，你纠正自己，开始渴求美善。一旦你开始渴求美善，你身上就会有纯洁的敬畏。什么是纯洁的敬畏？是害怕失去这些美善本身。注意！害怕主把你与魔鬼一起投入地狱是一回事，害怕主抛弃你是另一回事。你因害怕被与魔鬼一起投入地狱而敬畏主，这还不纯洁；因为它不是出于对主的爱，而是出于对惩罚的恐惧；但当你害怕主离开你时，你拥抱祂，你渴望享受天主本身。\par

6. 最好用两个已婚妇女的例子来说明这两种敬畏的区别，一种被爱德驱除，另一种纯洁、永远存续。假设一个妇女愿意犯奸淫，以邪恶为乐，只怕被丈夫定罪；她害怕丈夫，但因她仍爱邪恶，所以她害怕丈夫。对这位妇女而言，丈夫的出现并不令她喜悦，而是负担；如果她碰巧生活邪恶，她害怕丈夫，怕他来。这样的人害怕审判之日的来临。设另一位妇女爱丈夫，觉得自己欠他纯洁的拥抱，没有以任何奸淫的污秽玷污自己；她渴望丈夫的到来。这两种敬畏如何区分？一个妇女害怕，另一个也害怕。问她们：似乎回答相同；问第一个：\Quote{你怕丈夫吗？}她回答：\Quote{我怕。}问另一个是否怕丈夫，她回答：\Quote{我怕。}声音相同，心思不同。那么现在问她们\Quote{为什么}？一个说：\Quote{我怕丈夫，怕他回来。}另一个说：\Quote{我怕丈夫，怕他离开我。}一个说：\Quote{我怕被定罪。}另一个说：\Quote{我怕被抛弃。}让同样的事发生在基督徒的心里，你就会发现一种爱德所驱除的敬畏，和另一种纯洁的、永远存续的敬畏。\par

7. 那么，让我们先对那些敬畏天主，却像那个以罪恶为乐的女人一样的人说话；也就是说，她害怕丈夫，怕他定她的罪；让我们先对这样的人说话。哦，灵魂啊，你敬畏天主，怕祂定你的罪，就像那以罪恶为乐的女人一样：她害怕丈夫，怕被丈夫定罪：正如你厌恶这个女人，也要厌恶你自己。如果你偶然有妻子，你愿意你的妻子这样怕你，只因为怕被你定罪吗？她以罪恶为乐，却只靠对你的恐惧的压制来约束，而不是因为她对自己罪孽的谴责？你希望她贞洁，好让她爱你，而不是怕你。你要在天主面前如此，就像你希望你的妻子在你面前那样。如果你还没有妻子，并且希望有一个，你会希望她如此。然而，弟兄们，我们说的是什么呢？那个女人，她对丈夫的恐惧就是被丈夫定罪，也许她没有犯奸淫，免得用什么方式传到丈夫耳中，被他剥夺这暂时的生命之光；但丈夫可能受骗，也可能被蒙在鼓里；因为他不过是人，正如那个能欺骗他的女人一样。她害怕他，而他却可能看不见她：难道你不害怕你的丈夫永远注视你的面容吗？\BibleQuote{恶人的面容必被主抛弃。} 她趁丈夫不在，也许被奸淫之乐所诱惑；但她对自己说：我不这样做；他确实不在，但很难避免以某种方式让他知道。她克制自己，免得被一个必朽的人知道，他也有可能永远不知道，也可能受骗，以至于把一个坏女人当做好女人，把奸妇当做贞洁的妇女：而你不害怕那无人能欺骗的眼睛吗？你不害怕那不能离开你面前的天主吗？求天主垂顾你，转面不看你的罪恶；\BibleQuote{求祢转面不看我的罪恶。} 但你凭什么配得祂转面不看你的罪恶，如果你自己不转面不看你的罪恶？因为圣咏中同样的话语说：\BibleQuote{因为我承认我的过犯，我的罪恶常在我的眼前。} 承认吧，祂必宽恕。\par

8. 我们已经对那些尚存那不能存留到永远的恐惧的灵魂说话，这种恐惧被爱德驱逐出去；现在让我们对那拥有贞洁的、存留到永远的恐惧的灵魂说话。我们会找到那个灵魂吗，你想想，好让我们对它说话？你认为，它在这会众中吗？你认为，它在这祭台间吗？你认为，它在地上吗？它必然存在，只是隐藏着。现今是冬天；内在的根是青绿的。或许我们能触及那灵魂的耳朵。但无论那灵魂在哪里，哦，但愿我能找到它，不是让它听我说，而是我倾听它！它应该教导我一些东西，而不是向我学习！一个圣洁的灵魂，火热的灵魂，渴慕天主的国；那个灵魂，不是我对它说话，而是天主亲自对它说话，并这样安慰它，当它耐心地忍受在世上的生活时：\BibleQuote{你希望我现在就来，我知道你希望我现在就来：我知道你是什么样的，以至于你能毫无恐惧地等待我的来临；我知道这对你是一种烦恼；但你还要再等，忍受；我来了，而且很快就要来。} 但对于那爱慕的灵魂，时间过得很慢。听她歌唱，像一朵在荆棘中的百合；听她叹息着说：\BibleQuote{我要歌唱，并以无瑕的方式领会：你何时会到我这里来？} 但她以无瑕的方式也许并不害怕；因为\BibleQuote{完美的爱德驱逐恐惧。} 当祂来拥抱她时，她仍然害怕，但以一种感到安全的方式。她害怕什么？她要当心，谨慎提防自己的罪孽，以免再犯罪；不是怕被丢进火里，而是怕被祂抛弃。她里面将有什么？那\BibleQuote{贞洁的、存留到永远的恐惧。} 我们已经听到了两支笛子和谐地吹奏。那支讲到恐惧，这支也讲到恐惧：但那是灵魂怕被定罪的恐惧；这是灵魂怕被抛弃的恐惧。那是被爱德驱逐的恐惧；这是存留到永远的恐惧。\par

9. \BibleQuote{我们应当爱，因为祂先爱了我们。} \BibleRef{若一 4:19} 因为我们怎能爱，除非祂先爱了我们？因爱我们成为朋友，但祂爱我们还是仇敌时，好使我们成为朋友。祂先爱了我们，并赐给我们爱祂的恩赐。我们尚未爱祂：因爱我们成为美丽的。如果一个畸形丑陋的人爱一个美丽的女子，他能做什么？或者一个女子，如果她畸形丑陋、面色黝黑，爱一个英俊的男子，她能做什么？因爱她能变得美丽吗？他因爱能变得英俊吗？他爱一个美丽的女子，当他在镜子中看见自己时，便羞于向她——他所倾慕的这位可爱者——抬起脸来。他怎样才能美丽？他等待美貌来临吗？不，反而因等待，年岁增加，使他更丑陋。那么没有可做的，没有方法可以劝告他，只有他应当自制，不妄图不相称的爱；或者，如果他确实爱她，并希望娶她为妻，那么在他身上，让他爱贞洁，而非血肉的面容。但是，我的弟兄们，我们的灵魂因不义而不可爱；因爱天主而变得可爱。那使爱者美丽的爱，该是怎样的爱啊！但天主永远是可爱的，从不可恶，从不改变。永远可爱的那位先爱了我们；祂爱我们时，我们是什么，岂非污秽和不可爱吗？但并非要我们停留在污秽中；不，而是要改变我们，从不可爱变为可爱。我们如何变得可爱？藉着爱那永远可爱者。随着你里面的爱增长，可爱也增长：因为爱本身就是灵魂的美。\BibleQuote{我们应当爱，因为祂先爱了我们。} 听宗徒保禄说：\BibleQuote{但天主在我们身上显示了祂的爱，因为当我们还是罪人的时候，基督为我们死了：} \BibleRef{罗 5:8} 义人代替不义的人，美丽的代替污秽的。我们如何找到耶稣是美丽的？\BibleQuote{你在美貌中超越世人，恩宠倾注在你的唇上。} 为何如此？再看为何祂是俊美的；\BibleQuote{你在美貌中超越世人：} 因为\BibleQuote{在起初已有圣言，圣言与天主同在，圣言就是天主。} \BibleRef{若 1:1} 但祂取了血肉，祂取了，可以说，你的丑陋，\Emph{即}你的必朽性，以便祂能适应你，与你相配，并激发你爱内在的美。那么，在圣经中何处我们找到耶稣不俊美、畸形，如我们找到祂俊美且\BibleQuote{你在美貌中超越世人}呢？何处我们也找到祂畸形？问依撒依亚：\BibleQuote{我们看见了祂，祂没有形态，也没有俊美。} \BibleRef{依 53:2} 现在有两支笛子，似乎发出不协调的声音；但同一圣神吹入二者。藉此说：\BibleQuote{你在美貌中超越世人：} 凭那说在依撒依亚：\BibleQuote{我们看见了祂，祂没有形态，也没有俊美。} 同一圣神充满两支笛子，它们没有不谐。不要转开你的耳朵，运用理解。让我们问宗徒保禄，让他为我们解释两支笛子的和谐。让他为我们吹奏这一音符：\BibleQuote{你在美貌中超越世人。——祂虽具有天主的形态，却不以自己与天主同等为应当把持的。} \BibleRef{斐 2:6} 也让他为我们吹奏这一音符：\BibleQuote{我们看见了祂，祂没有形态，也没有俊美。——祂使自己空虚，取了奴仆的形态，成了人的模样，且以人的形状出现。祂没有形态，也没有俊美，} \BibleRef{斐 2:6} 为赐给你\OriginalTerm{形态}{form}和\OriginalTerm{俊美}{comeliness}。什么形态？什么俊美？那在\OriginalTerm{爱德}{caritas}中的爱：好使你们爱慕，奔跑；\BibleRef{歌 1:4} 奔跑，爱慕。你们现在美丽了：但不要将目光停留在自己身上，免得失去你所领受的；让你的目光止于那使你美丽者。只为祂爱你而美丽。但你要将一切目标指向祂，奔向祂，寻求祂的拥抱，害怕离开祂；好使在你内常有那\OriginalTerm{纯洁的敬畏}{chaste fear}，它永远长存。\BibleQuote{我们应当爱，因为祂先爱了我们。}\par

10. \Quote{若有人说：我爱天主。} \BibleRef{若一 4:20} 哪一个天主？我们为什么爱？\Quote{因为祂先爱了我们，}并赐给我们去爱。祂爱我们这些不敬虔的人，为使我们成为敬虔；爱我们这些不义的人，为使我们成为义人；爱我们这些患病的人，为使我们痊愈。问每一个人，让他告诉你他是否爱天主。他呼喊，他承认：\Emph{我爱}，天主知道。还有另一个问题要问。\Quote{若有人说：我爱天主，却恨他的弟兄，便是撒谎。} 你凭什么证明他是撒谎者？请听。\Quote{因为那不爱所见之弟兄的，怎能爱那未见的天主呢？} 那么怎样？爱弟兄的人也爱天主吗？他必然爱天主，必然爱那爱本身。一个人能爱弟兄，却不爱爱吗？他必然爱爱。那么怎样？因为他爱爱，就能推论他爱天主吗？当然能推论。在爱爱时，他爱天主。或者你忘了你刚才所说的话：\Quote{爱就是天主。} \BibleRef{若一 4:8} 如果\Quote{爱就是天主，}那么谁爱爱，谁就爱天主。所以，爱你的弟兄，你就可安心。你不能说：\Quote{我爱我的弟兄，但我不爱天主。} 正如你说\Quote{我爱天主}而你不爱弟兄时你在撒谎，同样，当你说\Quote{我爱我的弟兄}却认为你不爱天主时，你也被欺骗了。凡爱弟兄者，必然爱爱本身；但\Quote{爱就是天主：}所以，凡爱弟兄者必然爱天主。但如果你不爱所见之弟兄，怎能爱那未见的天主呢？他为什么看不见天主？因为他没有爱本身。他看不见天主，是因为他没有爱；他没有爱，是因为他不爱弟兄。因此，他看不见天主的原因，是他没有爱。因为如果他有了爱，他就看见天主，因为\Quote{爱就是天主：}并且那眼睛正藉着爱日益洁净，以便看见那不变的本体，在祂面前他将永远喜乐，永远享受祂，当他与天使结合时。只是，他现在要奔跑，以便最终在自己的家乡得享喜乐。让他不要爱他的流放，不要爱道路；让一切都成为苦的，除了那呼唤我们的那一位，直到我们紧紧抓住祂，并说出圣咏中所说的：\Quote{祢消灭了一切远离祢而行淫的人}——那些行淫的是谁？就是那些离开并爱世界的人。但你要做什么？祂继续说：——\Quote{但我亲近天主为善。} 我的一切好处就是紧紧依附于天主，白白地。因为如果你问他，说：你为什么依附祂？他应该说：为使祂给我——给你什么？是祂创造了天，祂创造了地：祂要给你什么？你已经依附于祂了：找更好的东西，祂会赐给你。\par

11. \Quote{因为那不爱所见之弟兄的，怎能爱那未见的天主呢？我们从祂领受了这命令：爱天主的，也当爱他的弟兄。} \BibleRef{若一 4:20} 你所说的话真是奇妙的高论：\Quote{我爱天主，}却恨你的弟兄！哦，杀人者，你怎么爱天主？难道你没有在这同一书信的上文听到：\Quote{凡恨自己弟兄的，就是杀人者} \BibleRef{若一 3:15}？是的，但我确实爱天主，尽管我恨我的弟兄。你确实不爱天主，如果你恨你的弟兄。现在我用另一个证明来说明这一点。这同一位宗徒说：\Quote{祂命令我们应彼此相爱。} 你怎么能说你爱祂，却恨祂的命令？谁能说：\Quote{我爱皇帝，但我恨他的法律？} 皇帝正是藉此知道你爱他：即他的法律在各省得到遵守。我们皇帝的法律是什么？\Quote{我给你们一条新命令：你们当彼此相爱。} \BibleRef{若 13:34} 那么，你说你爱基督：遵守祂的命令，爱你的弟兄。但如果你不爱你的弟兄，你怎么能说你爱祂，却藐视祂的命令？弟兄们，我奉主的名谈论爱德，从不厌倦。你们对这件事有多少无止境的渴望，我们也希望这件事本身在你们内增长，并驱逐恐惧，好使那纯洁的、永远长存的敬畏存留。让我们忍受世界，忍受磨难，忍受诱惑的绊脚石。让我们不要离开道路；让我们持守教会的合一，持守基督，持守爱德。让我们不要被从祂净配的肢体中夺去，不要从信德中被夺去，好使我们在祂来临时得享光荣。并且我们要安全地住在祂内，现今凭着信德，那时凭着目睹，我们拥有如此大的保证，就是圣神的恩赐。\par

\SourceRecord{Translated by H. Browne.；Nicene and Post-Nicene Fathers, First Series；Vol. 7.；Edited by Philip Schaff.；Buffalo, NY: Christian Literature Publishing Co.,；1888.；<http://www.newadvent.org/fathers/170209.htm>.}

% structure: p0001 source=h1 target=section reason=page-title

\section{若望一书第十篇讲道}

% structure: p0002 source=h2 target=subsection reason=relative-heading-rank; base=h2

\subsection{若望一书 5:1-3}

\BibleQuote{凡相信耶稣就是基督的，是由天主而生的；凡爱那生祂的，也爱那由祂而生的。我们由此知道我们爱天主的子女，是因为我们爱天主，并遵行祂的诫命。因为这就是对天主的爱：我们遵守祂的诫命。}\par

1. \Emph{我想}你们还记得，你们当中昨天在场的人，我们的讲解进行到这封书信的什么地方了：即\BibleQuote{不爱他所看见的弟兄，怎能爱他所看不见的天主呢？并且我们从祂领受了这条诫命：爱天主的，也当爱他的弟兄。} \BibleRef{若一 4:20} 到此为止我们讲论了。那么让我们看看接下来是什么。\BibleQuote{凡相信耶稣就是基督的，是由天主而生的。} \EditorialNote{若望一书5:1} 谁是不相信耶稣就是基督的呢？就是那不像基督所命令的那样生活的人。因为许多人说：\Quote{我相信}，但没有行为的信德不能得救。如今信德的行为就是爱德，正如宗徒保禄所说：\BibleQuote{以爱德行事之信德}。\BibleRef{迦 5:6} 你们过去的行为，在你们相信之前，要么没有，要么看似良善，却毫无价值。因为如果没有行为，你们就像一个没有脚或脚有伤不能行走的人；但如果行为看似良善，在你们相信之前，你们确实在奔跑，但跑偏了道路，反而迷了路，不能到达终点。所以，我们既要奔跑，也要在道路上奔跑。偏离道路奔跑的，跑也徒然，更确切地说，只是白费力气。他越偏离道路，就越是迷路。我们奔跑的道路是什么？基督告诉我们：\BibleQuote{我就是道路。} \BibleRef{若 14:6} 我们奔向的家园是什么？\BibleQuote{我就是真理。} 借着祂你们奔跑，向祂你们奔跑，在祂内你们安息。但为了使我们能借着祂奔跑，祂甚至伸向我们：因为我们原是遥远的、在异乡的客旅。不仅我们身处异乡，我们还软弱无力，动弹不得。祂是医生，来到病人中间；祂是道路，将自己延伸至远方的人。让我们由祂得救，让我们在祂内行走。这就是\Quote{相信耶稣就是基督}，如同基督徒那样相信，他们不只在名义上是基督徒，更是以实际行动和生命，不像魔鬼那样相信。因为圣经告诉我们：\BibleQuote{魔鬼也信，且战栗。} \BibleRef{雅 2:19} 魔鬼所能相信的，难道比他们说\BibleQuote{我们知道你是谁，你是天主子}还多吗？魔鬼所说的，伯多禄也说了同样的话。当主问他们祂是谁，人们说祂是谁，门徒回答祂说：\BibleQuote{有人说你是若翰洗者，有人说是厄里亚，还有人说是耶肋米亚，或是先知中的一位。\InnerQuote{祂对他们说：但你们说我是谁？}伯多禄回答说：\InnerQuote{你是基督，永生天主之子。}} \BibleRef{玛 16:13} 这是他从主那里听到的：\BibleQuote{西满·巴尔约纳，你是有福的，因为不是血肉启示了你，而是我在天之父。} 看，这种信德带来何等的赞美。\BibleQuote{你是伯多禄，在这磐石上我要建立我的教会。} \Quote{在这磐石上我要建立我的教会}是什么意思？是在这种信德上；是在所说的\BibleQuote{你是基督，永生天主之子}上。祂说：\BibleQuote{在这磐石上，我要建立我的教会。} 多么大的赞美！于是，伯多禄说：\BibleQuote{你是基督，永生天主之子。} 魔鬼也说：\BibleQuote{我们知道你是谁，你是天主子，天主的圣者。} 伯多禄说了这话，魔鬼也说了同样的话：言语相同，心思却不同。伯多禄说这话如何显得出于爱德呢？因为基督徒的信德带有爱德，而魔鬼的信德没有爱德。怎么没有爱德呢？伯多禄说这话，是为拥抱基督；魔鬼说这话，是为让基督离开他们。因为在他们说\BibleQuote{我们知道你是谁，你是天主子}之前，他们说：\BibleQuote{我们与你有何相干？你为何在时候未到之前来毁灭我们？} 所以，承认基督是为了抓住基督是一回事，承认基督是为了驱赶基督是另一回事。所以你们看，他在这里说\Quote{凡相信的}，是指一种属于自己的信德，而不是与许多人共有的信德。因此，弟兄们，不要让任何异端者对你们说：\Quote{我们也相信。} 为此，我以魔鬼的例子给你们一个实例，叫你们不要因相信的话语而欢喜，而要仔细查考生活的行为。\par

2. 那么让我们看看相信基督是什么意思；相信耶稣就是基督是什么意思。他接着说：\BibleQuote{凡相信耶稣就是基督的，是由天主而生的。} 但相信那是什么意思？\BibleQuote{凡爱那生祂的，也爱那由祂所生的。} 他立刻将爱德与信德结合，因为没有爱德的信德毫无价值。有爱德，是基督徒的信德；没有爱德，是魔鬼的信德；但不信的人比魔鬼更糟糕，比魔鬼更愚笨。有人不信基督：那么，他甚至比不上魔鬼。有人相信基督，却恨基督：他怀有信德的宣认是出于害怕惩罚，而非出于对冠冕的爱德；魔鬼也害怕受罚。给这信德加上爱德，使它成为宗徒保禄所说的那种信德，即\BibleQuote{以爱德行事之信德} \BibleRef{迦 5:6}。你找到了一个基督徒，找到了耶路撒冷的市民，找到了天使的同乡，找到了在路上叹息的朝圣者：与他结合，他是你的旅伴，与他一同奔跑，如果你也是这样。\BibleQuote{凡爱那生祂的，也爱那由祂所生的。} 谁\Quote{生了}？圣父。谁\Quote{被生}？圣子。那么他说什么？\Quote{凡爱圣父的，也爱圣子。}\par

3. \BibleQuote{凭此我们就知道我们爱天主的子女。} \BibleRef{若一 4:2} 弟兄们，这是什么意思？刚才他还在谈论天主子，而不是天主的子女；看，一位基督呈现在我们面前默想，我们被告知：\BibleQuote{凡信耶稣为基督的，是生自天主的；凡爱生他之父的，}即圣父，\BibleQuote{也爱那由他所生的，}即圣子，我们的主耶稣基督。他接着说：\BibleQuote{凭此我们就知道我们爱天主的子女；}仿佛他正要说话：\Quote{凭此我们就知道我们爱天主子。}他说\Quote{天主的子女，}而刚才他却在谈论天主子——因为天主的子女是天主独生子的身体，当他作为头，我们作为肢体时，便是一个天主子。因此，凡爱天主的子女的，就爱天主子；凡爱天主子的，就爱圣父；没有人能爱圣父而不爱圣子，凡爱子女的，也爱天主子。什么是天主的子女？就是天主子的肢体。借着爱，他自身也成为肢体，并通过爱进入基督身体的构造中，于是便有一位基督，爱他自己。因为当肢体彼此相爱时，身体就爱自己。\BibleQuote{若是一个肢体受苦，所有肢体都一同受苦；若是一个肢体得荣耀，所有肢体都一同欢乐。} \BibleRef{格前 12:26} 然后他接着说：\Quote{如今你们是基督的身体，各自为肢体。}若望刚才在谈论兄弟之爱，说：\BibleQuote{那不爱他所看见的弟兄的，怎能爱他所看不见的天主呢？} \BibleRef{若一 4:20} 但如果你爱你的弟兄，也许你爱你的弟兄却不爱基督？怎么可能呢，当你爱基督的肢体时？因此，当你爱基督的肢体时，你就爱基督；当你爱基督时，你就爱天主子；当你爱天主子时，你也爱圣父。因此，爱不能被分割成部分。选择你所爱的，其余的会跟随你。假设你说：我只爱天主，天主圣父。你在撒谎：如果你爱，你就不是只爱他；但如果你爱圣父，你也爱圣子。看，你说：我爱圣父，我也爱圣子：但仅仅是这样，圣父天主和圣子天主，我们的主耶稣基督，他升了天，坐在圣父的右边，那圣言，万物藉他而造，并且\BibleQuote{圣言成了血肉，住在我们中间：}我只爱这个。你在撒谎；因为如果你爱头，你也爱肢体；但如果你不爱肢体，你也不爱头。难道你不畏惧那头从天上为他的肢体发出的声音吗：\BibleQuote{扫禄，扫禄，你为什么迫害我？} \BibleRef{宗 9:4} 迫害他肢体的，他称为迫害他的人；爱他的，就是爱他肢体的人。现在，什么是他的肢体？你们知道，弟兄们，不外乎天主的教会。\Quote{凭此我们就知道我们爱天主的子女，即我们爱天主。}如何？难道天主的子女是一回事，天主自己是另一回事吗？但那爱天主的，也爱他的诫命。天主的诫命是什么？\BibleQuote{我给你们一条新命令，你们要彼此相爱。} \BibleRef{若 13:34} 谁也不要以另一种爱为自己开脱，因为另一种爱；只有这种爱是这样：正如爱本身凝聚为一，所有依附于它的也被合而为一，如同火将它们熔化为一。它是金子：金块熔化后成为某一种东西。但除非\OriginalTerm{爱德}{charity}的热火被应用，否则众多碎片无法熔合为一。\Quote{我们爱天主，}凭此\Quote{我们就知道我们爱天主的子女。}\par

4. 我们又凭什么知道我们爱天主的子女呢？凭这一点：\Quote{我们爱天主，并遵守他的诫命。}我们在此叹息，因为遵守天主的诫命实属艰难。请听下文。人啊，你在爱什么中劳苦？在爱贪婪中。你所爱的事物是以劳苦为代价；爱天主却毫无劳苦。贪婪会命令你劳作、冒险、忍受严酷的艰辛和磨难；而你却会服从它。目的是什么？为了让你拥有可以装满箱子、失去内心平安的东西。也许你此前比拥有它时更感到安全。看，贪婪命令了你什么。你装满了房屋，却惧怕盗贼；你得到了黄金，却失去了睡眠。看，贪婪命令了你什么。你做了，而且做到了。天主命令你什么呢？爱我吧。你爱黄金，你会寻找黄金，也许找不到；无论谁寻求我，我都与他同在。你爱荣誉，也许达不到；有谁爱过我却未达到呢？天主对你说：你想找一个保护者，或一个有权势的朋友；你通过一个较低级的人寻求他的恩宠。爱我吧，天主对你说：与我亲近，不需要通过别人来说项；你的爱本身就使我临在于你。还有什么比这种爱更甜蜜的呢，弟兄们？你们刚才在圣咏中听到，并非没有理由：\Quote{不义者向我谈论喜乐，但不像你的法律，上主。}天主的法律是什么？天主的诫命。天主的诫命是什么？那\Quote{新命令，}它被称为新，因为它使人更新：\BibleQuote{我给你们一条新命令，你们要彼此相爱。}\BibleRef{若 13:34} 听，因为这就是天主的法律。宗徒说：\BibleQuote{你们要彼此分担重担，这样你们就满全了基督的法律。}\BibleRef{迦 6:2} 这就是，这正是我们一切工作的总归：爱。爱是终点：我们为此奔跑；我们奔向它；当我们到达时，我们必得安息。\par

5. 你们在圣咏中听见：\Quote{我看见了一切成全的终向。} 祂说：我看见了一切成全的终向：祂看见了什么？请想想，祂是否登上一座极高而陡峭的山巅，从那里眺望，看见了大地的周围和圆世界的环，因此说：\Quote{我看见了一切成全的终向}？如果这是一件可赞美的事，让我们求主赐予肉身的眼睛如此锐利，我们只需要在地上找一座极高的山，从山巅上便能看见一切成全的终向。不要走远：看，我对你们说，它就在这里；登山，并看见终向。基督就是那座山；来到基督跟前，你们从那里看见一切成全的终向。这个终向是什么？请问保禄：\BibleQuote{但诫命的终向是爱德，出自纯洁的心、良好的良心和真诚的信德。} \BibleRef{弟前 1:5} 又在另一处：\BibleQuote{爱德是法律的圆满，}或成全，\BibleQuote{法律的。} 有什么像\Quote{圆满}一样既完成又终止呢？因为，弟兄们，宗徒在这里以赞美的方式使用\OriginalTerm{终向}{end}。不要把它想成消灭，而是成全。因为，在一方面有人说：我吃完了我的饼；在另一方面有人说：我完成了我的外衣。我吃完了饼，是吃尽了；完成了外衣，是制成它。两处都用\Quote{终了}、\Quote{完成}这个词：但饼的终了是因其被耗尽，外衣的终了是因它被制成；饼如此便不再存在，外衣如此便得以完成。因此，当你们诵读圣咏并听见\Quote{终向，达味的圣咏}时，也要以此意理解\Quote{终向}这个词。你们常在圣咏中听见这话，应该明白所听见的是什么。\Quote{关于终向}是什么意思？——\BibleQuote{因为基督是法律的终向，使每个相信者获得正义。} \BibleRef{罗 13:10} 而\Quote{基督是终向}是什么意思？因为基督是天主，而\Quote{诫命的终向是爱德}，且\Quote{天主是爱}：因为父、子、圣神是一体。祂在那里是你们的终向；在别处祂是道路。不要困在道路上，以至永远不能到达终向。无论你们到达什么别的东西，都要超越它，直到你们到达终向。终向是什么？\Quote{我紧紧依附于天主，为我正是好事。} 你们若紧紧抓住了天主，便走完了道路；你们将在自己的家乡安息。注意！有人寻求金钱：不要让它成为你们的终向；继续前行，如同异乡中的旅客。但若你们爱它，你们就被贪婪所缠；贪婪将成为你们脚上的镣铐；你们便不能再前进。因此，你们也要超越它：寻求终向。你们寻求身体的健康：仍不要停在那里。这身体的健康是什么呢？是死亡使之终结、疾病使之衰弱、脆弱、有死、短暂的东西吗？寻求它，诚然，免得健康不佳妨碍你们的善工；但正是因为如此，终向不在那里，因为它是为了别的事而被寻求。凡是为了别的事而被寻求的，终向就不在那里；凡是为了自身而被爱、且被自由地爱着的，终向就在那里。你们寻求荣誉；或许是为了完成某事而寻求它，好成就某事，从而悦乐天主：不要爱荣誉本身，免得你们停在那里。你们寻求赞美？若寻求天主的赞美，你们做得对；若寻求自己的赞美，你们做得不对；你们在半途就停滞了。但是，看，你们被爱，被赞美：不要把喜乐放在你们自己被赞美上；要在主内被赞美，好使你们能咏唱：\Quote{在主内我的灵魂要受赞美。} 你们发表了一篇好的言论，你们的言论受到赞美。不要让它作为你们自己的而受赞美，终向不在那里。若你们把终向放在那里，那就是你们的终了：但这不是成全，而是消灭。因此，不要因你们自己、作为你们的，而让你们的言论受赞美。但应如何受赞美呢？如圣咏所说：\Quote{在天主内我要赞美言论，在天主内我要赞美话语。} 如此，那里紧接着的经文将在你们身上成就：\Quote{我在天主内寄望，不怕人能如何对待我。} 因为当所有属于你们的东西都在天主内被赞美时，就无需害怕你们的赞美会失落，因为天主不会失落。因此，你们也要超越这一点。\par

6. 弟兄们，请看，我们经过了多少事物，其中并没有终点。我们使用它们如同行路；我们在旅途的歇脚处稍作休憩，然后继续前行。那么终点在哪里？\Quote{亲爱的，我们现在是天主的子女，但我们将来如何，还没有显明；}\BibleRef{若一 3:2}这话就在这封书信里。因此，我们尚在旅途；无论我们来到何处，都必须继续前行，直到我们抵达某个终点。\Quote{我们知道，当祂显现时，我们将相似祂，因为我们将看见祂实在怎样。}那就是终点；那里有永恒的赞美，那里有永远不辍的阿肋路亚。这就是他在圣咏中所说的终点：\Quote{我看见了万有的终向：}好像有人问他：你看见的终向是什么？\Quote{祢的诫命极其宽广。}这就是终点：诫命的宽广。诫命的宽广就是爱德，因为哪里有爱德，那里就没有狭窄。在这宽广中，在这宽阔之处，宗徒曾说：\Quote{我们向你们，格林多人，口是张开的，心是宽大的：你们在我们内并不狭窄。}\BibleRef{格后 6:11}因此，\Quote{祢的诫命极其宽广。}什么是宽广的诫命？\Quote{我给你们一条新命令：你们该彼此相爱。}所以，爱德并不狭窄。你愿意在地上不被狭窄所困吗？请住在宽广之处。因为无论人怎样对你，他不能使你狭窄；因为你所爱的是人不能伤害的：你爱天主，爱兄弟之爱，爱天主的法律，爱天主的教会：它将永远常存。你在地上劳苦，但你会来到所应许的享受。谁能夺走你所爱的呢？若无人能夺走你所爱的，你便安然入睡，或者更确切地说，你安然醒着，免得因沉睡而失去你所爱的。因为圣经并非无缘无故地说：\Quote{照亮我的眼睛，免得我沉睡至死。}那些闭眼不看爱德的人，便沉溺于肉欲的享乐而睡着了。因此，要警醒。因为享乐是：吃喝、奢侈、嬉戏、打猎；这些虚幻的浮华之后，一切邪恶随之而来。难道我们不知道它们是享乐吗？谁能否认它们令人愉悦？但更可爱的是天主的法律。要呼喊反对那些诱惑者：\Quote{不法之徒向我讲论享乐，但祢的法律，上主，并非如此。}这种喜乐存留。它不仅是你可到达的终点，而且当你逃离时，它也呼唤你回来。\par

7. \Quote{因为爱天主，就是遵行祂的诫命。}\BibleRef{若一 4:3}你们已经听过：\Quote{全部法律和先知都系于这两条诫命。}请看，祂不愿你们将自己分散在众多书卷中：\Quote{全部法律和先知都系于这两条诫命。}哪两条诫命？\Quote{你当全心、全灵、全意爱上主你的天主。第二条是：你当爱近人如你自己。全部法律和先知都系于这两条诫命。}\BibleRef{玛 22:37}请看，这整封书信所谈论的就是这些诫命。因此，要持守爱德，安心度日。你为什么害怕会向某人作恶呢？谁爱一个人却向他作恶呢？爱：不可能不做好事。但或许你会责斥。那是出于仁慈，而非严厉。但或许你会责打。你这样做是为了管教，因为你的仁慈之爱不容许他无人管教。确实，有时会产生这些相反的结果：有时仇恨采用奉承的方式，而爱德却表现得严厉。某人恨他的仇人，却对他佯装友好：见他做坏事，便称赞他；希望他一意孤行，希望他盲目地堕入情欲的深渊，或许永不回头；他称赞他，\Quote{因为罪人随从自己的欲望而受赞美；}他以谄媚的油膏抹他；看，他恨他，却赞美他。另一个人看见朋友做了类似的事，便将他唤回；若他不听，甚至用责备的话，他责骂，他争吵；有时甚至会到争吵的地步！看，仇恨表现得很温和，而爱德却争吵！不要注视那看似温和的话语或斥责的残酷；要洞察它们的来源；寻找它们的根。一个是温柔和顺以欺骗，另一个是争吵以改正。所以，弟兄们，我们并非要使你们的心放宽；要从天主那里求得彼此相爱的恩赐。要爱所有的人，甚至你们的仇人，不是因为他们已经是你们的弟兄，而是为了使他们成为你们的弟兄；好使你们时刻燃起兄弟之爱，无论是对已经成了你们弟兄的人，还是对你们的仇人，使他因被爱而成为你们的弟兄。无论你们在何处爱一个弟兄，你们就是爱一个朋友。现在他与你们同在，现在他与你们在合一中，是的，在至公的合一中紧密结合。如果你们生活正直，你们会爱一个从仇人转变为弟兄的人。但你们爱一个人，他尚未信仰基督，或者即使信仰，也像魔鬼那样信：你们责备他的虚妄。要爱，且以兄弟之爱去爱：他尚不是弟兄，但你们爱他是为了使他成为弟兄。因此，我们所有的爱都是兄弟之爱，对基督徒，对祂的所有肢体。爱德的纪律，弟兄们，它的力量、花朵、果实、美丽、欢愉、食粮、饮料、食物、拥抱，在其中没有厌腻。如果它在我们客居他乡时如此令我们喜悦，那么在我们自己的家乡，我们将如何欢乐呢！\par

8. 那么，我的弟兄们，让我们奔跑，让我们奔跑，并爱基督。哪一位基督？耶稣基督。祂是谁？天主的圣言。祂如何来到患病者中间？\Quote{圣言成了血肉，居住在我们中间。}\BibleRef{若 1:14} 这样，圣经预言的\Quote{基督必须受苦，第三天从死者中复活}\BibleRef{路 24:46} 便完全了。祂的身体在哪里？祂的肢体在哪里劳苦？你必须在哪里，才能在你的头之下？\Quote{并且必须从耶路撒冷开始，因祂的名向万邦宣讲悔改和赦罪。}\BibleRef{路 24:47} 让你的爱德在那里广传。基督说，并且\BibleBook{}{圣咏}——\Emph{即}天主圣神——也说：\Quote{你的诫命极其宽广。} 诚然，有人竟想把爱德局限在非洲！你若爱基督，就要将你的爱德扩展到普世，因为基督的肢体遍布普世。你若只爱一部分，你就是分裂的；你若分裂，就不在身体内；你若不在身体内，就不在头之下。你相信却亵渎，这于你何益？你在头上朝拜祂，却在身体上亵渎祂。祂爱祂的身体。你若从祂的身体上切断自己，头并没有从祂的身体上切断自己。你的头从高处向你呼喊：你尊敬我是枉然的，你尊敬我是枉然的。这好比一个人吻你的头，却踩你的脚：他或许穿着钉靴压碎你的脚，同时却抱住你的头亲吻；你难道不会在他恭维你的话语中大喊说：\Quote{你在做什么，人哪？你踩着我！} 你不能说\Quote{你踩着我的头}，因为头是他尊敬的部分；但头为被踩的肢体呼喊，比为自己被尊敬呼喊得更厉害。头岂不自己呼喊：\Quote{我不要你的尊敬，不要踩我！} 现在，你若能，就说：\Quote{我怎么踩了你？} 你对头说：\Quote{我想吻你，我想拥抱你。} 可是，愚人哪，你难道看不出，你所要拥抱的，因着那将全身连接在一起的某种合一，也延伸到你所踩的部分吗？上面你尊敬我，下面你踩我。你踩的那部分，比你所尊敬的那部分更感到疼痛。舌头是怎样呼喊的？\Quote{它伤了我。} 它不说：\Quote{它伤了我的脚，} 而是说：\Quote{它伤了我。} 哦，舌头，谁碰了你？谁打了你？谁刺了你？谁扎了你？没有人，但我与那被踩的部分是连在一起的。既然我不分离，你怎能叫我不痛？\par

9. \Person{我们的主耶稣基督}在第四十日升天时，为此原因将自己的身体托付给我们，使身体继续躺卧在地，因为他看到许多人会因他升天而尊敬他，并且看到如果他们在世上践踏他的肢体，这种尊敬便是无用的。为避免任何人犯错，当他敬拜天上的头时，却践踏地上的脚，他告诉我们他的肢体将在何处。因为他即将升天，他说了他在世上的最后话语；此后他不再在世上说话。将要升天的头将地上的肢体托付给我们，然后离开了。从此你找不到基督在世上说话；你发现他说话，却是从天上。甚至从天上，为何？因为他在地上的肢体被践踏了。对迫害者扫禄，他从高处说：\BibleQuote{扫禄，扫禄，你为什么迫害我？}\BibleRef{宗 9:4} 我已升天，但仍躺在地上；我在此坐在圣父的右边，但在那里我仍然饥饿、口渴、作客旅。那么，当他即将升天时，他是如何将他的身体托付给我们的呢？当门徒问他，说：\BibleQuote{主，你要在这时候复兴以色列国吗？}\BibleRef{宗 1:6} 他即将离去时回答说：\Quote{父以自己的权柄所定的时候、日期，不是你们可以知道的；但圣神降临在你们身上，你们就必得着能力，并要作我的见证人。} 看哪，他的身体如何扩展，看哪，他将不被践踏：\Quote{你们要作我的见证人，在耶路撒冷、犹太全地，直到地极。} 看哪，我升天时躺卧之处！因为我升天，因为我是头；我的身体仍躺卧在下。何处躺卧？在整个大地。当心你们不要击打，当心不要伤害，当心不要践踏：这是基督即将进入天堂的最后话语。看看一个患病的人，卧病在床，躺在家里，被疾病折磨得奄奄一息，灵魂仿佛已在齿间：他可能对某件他所珍爱、深爱的东西焦虑，想起这事，便召唤他的继承人，对他们说：求你们做这件事。他仿佛用强力留住自己的灵魂，以免在这些话确定之前离去。当他口述完这些最后的话后，他呼出灵魂，被抬着尸体送往坟墓。他的继承人如何记得垂死之人的最后话语？如果有人站起来对他们说：不要这样做，他们会怎么说？\Quote{什么？难道我不该做我父亲临终时用最后一口气命令我的事吗？那是他在我耳边最后的话语，当我父亲离开此世时？无论他其他的话我可能不理会，他最后的话对我有更强的约束：此后我再也见不到他，再也听不到他的声音。} 弟兄们，请以基督徒的心思考：如果对于一个人的继承人而言，他即将进入坟墓时所说的话是如此甜蜜、如此感激、如此沉重，那么对基督的最后话语，我们该如何看待？他不是要回到坟墓，而是要升天！至于那活着却死了的人，他的灵魂被匆忙带到别处，身体被埋入大地，无论他的话是否被遵行，对他都无关紧要：他现在有别的事要做，或要承受别的事：或在亚巴郎的怀中喜乐，或在地狱之火中渴望一滴水，而他的尸体无知觉地躺在坟墓中；然而垂死之人的最后话语仍被遵守。那些不遵守坐在天上的那一位的最后话语的人，又有什么期待呢？他从高处看着这些话是被轻视还是不被轻视？正是那一位的话语，他曾说：\BibleQuote{扫禄，扫禄，你为什么迫害我？} 他记下他看见他的肢体所受的一切痛苦，直到审判之时。\par

10. 他们又说：我们做了什么？我们是受迫害者，不是迫害者。你们是迫害者，可怜的人。首先，因为你们分裂了教会。舌头的剑比钢铁的剑更锋利。撒辣的婢女阿加尔骄傲了，她因骄傲被主母责罚。那是管教，不是惩罚。因此，当她离开主母时，天使对她说：\BibleQuote{回到你主母那里去。}\BibleRef{创 16:4} 那么，属肉体的灵魂，如同骄傲的婢女，假设你为了管教而遭受了任何麻烦，你为什么发狂？\Quote{回到你主母那里去，} 持守教会的和平。看哪，福音被提出，我们读到教会扩展到哪里。人们与我们争辩，对我们说：\Quote{出卖者！} 出卖什么？基督将他的教会托付给我们，你们不信：难道我要相信你们，当你们毁谤我的父母时？难道你们要我因\Quote{出卖者}而相信你们？先相信基督。什么值得相信？基督是天主，你们是人：哪个更应首先相信？基督已使他的教会扩展到全地：我说了——你们轻视我；福音说了——当心。福音说什么？\BibleQuote{基督必须受苦，第三日从死者中复活，并且要奉他的名宣讲悔改及罪过的赦免。}\BibleRef{路 24:47} 那里有罪过的赦免，那里就有教会。因为对她说了：\BibleQuote{我要将天国的钥匙交给你；凡你在地上所束缚的，在天上也要被束缚；凡你在地上所释放的，在天上也要被释放。}\BibleRef{玛 16:19} 这罪过的赦免扩展到何处？\Quote{从耶路撒冷起，直到万邦。} 看哪，相信基督！但是，因为你们清楚知道，如果你们相信基督，你们将不再有关于\Quote{出卖者}的任何话可说，你们却宁愿让我在你们毁谤我的父母时相信你们，而不愿你们自己相信基督所预言的！\par

\EditorialNote{本讲道的剩余部分在所有手稿中均缺失。此外，圣奥思定似乎未能如他所承诺的那样完成整篇书信的注释：至少波西狄乌斯将这部作品列为《论若望宗徒书致帕提亚人十篇讲道》，而整个第五章在这第十篇讲道中得到阐释的可能性很小。关于讲道集，没有关于这封书信剩余部分的讲道：圣奥思定其他著作中的以下摘录将提供最需要的部分：即他对\Quote{三个见证}、至于死的罪和第二十节经文的注释。}\par

\SourceRecord{Translated by H. Browne.；Nicene and Post-Nicene Fathers, First Series；Vol. 7.；Edited by Philip Schaff.；Buffalo, NY: Christian Literature Publishing Co.,；1888.；<http://www.newadvent.org/fathers/170210.htm>.}
